\documentclass[11pt,a4paper,oneside]{ctexbook}

\usepackage[a4paper,margin=1.8cm]{geometry}
\usepackage{amsmath,amssymb}
\usepackage{booktabs}
\usepackage{array}
\usepackage{longtable}
\usepackage{tabularx}
\usepackage{xcolor}
\usepackage{hyperref}
\usepackage{enumitem}
\usepackage{listings}
\usepackage{fancyhdr}
\usepackage{titlesec}
\usepackage{tikz}

\setlength{\headheight}{14pt}

\hypersetup{
  colorlinks=true,
  linkcolor=blue!60!black,
  urlcolor=blue!60!black,
  citecolor=blue!60!black
}

\pagestyle{fancy}
\fancyhf{}
\lhead{CCF CSP 考场手册}
\rhead{\leftmark}
\cfoot{\thepage}

\setlist[itemize]{leftmargin=2em,itemsep=0.2em,topsep=0.2em}
\setlist[enumerate]{leftmargin=2em,itemsep=0.2em,topsep=0.2em}

\definecolor{DangerRed}{RGB}{170,30,30}
\definecolor{NoteBlue}{RGB}{30,80,150}
\definecolor{GoodGreen}{RGB}{40,120,70}
\definecolor{SoftGray}{RGB}{245,245,245}
\definecolor{CoverNavy}{RGB}{18,42,72}
\definecolor{CoverBlue}{RGB}{37,99,160}
\definecolor{CoverLine}{RGB}{214,224,235}
\definecolor{CoverSoft}{RGB}{246,249,252}
\definecolor{CoverText}{RGB}{38,50,65}

\lstdefinestyle{cppstyle}{
  language=C++,
  basicstyle=\ttfamily\small,
  keywordstyle=\color{blue!70!black},
  commentstyle=\color{GoodGreen},
  stringstyle=\color{DangerRed},
  numbers=left,
  numberstyle=\tiny\color{gray},
  stepnumber=1,
  numbersep=8pt,
  showstringspaces=false,
  breaklines=true,
  frame=single,
  columns=fullflexible,
  tabsize=4
}
\lstset{style=cppstyle}

\newcommand{\Key}[1]{\textbf{\color{NoteBlue}#1}}
\newcommand{\Warn}[1]{\textbf{\color{DangerRed}#1}}
\newcommand{\Good}[1]{\textbf{\color{GoodGreen}#1}}
\newcommand{\Code}[1]{\texttt{#1}}

\begin{document}

\frontmatter
\begin{titlepage}
  \thispagestyle{empty}
  \newgeometry{left=2.2cm,right=2.2cm,top=2.0cm,bottom=2.0cm}

  \begin{tikzpicture}[remember picture,overlay]
    \fill[CoverSoft] (current page.south west) rectangle (current page.north east);
    \fill[CoverNavy] (current page.north west) rectangle ([yshift=-8.15cm]current page.north east);
    \fill[CoverBlue] ([yshift=-8.15cm]current page.north west) rectangle ([yshift=-8.33cm]current page.north east);
    \draw[CoverLine,line width=0.9pt]
      ([xshift=1.45cm,yshift=-1.45cm]current page.north west)
      rectangle
      ([xshift=-1.45cm,yshift=1.45cm]current page.south east);
    \draw[CoverBlue!55,line width=1.2pt]
      ([xshift=2.0cm,yshift=-9.7cm]current page.north west) --
      ([xshift=-2.0cm,yshift=-9.7cm]current page.north east);
    \fill[white,opacity=0.12]
      ([xshift=-5.8cm,yshift=-1.45cm]current page.north east) --
      ([xshift=-1.45cm,yshift=-1.45cm]current page.north east) --
      ([xshift=-1.45cm,yshift=-5.8cm]current page.north east) -- cycle;
    \draw[CoverBlue!45,line width=0.8pt]
      ([xshift=2.0cm,yshift=-10.45cm]current page.north west) --
      ([xshift=8.1cm,yshift=-10.45cm]current page.north west);
  \end{tikzpicture}

  \begin{center}
    \vspace*{0.45cm}
    {\Large\color{CoverLine}\bfseries CCF CSP Handbook}\\[1.05cm]
    {\fontsize{36}{44}\selectfont\bfseries\color{white} CCF CSP 考场手册}\\[0.55cm]
    {\Large\color{CoverLine} 题型决策树、C++ 模板与真题索引体系}
  \end{center}

  \vspace{3.15cm}

  {\Large\bfseries\color{CoverNavy} 面向考场检索的开源备考资料}\\[0.65cm]
  {\color{CoverText}
  本手册围绕“题目识别、数据范围判断、算法模板选择、易错点检查、真题反查”组织内容，
  目标是在复习和临场翻阅时都能快速定位可执行方案。
  }

  \vspace{1.05cm}

  \begin{center}
    \setlength{\tabcolsep}{8pt}
    \renewcommand{\arraystretch}{1.35}
    \begin{tabular}{>{\bfseries\color{CoverBlue}}r p{9.3cm}}
      判断 & 从题面信号和数据范围定位题型。\\
      模板 & 汇总 C++、STL、算法和大模拟骨架。\\
      检查 & 覆盖下标、溢出、取整、格式和部分分策略。\\
      校准 & 用历年真题反查判断树和模板入口。\\
    \end{tabular}
  \end{center}

  \vfill

  \begin{center}
    {\color{CoverLine}\rule{0.72\textwidth}{0.9pt}}\\[0.55cm]
    {\large\color{CoverNavy} 开源备考资料 \quad · \quad 可打印 \quad · \quad 可持续修订}\\[0.35cm]
    {\color{CoverBlue} GitHub 发布版}
  \end{center}
  \restoregeometry
\end{titlepage}

\chapter*{使用前说明}
\addcontentsline{toc}{chapter}{使用前说明}

这本手册的目标不是覆盖所有算法竞赛知识，而是在 CCF CSP 考场中帮助你快速完成四件事：

\begin{enumerate}
  \item 根据题面关键词和数据范围判断题型。
  \item 找到可用算法、C++ 语法、STL 容器和代码模板。
  \item 避免输出格式、下标、溢出、初始化等低级错误。
  \item 遇到难题时先写暴力或特殊情况，尽量拿部分分。
\end{enumerate}

\Warn{重要提醒：}正式打印和考试前，必须重新核对 CCF CSP 官网的最新规则、语言标准、考场环境和允许携带材料说明。

\tableofcontents

\mainmatter

\part{导航与判断}
\chapter{全书总览与使用路径}

\section{本手册的工作方式}

这本手册按“题目识别”而不是按“算法教材目录”组织。考场上你通常不是先想起某个算法，再去找题目匹配它；更常见的是先看到题面中的若干信号，例如区间、多次询问、最短、连通、排序、字符串、日期、进制、命令序列，然后再判断应该翻到哪个章节。

因此，本手册的每个决策入口都尽量回答下面七个问题：

\begin{enumerate}
  \item 题面中出现了什么信号？
  \item 数据范围允许什么复杂度？
  \item 它更像 A、B、C、D、E 中哪类题？
  \item 先尝试哪个模板？
  \item 需要使用哪些 C++ 语法或 STL 容器？
  \item 最容易错在哪里？
  \item 如果满分做法不会，能否先写部分分？
\end{enumerate}

\Warn{使用原则：}不要在考场上临时学习一个完全没练过的模板。模板页是用来提醒和防错的，不是用来从零学习的。

\section{全书结构路线图}

本书按考场使用顺序重排为四个分部。先用索引定位，再用题型章节展开，最后用模板和真题反向校准。

\begin{longtable}{p{3.2cm}p{4.2cm}p{6.6cm}}
\toprule
分部 & 章节 & 用法 \\
\midrule
导航与判断 & 第 1--3 章 & 先看全书用法，再用 1 分钟极速索引定位，必要时进入详细判断树。 \\
基础与题型 & 第 4--8 章 & 先补 C++ 和 STL 基础，再按 A/B/C/D/E 题型查看策略和高频入口。 \\
模板与检查 & 第 9--10 章 & 找可复制模板，并在提交前扫易错点。 \\
真题校准 & 第 11 章 & 用历年题反查模板和判断树，发现误判后回到第 3 章修正。 \\
\bottomrule
\end{longtable}

推荐阅读路线：

\begin{enumerate}
  \item \Key{考前通读：}第 1 章、第 4 章、第 5--8 章、第 10 章。
  \item \Key{刷题复盘：}第 3 章判断路径、第 9 章模板、第 11 章真题索引。
  \item \Key{考场翻查：}第 2 章极速索引；定位不准时翻第 3 章；写代码时翻第 9 章；提交前翻第 10 章。
\end{enumerate}

\section{官方规则速查}

本节只作为考前检查清单。正式打印前必须重新访问 CCF CSP 官网核对最新版本。

\begin{longtable}{p{3.2cm}p{10.8cm}}
\toprule
项目 & 当前应记录的信息 \\
\midrule
官方网站 & \url{https://www.cspro.org/} \\
认证形式 & 上机编程，黑盒测试，按测试点得分。 \\
题目数量 & 通常为 5 道编程题。每题 100 分，总分 500 分。 \\
考试时间 & 通常为 240 分钟。 \\
输入输出 & 使用标准输入、标准输出。不要输出提示语、调试信息或多余空格。 \\
提交方式 & 每题提交一个源文件；同一题多次提交时，一般以最后一次提交为准。 \\
语言环境 & 以官网报名系统和当次考试公告为准。C++ 版本、编译器和操作系统可能变化。 \\
纸质资料 & 根据 CCF 官网答疑和考生须知，应考时可携带纸质书籍、打印资料、手写资料；不得使用通讯工具、存储设备。考前必须再次核对。 \\
程序长度 & 官网注意事项中提到程序长度限制，准备模板时不要复制大量无关代码。 \\
\bottomrule
\end{longtable}

建议在打印版首页手写确认：

\begin{itemize}
  \item 我已经在考前一周核对了官网规则。
  \item 我已经确认了本次考试允许携带的纸质资料范围。
  \item 我已经确认了本次考试 C++ 编译标准。
  \item 我知道不能携带电子资料、存储设备或联网设备。
\end{itemize}

\section{考场总流程}

\subsection{开场 10 分钟}

开场不要立刻写代码。先快速浏览 5 道题，目标是建立时间分配和得分路线。

\begin{enumerate}
  \item 看每道题的题面长度。题面短不一定简单，题面长往往意味着模拟或细节多。
  \item 看每道题的数据范围。数据范围比题面故事更重要。
  \item 判断 A、B 是否可以快速完成。
  \item 判断 C 是算法题还是大模拟。
  \item 判断 D、E 是否有子任务、小数据、特殊性质或熟悉模板。
  \item 在草稿纸上写下计划：先做哪题、卡住多久换题、最后回来看哪题。
\end{enumerate}

\subsection{每题固定流程}

每道题都按同一套流程处理：

\begin{enumerate}
  \item 读题面，圈出关键词。
  \item 抄下所有数据范围。
  \item 估算可接受复杂度。
  \item 查本章总决策树，定位章节。
  \item 写出暴力思路，用来理解题意和造样例。
  \item 如果满分模板明确，再写正式代码。
  \item 先过样例，再造边界样例。
  \item 删除调试输出后提交。
\end{enumerate}

\Warn{不要跳过复杂度估算。}很多新手失分不是因为不会算法，而是把 $O(n^2)$ 写到了 $n=10^5$ 的题里。

\subsection{建议时间分配}

\begin{longtable}{p{2cm}p{3cm}p{5cm}p{4cm}}
\toprule
题号 & 建议时间 & 目标 & 卡住时的处理 \\
\midrule
A & 20--35 分钟 & 稳定满分 & 15 分钟没思路就重新读题；30 分钟没过样例先跳。 \\
B & 40--70 分钟 & 尽量满分 & 优先查前缀和、差分、排序、哈希、BFS。 \\
C & 60--100 分钟 & 高分或满分 & 如果是大模拟，先拆结构体和函数，不要全写在主函数。 \\
D & 40--80 分钟 & 部分分到满分 & 先找小数据和熟悉模板。不会满分也要写暴力。 \\
E & 剩余时间 & 尽量拿分 & 重点看子任务、特殊性质、样例规律。 \\
\bottomrule
\end{longtable}

时间分配不是固定规则。若 A、B 很顺，可以把更多时间留给 C、D；若 A、B 卡住，不要让前两题吞掉整场考试。

\section{数据范围与复杂度}

\subsection{复杂度速查表}

\begin{longtable}{p{4cm}p{5cm}p{5cm}}
\toprule
输入规模 & 通常可考虑 & 常见模板 \\
\midrule
$n \le 20$ & 指数搜索、状态压缩、全排列、DFS 暴力 & DFS、状压 DP、枚举子集 \\
$n \le 100$ & $O(n^3)$ 通常可试 & Floyd、区间 DP、三重枚举 \\
$n \le 500$ & $O(n^3)$ 可能可行，需看时限 & Floyd、部分 DP \\
$n \le 1000$ & $O(n^2)$ 通常可行 & 双重循环、二维 DP、朴素图算法 \\
$n \le 5000$ & $O(n^2)$ 需谨慎 & 优化前的 DP、枚举配合剪枝 \\
$n \le 10^5$ & $O(n \log n)$ 或 $O(n)$ & 排序、二分、前缀和、差分、堆、树状数组 \\
$n \le 10^6$ & $O(n)$ 或较轻的 $O(n \log n)$ & 线性扫描、计数、前缀和 \\
$n \ge 10^7$ & 通常只能线性或公式 & 数学、简单扫描、避免大对象数组 \\
\bottomrule
\end{longtable}

\subsection{复杂度校验规则}

\begin{itemize}
  \item 看到双重循环，先计算最坏次数。
  \item 有多组数据时，总复杂度要乘以组数，或者看所有 $n$ 的总和限制。
  \item 排序一般是 $O(n \log n)$，$n=10^5$ 常见可行。
  \item Floyd 是 $O(n^3)$，只适合点数较小的图。
  \item Dijkstra 配合堆是 $O((n+m)\log n)$，适合正权稀疏图。
  \item BFS 适合边权全为 1 或每步代价相同的最短步数。
  \item 前缀和适合不修改的区间查询。
  \item 差分适合批量区间修改后统一查询。
  \item 树状数组和线段树适合修改与查询交织出现。
\end{itemize}

\Warn{溢出提醒：}只要出现求和、乘法、路径长度、方案数、时间累计，默认先考虑 \Code{long long}。两个 \Code{int} 相乘时写成 \Code{1LL * a * b}。

\section{A/B/C/D/E 总体策略}

\subsection{A 题}

A 题通常考察基本实现能力。它可能是简单模拟、统计、字符串、日期、进制、排序或公式题。

\begin{itemize}
  \item 目标：尽量一次 AC。
  \item 优先检查：输入格式、输出格式、整数溢出、小数精度、边界值。
  \item 常用章节：C++ 基础、数组统计、字符串、进制、日期、排序。
  \item 不建议：为了追求代码短而写不熟悉的技巧。
\end{itemize}

\subsection{B 题}

B 题常常是第一个真正需要算法识别的题，但算法通常不会太深。

\begin{itemize}
  \item 高频方向：排序、前缀和、差分、哈希、二分、BFS、简单图、简单 DP。
  \item 读题重点：是否有多次询问，是否有区间，是否需要快速统计。
  \item 常见错误：用暴力过样例但过不了大数据。
\end{itemize}

\subsection{C 题}

C 题经常是分水岭。它可能是大模拟，也可能是中等算法题。

\begin{itemize}
  \item 如果题面很长，先判断是否是规则模拟。
  \item 如果输入是命令、事件、状态变化，优先设计结构体和函数。
  \item 如果出现图、网格、路径、依赖关系，再进入对应算法分支。
  \item 大模拟题不要把所有逻辑塞进 \Code{main}。
\end{itemize}

\subsection{D/E 题}

D、E 题要同时考虑满分算法和部分分。

\begin{itemize}
  \item 先看是否给出子任务或特殊数据范围。
  \item 先写能过小数据的暴力，避免空题。
  \item 若满分算法明显且练过，再冲满分。
  \item 若 20 分钟没有形成可写方案，转为部分分策略。
\end{itemize}

\section{总决策树}

\subsection{第一层：看题面信号}

\begin{longtable}{p{4.2cm}p{4.8cm}p{5cm}}
\toprule
题面信号 & 优先怀疑 & 先翻章节 \\
\midrule
字符、单词、路径、命令、表达式、括号 & 字符串处理或解析模拟 & 字符串专题、栈、大模拟 \\
二进制、十六进制、位、补码、编码 & 进制和位运算 & 进制专题、位运算 \\
日期、时间、星期、天数、日历 & 日期时间处理 & 日期专题 \\
区间和、多次查询、子矩阵和 & 前缀和 & 前缀和模板 \\
区间加、范围修改、最后输出 & 差分 & 差分模板 \\
修改和查询交替出现 & 树状数组或线段树 & 数据结构模板 \\
最大值、最小值动态维护 & 堆、集合、单调结构 & 优先队列、set、单调队列 \\
排序、排名、优先级、相同条件按某字段 & 自定义排序 & sort 与 lambda \\
去重、出现次数、频率统计 & 哈希或 map & map、unordered map、计数数组 \\
最少步数、迷宫、网格移动 & BFS & BFS、网格 BFS \\
正权最短路、道路费用、距离最小 & Dijkstra & 最短路模板 \\
任意两点最短路且点数小 & Floyd & Floyd 模板 \\
合并集合、是否连通、同一类 & 并查集 & 并查集模板 \\
依赖、先后顺序、任务调度 & 拓扑排序 & 拓扑排序模板 \\
最大化最小值、最小化最大值、答案可行性 & 二分答案 & 二分答案模板 \\
容量、价值、选择若干个、最多或恰好 & 背包 DP & 动态规划模板 \\
规则很多、对象很多、按时间变化 & 大模拟 & 大模拟专题 \\
\bottomrule
\end{longtable}

\subsection{第二层：用数据范围排除错误方向}

\begin{longtable}{p{5cm}p{4.5cm}p{4.5cm}}
\toprule
如果你想写 & 但数据范围是 & 应改查 \\
\midrule
枚举所有点对 $O(n^2)$ & $n=10^5$ & 排序、哈希、前缀和、树状数组、双指针 \\
枚举所有区间 $O(n^2)$ & $n=10^5$ & 前缀和、双指针、单调队列、DP 优化 \\
Floyd $O(n^3)$ & $n>500$ 或边很多 & Dijkstra、BFS、图的特殊性质 \\
每次查询都扫描数组 & 查询数 $q=10^5$ & 前缀和、树状数组、线段树、离线处理 \\
递归搜索所有方案 & $n>25$ 且无剪枝 & DP、贪心、数学、二分答案 \\
字符串反复插入删除 & 字符串长度和操作数很大 & 分块、链表思想、栈、模拟优化 \\
\bottomrule
\end{longtable}

\subsection{第三层：选择满分或部分分}

\begin{enumerate}
  \item 如果模板明确、数据范围匹配、自己练过，写满分做法。
  \item 如果模板大致明确但细节多，先写能跑的小版本，再逐步补全。
  \item 如果满分算法不明确，找子任务和小数据，写暴力或特殊情况。
  \item 如果题目很难但 A/B/C 还没稳，不要长时间停留在 D/E。
\end{enumerate}

\section{关键词到章节索引}

这张表用于快速翻书。后续章节完成后，需要把“章节编号”和“模板编号”补齐。

\begin{longtable}{p{4cm}p{4.5cm}p{5.2cm}}
\toprule
关键词 & 翻到哪里 & 核心提醒 \\
\midrule
字符串、字符、大小写 & 字符串专题 & 注意 \Code{cin} 和 \Code{getline} 混用。 \\
空格分割、命令行 & 字符串分割 & 用 \Code{stringstream} 或手写分割。 \\
括号、嵌套、表达式 & 栈、表达式解析 & 注意优先级和匹配失败。 \\
进制转换 & 进制专题 & 判断是否超过 \Code{long long}。 \\
二进制位 & 位运算 & 大位数用 \Code{1LL << k}。 \\
日期、星期 & 日期专题 & 先确认闰年和月份天数。 \\
区间查询 & 前缀和 & 不修改时优先前缀和。 \\
区间修改 & 差分 & 修改后统一查询时优先差分。 \\
在线修改查询 & 树状数组、线段树 & 看是单点还是区间。 \\
排序、排名 & sort 与 lambda & 相同条件要写完整。 \\
频率、出现次数 & map、unordered map、计数数组 & 值域小用数组更快。 \\
最短步数 & BFS & 边权相同才用 BFS。 \\
正权最短路 & Dijkstra & 距离用 \Code{long long}。 \\
连通、合并 & 并查集 & 路径压缩和按大小合并。 \\
依赖、先后 & 拓扑排序 & 入度为 0 入队。 \\
最优选择 & DP 或贪心 & 先看是否有状态转移。 \\
答案范围 & 二分答案 & 必须能写 \Code{check}。 \\
大量规则 & 大模拟 & 结构体、函数、状态表。 \\
\bottomrule
\end{longtable}

\section{低级错误总检查}

提交前至少检查下面内容：

\begin{itemize}
  \item 是否删除所有调试输出。
  \item 是否使用标准输入输出。
  \item 输出末尾是否多了不该有的内容。
  \item 是否把答案变量开成 \Code{long long}。
  \item 数组是否开到最大范围。
  \item 下标是从 0 还是从 1 开始。
  \item 循环边界是否包含末尾。
  \item 多组数据时是否清空数组、图、队列、map。
  \item 比较器是否满足严格弱序。
  \item \Code{lower\_bound} 或 \Code{find} 的返回值是否检查。
  \item \Code{priority\_queue} 是否需要小根堆。
  \item BFS 是否在入队时标记访问。
  \item Dijkstra 是否跳过过期状态。
  \item 递归深度是否可能爆栈。
\end{itemize}

\section{部分分总策略}

当满分做法暂时不会时，按这个顺序找分：

\begin{enumerate}
  \item \Key{小数据暴力：}如果 $n \le 1000$ 的测试点存在，写 $O(n^2)$；如果 $n \le 100$，考虑 $O(n^3)$。
  \item \Key{特殊结构：}图是树、链、完全图、DAG、无权图时，可能有更简单做法。
  \item \Key{特殊参数：}如果某个参数为 0、1、很小，单独写分支。
  \item \Key{无修改版本：}有修改和查询时，若部分数据没有修改，先用前缀和或静态算法。
  \item \Key{少量查询或少量修改：}查询少就每次暴力；修改少就重算或离线处理。
  \item \Key{样例和明显边界：}最后时间不足时，也要保证样例和简单边界正确。
\end{enumerate}

\Warn{部分分代码也要干净。}不要写一堆互相冲突的特判。特判越多，越容易把能拿的分弄丢。

\section{本章完成后的后续连接}

后续章节需要逐步补齐以下链接：

\begin{itemize}
  \item 将“关键词到章节索引”中的章节编号补全。
  \item 将每个决策树叶子节点连接到具体模板编号。
  \item 将真题索引中的每道题反向连接到本章的某个判断路径。
  \item 将个人错题中的错误补入“低级错误总检查”。
\end{itemize}

\chapter{考场 1 分钟极速索引}\label{ch:02}

\section{这一章怎样查}\label{sec:02-1}

赶时间时可以先翻本章，按题目中的操作查模板编号。需要进一步判断时看第\ref{ch:03}章，具体代码在第\ref{ch:09}章，练习题索引在第\ref{ch:11}章。

表中的关键词只是线索，还要看数组会不会修改、边权有什么限制、数据有多大。

\section{1 分钟总流程}\label{sec:02-2}

\begin{longtable}{p{2.4cm}p{6.0cm}p{5.6cm}}
\toprule
时间 & 做什么 & 目的 \\
\midrule
\endfirsthead
\toprule
时间 & 做什么 & 目的 \\
\midrule
\endhead
前 10 秒 & 看题号：A/B/C/D/E & 决定稳分、冲分、拆题、部分分。 \\
第 10-25 秒 & 圈出输入主体：数组、矩阵、字符串、图、树、对象、事件 & 确定大类。 \\
第 25-40 秒 & 圈出操作：查询、修改、排序、最优、路径、状态变化 & 确定模板方向。 \\
第 40-55 秒 & 看数据范围：$n,m,q$、值域、字符串总长 & 排除不可能复杂度。 \\
最后 5 秒 & 翻模板并写骨架 & 不在脑子里硬背。 \\
\bottomrule
\end{longtable}

\section{题号策略速查}\label{sec:02-3}

\begin{longtable}{p{2cm}p{4.2cm}p{4.2cm}p{3.8cm}}
\toprule
题号 & 目标 & 常见题型 & 第一动作 \\
\midrule
\endfirsthead
\toprule
题号 & 目标 & 常见题型 & 第一动作 \\
\midrule
\endhead
A & 争取满分 & 统计、公式取整、排序、字符串、日期、二维数组 & 先写暴力扫描，检查输出格式。 \\
B & 尽量满分 & 前缀、差分、哈希、排序扫描、二分、简单 DP & 先写暴力，再找重复计算。 \\
C & 尽量拿高分 & 大模拟、字符串解析、状态机、事件、树/目录 & 先建对象和状态，不急着写代码。 \\
D & 能拿部分分就拿 & 图、树、DP、数据结构、数学 & 先看子任务和特殊性质。 \\
E & 尝试明确的子任务 & 高级算法或综合题 & 写小数据暴力，不耗尽时间。 \\
\bottomrule
\end{longtable}

\section{题面主体到模板}\label{sec:02-4}

\begin{longtable}{p{3.2cm}p{5.8cm}p{4.8cm}}
\toprule
主体 & 题面信号 & 模板/章节 \\
\midrule
\endfirsthead
\toprule
主体 & 题面信号 & 模板/章节 \\
\midrule
\endhead
数组 & 求和、计数、相邻、分段、排名、分组 & A 题、\Tref{101}、\Tref{105}、\Tref{106}、\Tref{510} \\
矩阵 & 图像、矩形、邻域、旋转、重塑 & \Tref{102}、\Tref{104}、矩阵扫描 \\
字符串 & 分隔符、路径、URL、括号、表达式 & \Tref{501}、\Tref{606}、\Tref{607}、\Tref{608}、\Tref{609} \\
坐标点 & 点集、距离、矩形、覆盖、邻居 & \Tref{106}、\Tref{109}、坐标公式 \\
图 & 最短路、连通、依赖、最小代价 & \Tref{201}、\Tref{203}、\Tref{204}、\Tref{205}、\Tref{207} \\
树 & 子树、祖先、路径、目录、合并 & DFS、\Tref{208}、\Tref{402} \\
多命令 & create、delete、query、request、release & \Tref{601}、\Tref{602}、\Tref{603} \\
时间事件 & 到期、归还、租约、按时刻发生 & \Tref{110}、\Tref{514}、\Tref{603} \\
选择物品 & 容量、费用、价值、至少达到 & \Tref{301}、\Tref{302} \\
最优答案 & 最小时间、最大阈值、最小最大 & \Tref{107} 或 DP \\
\bottomrule
\end{longtable}

\section{数据范围秒判}\label{sec:02-5}

\begin{longtable}{p{3.6cm}p{5.2cm}p{4.8cm}}
\toprule
范围 & 可考虑的做法 & 通常不合适的做法 \\
\midrule
\endfirsthead
\toprule
范围 & 可考虑的做法 & 通常不合适的做法 \\
\midrule
\endhead
$n \le 20$ & 枚举、DFS、状压 & 复杂数据结构优先 \\
$n \le 100$ & $O(n^3)$、Floyd & 指数暴力 \\
$n \le 1000$ & $O(n^2)$ & $O(n^3)$ 大概率危险 \\
$n \le 10^5$ & $O(n\log n)$、$O(n)$ & $O(n^2)$ \\
$q \le 10^5$ & 每次 $O(\log n)$ 或 $O(1)$ & 每次扫描全体 \\
值域小 & 计数数组、桶 & map 不一定必要 \\
值域大但数量少 & map、hash、离散化 & 直接开值域数组 \\
\bottomrule
\end{longtable}

\section{B 题优化入口}\label{sec:02-6}

\begin{longtable}{p{4.5cm}p{5.3cm}p{4.0cm}}
\toprule
暴力写法 & 应想到 & 模板 \\
\midrule
\endfirsthead
\toprule
暴力写法 & 应想到 & 模板 \\
\midrule
\endhead
每次区间求和都循环 & 前缀和 & \Tref{101} \\
每次矩形求和都双重循环 & 二维前缀和 & \Tref{102} \\
每次区间加都逐个改 & 差分 & \Tref{103} / \Tref{104} \\
每次查是否出现都循环 & set / unordered\_set & \Tref{106} \\
排序后每次找位置 & 二分 / lower\_bound & C++ 语法章 \\
两个有序表逐个配对 & 双指针 & \Tref{108} \\
最小时间、最大阈值 & 二分答案 & \Tref{107} \\
选若干个凑金额 & 背包 & \Tref{301} \\
前置任务、依赖关系 & 拓扑排序 & \Tref{205} \\
\bottomrule
\end{longtable}

\section{C 题建模速查}\label{sec:02-7}

C 题先填这张表，再写代码。

\begin{longtable}{p{3cm}p{5.4cm}p{5.2cm}}
\toprule
要素 & 题面中找什么 & 代码怎么落地 \\
\midrule
\endfirsthead
\toprule
要素 & 题面中找什么 & 代码怎么落地 \\
\midrule
\endhead
对象 & 用户、文件、任务、节点、请求、词条 & \Code{struct Obj} \\
编号 & 数字编号、字符串名、路径 & vector / map / unordered\_map \\
状态 & 有效、过期、占用、空闲、允许、拒绝 & enum/int 状态机 \\
命令 & 每行第一个单词或操作类型 & \Tref{602} 命令分发 \\
时间 & 当前时刻、持续时间、到期时间 & \Tref{110} / \Tref{514} \\
优先级 & “优先”“若相同”“同时发生” & sort 比较器或 priority\_queue \\
失败规则 & 不存在、非法、超额、冲突 & 先判断，后修改，必要时回滚 \\
\bottomrule
\end{longtable}

\section{D/E 部分分入口}\label{sec:02-8}

\begin{longtable}{p{4.0cm}p{5.4cm}p{4.6cm}}
\toprule
看到子任务 & 写什么 & 可能分数来源 \\
\midrule
\endfirsthead
\toprule
看到子任务 & 写什么 & 可能分数来源 \\
\midrule
\endhead
$n \le 20$ & 子集枚举、DFS、状压 & 小数据分 \\
$n \le 100$ & Floyd、$O(n^3)$、暴力 DP & 中小数据分 \\
链 & 当数组处理，前缀/差分 & 特殊性质分 \\
树 & 每次 DFS，或 DFS 序 & 树特殊性质分 \\
无修改 & 预处理后回答查询 & 静态分 \\
查询少 & 每次暴力 & 查询少子任务 \\
边权 1 & BFS & 图特殊性质分 \\
值域小 & 计数数组 & 小值域分 \\
\bottomrule
\end{longtable}

\section{模板编号速查}\label{sec:02-9}

\begingroup
\small
\setlength{\tabcolsep}{4pt}
\begin{longtable}{p{1.7cm}p{5cm}p{1.7cm}p{5cm}}
\toprule
编号/页 & 模板 & 编号/页 & 模板 \\
\midrule
\endfirsthead
\toprule
编号/页 & 模板 & 编号/页 & 模板 \\
\midrule
\endhead
\Tref{101} & 一维前缀和 & \Tref{102} & 二维前缀和 \\
\Tref{103} & 一维差分 & \Tref{104} & 二维差分 \\
\Tref{105} & 排序自定义比较器 & \Tref{106} & 哈希统计 \\
\Tref{107} & 二分答案 & \Tref{108} & 双指针 \\
\Tref{109} & 坐标离散化 & \Tref{110} & 事件排序 \\
\Tref{111} & 矩阵重塑与乘法顺序 & \Tref{112} & 稀疏向量点积 \\
\Tref{113} & 有限状态映射与倍增 & \Tref{201} & 网格 BFS \\
\Tref{202} & DFS 连通块 & \Tref{203} & 并查集 \\
\Tref{204} & Dijkstra & \Tref{205} & 拓扑排序 \\
\Tref{206} & Floyd & \Tref{207} & Kruskal \\
\Tref{208} & LCA 倍增 & \Tref{209} & 无递归树遍历与子树区间 \\
\Tref{301} & 01 背包 & \Tref{302} & 完全背包 \\
\Tref{303} & 恰好装满与至少达到 & \Tref{401} & 树状数组 \\
\Tref{402} & 线段树 & \Tref{403} & 单调栈 \\
\Tref{404} & 单调队列 & \Tref{405} & 静态区间最大值 \\
\Tref{501} & split 字符串分割 & \Tref{502} & KMP \\
\Tref{503} & Trie & \Tref{510} & 整除取整和标准取模 \\
\Tref{511} & 快速幂 & \Tref{512} & 素数和约数 \\
\Tref{513} & 进制转换 & \Tref{514} & 日期处理 \\
\Tref{515} & 质因数分解与筛素数 & \Tref{601} & 大模拟骨架 \\
\Tref{602} & 命令分发 & \Tref{603} & 状态机 \\
\Tref{604} & 懒删除优先队列 & \Tref{605} & LRU 缓存 \\
\Tref{606} & 路径规范化 & \Tref{607} & URL 参数解析 \\
\Tref{608} & 括号匹配 & \Tref{609} & 简单表达式求值 \\
\Tref{610} & 整数四则表达式解析 & -- & -- \\
\bottomrule
\end{longtable}
\endgroup

\section{C++ 语法救急}\label{sec:02-10}

\begin{longtable}{p{4.6cm}p{9.2cm}}
\toprule
需求 & 写法 \\
\midrule
\endfirsthead
\toprule
需求 & 写法 \\
\midrule
\endhead
快速输入输出 & \Code{ios::sync\_with\_stdio(false); cin.tie(nullptr);} \\
long long & 所有和、乘法、坐标平方、答案上界优先用 \Code{long long}。 \\
正整数向上取整 & $a \ge 0,b>0$ 时用 \Code{a/b+(a\%b!=0)}；有负数时查 \Tref{510}。 \\
标准非负余数 & \Code{r=x\%m; if (r<0) r+=m;}，用于负数取模下标。 \\
vector 排序 & \Code{sort(a.begin(), a.end());} \\
按结构体字段排序 & 用 \Code{sort} 加 lambda；完整写法见第\ref{ch:04}章排序小节。 \\
map 计数 & \Code{mp[x]++;} \\
set 判断存在 & \Code{if (st.count(x)) ...} \\
lower\_bound & 先排序；返回迭代器后检查是否等于 \Code{end()}。 \\
小根堆 & \Code{priority\_queue<int, vector<int>, greater<int>>}。 \\
保留小数 & \Code{cout << fixed << setprecision(k) << ans;} \\
\bottomrule
\end{longtable}

\section{最后提交前 30 秒}\label{sec:02-11}

\begin{enumerate}
  \item 数组是否开够，是否有 $n+1$、$m+1$、边界哨兵。
  \item 所有乘法、前缀和、答案是否用了 \Code{long long}。
  \item 多组数据是否清空 vector、map、队列、数组。
  \item 排序比较器是否满足严格弱序，相等时是否按题目要求。
  \item \Code{lower\_bound}、\Code{find}、map 查询是否检查不存在。
  \item 输出空格、换行、小数位是否和样例完全一致。
  \item C 题非法操作是否真的没有改状态。
  \item D/E 暴力是否至少能过样例和小数据。
\end{enumerate}

\chapter{题型详细判断树}

\section{本章定位}

本章是整本手册的导航核心。考场上不要从算法目录开始翻，而应先用本章判断：

\begin{enumerate}
  \item 这题属于哪一类？
  \item 数据范围允许什么复杂度？
  \item 题面信号对应哪个模板？
  \item 如果满分不会，有什么部分分做法？
\end{enumerate}

\Warn{判断树不是保证一眼看出满分算法的魔法。}它的作用是减少迷路：让你从题面信号稳定定位到最可能的算法和模板。

\section{考场三层分流法}

读任何题都按三层走。不要直接猜算法名。

\begin{longtable}{p{3cm}p{5.2cm}p{6.2cm}}
\toprule
层级 & 问什么 & 作用 \\
\midrule
第一层：题号和任务 & A/B/C/D/E？要求输出一个值、一组答案，还是处理很多命令？ & 决定目标：稳拿、冲分、拆模拟、拿部分分。 \\
第二层：题面结构 & 输入是数组、矩阵、图、字符串、对象、事件、树，还是混合规则？ & 决定先翻哪个章节。 \\
第三层：操作和范围 & 是否多次查询/修改？$n,m,q$ 多大？值域多大？ & 决定具体模板和复杂度。 \\
\bottomrule
\end{longtable}

\section{总判断流程}

\begin{longtable}{p{2.6cm}p{6.2cm}p{5.8cm}}
\toprule
步骤 & 问题 & 动作 \\
\midrule
1 & 这是第几题？ & A/B 先求稳，C 先拆题，D/E 先找部分分。 \\
2 & 输出形式是什么？ & 单个答案偏统计/公式；多行答案偏查询/命令；复杂文本偏解析。 \\
3 & 数据范围多大？ & 用复杂度表排除暴力或高级算法。 \\
4 & 有多次查询或修改吗？ & 查前缀和、差分、树状数组、线段树、离线处理。 \\
5 & 有图、网格、路径吗？ & 查 BFS、DFS、Dijkstra、并查集、拓扑排序。 \\
6 & 有字符串格式吗？ & 查 split、路径、URL、括号、表达式、KMP。 \\
7 & 有状态、命令、事件吗？ & 查大模拟、状态机、事件排序。 \\
8 & 有最优值、选择、容量吗？ & 查 DP、贪心、二分答案。 \\
9 & 满分算法不明确吗？ & 立即找小数据和特殊性质分。 \\
\bottomrule
\end{longtable}

\section{决策树叶子节点标准格式}

本章后面的详细决策树尽量使用同一种叶子格式。考场上每次走到叶子，都要确认 5 件事：

\begin{longtable}{p{2.8cm}p{11.2cm}}
\toprule
项目 & 必须回答的问题 \\
\midrule
题面信号 & 是什么词、输入结构或输出要求让我走到这里？ \\
数据范围 & 暴力能不能过？目标复杂度应是多少？ \\
模板入口 & 应翻第几章、哪个模板编号？ \\
易错点 & 下标、边界、溢出、排序规则、输出格式哪里最容易错？ \\
保底做法 & 满分不会时，能不能先写暴力、小数据或特殊性质版本？ \\
\bottomrule
\end{longtable}

\Warn{如果某个叶子节点无法回答“数据范围”和“模板入口”，不要立刻写代码，先回到数据范围判断树重新分流。}

\section{从根节点开始的总决策树}

下面这棵树按考场读题顺序走。先用它确定大方向，再进入 A/B/C/D/E 专项树。

\begingroup
\small
\setlength{\tabcolsep}{3pt}
\begin{longtable}{p{1.5cm}p{4.0cm}p{4.9cm}p{4.1cm}}
\toprule
节点 & 问题 & 是：走向 & 否：走向 \\
\midrule
R0 & 这是 A 题吗？ & 进入 A 题详细树，先求稳。 & 看 R1。 \\
R1 & 这是 B 题吗？ & 进入 B 题详细树，先找重复计算。 & 看 R2。 \\
R2 & 这是 C 题吗？ & 进入 C 题详细树，先建模。 & 看 R3。 \\
R3 & 这是 D/E 题吗？ & 先看子任务，进入部分分树。 & 回题面确认题号。 \\
R4 & 输出只有一个数或少量值吗？ & 优先统计、公式、排序、DP。 & 看是否多次查询或命令。 \\
R5 & 每次输入操作后都要输出吗？ & 命令分发、数据结构、状态机。 & 可能是最后统一输出的模拟/预处理。 \\
R6 & 有 $q$ 次查询/修改且 $q$ 大吗？ & 前缀、差分、树状数组、线段树、离线。 & 暴力或简单模拟可能可行。 \\
R7 & 有“最小/最大/阈值/最少资源”吗？ & 检查单调性，可能二分答案或 DP。 & 看图、字符串、状态。 \\
R8 & 有“路径/连通/最短/依赖”吗？ & 图论或拓扑。 & 看字符串和大模拟。 \\
R9 & 有复杂格式、路径、括号、表达式吗？ & 字符串解析、栈、递归解析。 & 普通数组/矩阵/对象模拟。 \\
\bottomrule
\end{longtable}
\endgroup

\section{详细叶子索引总表}

这张表是全局的“最终落点”索引。专项树走不清楚时，直接查这一张。

\begingroup
\small
\setlength{\tabcolsep}{3pt}
\begin{longtable}{p{2.8cm}p{3.2cm}p{3.1cm}p{3.0cm}p{2.3cm}}
\toprule
叶子 & 题面信号 & 数据范围判断 & 模板入口 & 保底做法 \\
\midrule
L01 统计扫描 & 求和、计数、最大最小 & $O(n)$ 可过 & A 题、T106 & 一遍循环。 \\
L02 排序排名 & 排名、中位数、多关键字 & $O(n\log n)$ 可过 & T105 & 先按单关键字过样例。 \\
L03 一维区间 & 多次区间和/数量 & 静态查询多 & T101 & 每次暴力扫小数据。 \\
L04 二维矩阵 & 子矩阵、邻域、图像 & $nm$ 可扫，多次查询用预处理 & T102/T104 & 小矩阵暴力。 \\
L05 哈希查重 & 是否出现、频率、点集 & 值域大或字符串键 & T106 & map 替代 unordered\_map。 \\
L06 坐标离散 & 坐标大、点少 & 不能直接开数组 & T109 & map 坐标。 \\
L07 二分答案 & 最小时间、最大阈值 & check 单调且快 & T107 & 枚举小范围答案。 \\
L08 背包选择 & 选若干、容量、金额 & 总容量可 DP & T301/T302 & DFS 小数据。 \\
L09 BFS & 最短步数、边权 1 & 状态数可控 & T201 & 小图暴力搜索。 \\
L10 Dijkstra & 正权最短路 & $n,m$ 大，边权非 1 & T204 & 小图 Floyd。 \\
L11 拓扑 & 前置、依赖、逻辑门 & DAG 或需判环 & T205 & 小图反复找入度 0。 \\
L12 并查集 & 合并集合、连通性 & 只合并查连通 & T203 & BFS/DFS 小数据。 \\
L13 事件模拟 & 按时间发生、到期 & 事件数可排序 & T110 & 按输入顺序模拟小数据。 \\
L14 状态机 & 对象状态变化 & 状态少且规则明确 & T603 & if/else 先过样例。 \\
L15 路径/URL & 路径、参数、配置 & 总长度可线性 & T606/T607 & 手写 split。 \\
L16 表达式 & 括号、优先级、化学式 & 嵌套深度可递归 & T608/T609 & 暴力解析小表达式。 \\
L17 树上路径 & 祖先、子树、路径 & 查询多需预处理 & T208/T402 & 每次 DFS。 \\
L18 动态区间 & 修改和查询交替 & $q$ 大 & T401/T402 & 修改少就重算。 \\
L19 大模拟 & 长题面、多对象规则 & 规则多于算法 & T601/T602 & 分命令逐步拿分。 \\
L20 D/E 小数据 & 子任务 $n$ 小或特殊性质 & 暴力能过一档 & 第 8 章 & 先交可运行版本。 \\
\bottomrule
\end{longtable}
\endgroup

\section{题面结构第一判断}

\begingroup
\small
\setlength{\tabcolsep}{3pt}
\begin{longtable}{p{3.1cm}p{4.1cm}p{4.2cm}p{3.0cm}}
\toprule
输入主体 & 题面信号 & 第一反应 & 翻到哪里 \\
\midrule
一串数字 / 数组 & 求和、最大、最小、相邻、分段、前缀 & 扫描、排序、前缀和 & A/B、T101、T105 \\
二维矩阵 & 图像、网格、矩形、邻域、旋转、重塑 & 下标转换、二维前缀、矩阵模拟 & A/B、T102、T104 \\
很多字符串行 & 路径、URL、JSON、Markdown、命令、表达式 & 字符串解析、大模拟 & C、T501、T606、T607 \\
点和边 & 连通、最短路、依赖、拓扑、最小代价 & 图论模板 & T201--T207 \\
树 & 子树、祖先、路径、删点、合并 & DFS、LCA、树形 DP、DFS 序 & C/D/E、T208 \\
多条操作 & 修改、查询、插入、删除、回滚 & 数据结构或离线处理 & T103、T401、T402 \\
时间顺序事件 & 某时刻发生、到期、优先级、调度 & 事件排序、状态机、堆 & C、T110、T603、T604 \\
选择若干物品 & 容量、费用、价值、至少/至多 & 背包 DP 或贪心 & T301、T302 \\
答案最优值 & 最小时间、最大阈值、最小最大 & 二分答案或 DP & T107、D/E \\
\bottomrule
\end{longtable}
\endgroup

\section{数据范围判断树}

\subsection{单变量规模}

\begin{longtable}{p{3.4cm}p{5.4cm}p{5.4cm}}
\toprule
范围 & 可考虑 & 不宜考虑 \\
\midrule
$n \le 10$ & 全排列、爆搜、手写模拟所有情况 & 不必急着写复杂优化 \\
$n \le 20$ & 子集枚举、状压 DP、折半搜索 & $O(n!)$ 仍危险 \\
$n \le 100$ & $O(n^3)$、Floyd、区间 DP & 指数算法 \\
$n \le 1000$ & $O(n^2)$、简单 DP、稠密图 & $O(n^3)$ 需谨慎 \\
$n \le 10^5$ & $O(n\log n)$、$O(n)$ & $O(n^2)$ \\
$n \le 10^6$ & $O(n)$、计数数组、简单筛法 & 多层 map 或重型结构 \\
$n \ge 10^7$ & 只考虑线性或数学公式 & 大数组套多维结构 \\
\bottomrule
\end{longtable}

\subsection{多变量规模}

\begin{longtable}{p{4.2cm}p{5.2cm}p{5.0cm}}
\toprule
范围组合 & 判断 & 常见模板 \\
\midrule
$n,m \le 500$ 且矩阵 & $O(nm)$ 或 $O(nm\log nm)$ 通常可行 & 二维扫描、T102、T104 \\
$n,m \le 2000$ 且矩阵 & 尽量 $O(nm)$，少做多次全图扫描 & 二维前缀、BFS \\
$q \le 10^5$ & 每次操作必须快 & 前缀、差分、树状数组、线段树 \\
值域 $\le 10^6$ & 可以开计数数组 & 计数、桶、前缀计数 \\
值域很大但数量少 & 不可直接开值域数组 & map、unordered\_map、离散化 \\
边数 $m \approx n$ & 稀疏图 & 邻接表、BFS、Dijkstra 堆优化 \\
边数 $m \approx n^2$ 且 $n$ 小 & 稠密图 & Floyd、邻接矩阵 \\
字符串总长 $\le 10^6$ & 按总长度线性处理 & split、KMP、Trie \\
\bottomrule
\end{longtable}

\Warn{数据范围优先级高于题面关键词。}例如题面写“最短路”，但 $n \le 100$ 且多源查询，可以用 Floyd；若 $n,m \le 10^5$，通常要 Dijkstra 或 BFS。

\section{关键词到模板总索引}

\begingroup
\small
\setlength{\tabcolsep}{3pt}
\begin{longtable}{p{3.5cm}p{3.8cm}p{4.2cm}p{2.9cm}}
\toprule
题面关键词 & 判断 & 模板 & 章节 \\
\midrule
多次区间和 & 静态区间查询 & T101 一维前缀和 & B / 模板库 \\
子矩阵和、邻域均值 & 二维静态查询 & T102 二维前缀和 & B / 模板库 \\
区间加，最后查询 & 离线修改 & T103 一维差分 & B / 模板库 \\
矩形加，最后输出 & 二维离线修改 & T104 二维差分 & B / 模板库 \\
排名、优先级、多关键字 & 比较器排序 & T105 排序比较器 & A/B/C \\
出现次数、频率、查重 & 统计 & T106 哈希统计 & A/B \\
最小化最大值、最大化最小值 & 二分答案 & T107 二分答案 & B/D \\
排序后找两数、合并稀疏向量 & 双指针 & T108 双指针 & B \\
坐标很大、点不多 & 离散化 & T109 坐标离散化 & B/D \\
按时间发生、到期、归还 & 事件模拟 & T110 事件排序 & B/C \\
最短步数、迷宫、扩散 & BFS & T201 网格 BFS & B/C \\
连通块、区域大小 & DFS/BFS & T202 DFS 连通块 & B/C \\
合并集合、同组、连通性 & 并查集 & T203 并查集 & B/D \\
正权最短路 & Dijkstra & T204 Dijkstra & B/D \\
依赖、先后顺序、任务计划 & 拓扑排序 & T205 拓扑排序 & B/D \\
任意两点最短路，小图 & Floyd & T206 Floyd & B/D \\
最小连接代价 & Kruskal & T207 Kruskal & D/E \\
树上祖先、路径 & LCA & T208 LCA & D/E \\
选或不选、容量 & 01 背包 & T301 01 背包 & B/D \\
可重复选择 & 完全背包 & T302 完全背包 & B/D \\
单点修改区间和 & 树状数组 & T401 树状数组 & D/E \\
区间修改查询 & 线段树 & T402 线段树 & D/E \\
最近更大/更小 & 单调栈 & T403 单调栈 & B/D \\
滑动窗口极值 & 单调队列 & T404 单调队列 & B/D \\
路径、配置分割 & split & T501 split & C \\
长串匹配 & KMP & T502 KMP & D/E \\
前缀查询 & Trie & T503 Trie & D/E \\
分组、分页、负数取模 & 整除取整 & T510 整除取整和标准取模 & A/B \\
快速幂、取模幂 & 数学 & T511 快速幂 & B/D \\
素数、约数、因子 & 数学 & T512 素数约数 & B/D \\
进制转换、编码 & 数字处理 & T513 进制转换 & A/B/C \\
日期、时间、周期 & 日期处理 & T514 日期处理 & A/C \\
长题面、多命令 & 大模拟 & T601 大模拟骨架 & C \\
每行操作名 & 命令分发 & T602 命令分发 & C \\
多状态对象 & 状态机 & T603 状态机 & C \\
动态删除仍取最优 & 懒删除堆 & T604 懒删除优先队列 & C/D \\
最近最少使用 & LRU & T605 LRU 缓存 & C \\
目录路径 & 路径规范化 & T606 路径规范化 & C \\
URL 参数 & URL 解析 & T607 URL 参数解析 & C \\
括号合法性 & 栈 & T608 括号匹配 & C \\
简单表达式 & 扫描求值 & T609 表达式入门 & C \\
\bottomrule
\end{longtable}
\endgroup

\section{操作模式判断树}

\subsection{只有查询，没有修改}

\begin{longtable}{p{4.6cm}p{5.4cm}p{4.0cm}}
\toprule
查询内容 & 优先算法 & 模板 \\
\midrule
区间和、区间数量 & 前缀和 / 前缀计数 & T101 \\
子矩阵和、邻域平均 & 二维前缀和 & T102 \\
最大最小，静态 & 排序 / 预处理 / 单调结构 & T105 / T403 \\
图上单源距离 & BFS / Dijkstra & T201 / T204 \\
任意两点距离，小图 & Floyd & T206 \\
字符串是否出现 & find / KMP & T502 \\
子树信息，静态 & DFS 序 + 前缀 / 树状数组 & T208 / T401 \\
\bottomrule
\end{longtable}

\subsection{有修改，有查询}

\begin{longtable}{p{4.6cm}p{5.4cm}p{4.0cm}}
\toprule
操作形态 & 优先算法 & 模板 \\
\midrule
单点修改，区间和查询 & 树状数组 & T401 \\
区间修改，单点查询 & 差分思想 / 树状数组 & T103 / T401 \\
区间修改，区间查询 & 线段树 & T402 \\
动态插入删除，查最大最小 & set / 堆 / 懒删除 & T604 \\
修改很少 & 每次重算 & 部分分策略 \\
查询很少 & 每次暴力 & 部分分策略 \\
操作可以离线 & 排序事件、差分、并查集倒序 & T110 / T103 / T203 \\
\bottomrule
\end{longtable}

\subsection{多轮模拟}

\begin{longtable}{p{4.6cm}p{5.4cm}p{4.0cm}}
\toprule
模拟形态 & 判断重点 & 模板 \\
\midrule
每轮所有对象同时变化 & 先存变化量，再统一更新 & T601 / T603 \\
事件按时间发生 & 先排序，时间相同按题目优先级 & T110 \\
状态会过期 & 每次操作前处理过期，或懒处理 & T603 / T604 \\
每次选最优对象 & priority\_queue / set & T604 \\
命令种类多 & 每种命令写独立函数 & T602 \\
\bottomrule
\end{longtable}

\section{A 题快速判断树}

A 题目标是稳。判断时按下面顺序，不要一开始就想复杂算法。

\begin{enumerate}
  \item 是否只要求一个数？优先一遍扫描、公式、计数。
  \item 是否有排序、排名、最大最小？用 \Code{sort} 后处理。
  \item 是否有字符或字符串？按字符扫描，先别写复杂解析。
  \item 是否有二维数组？先画 $2 \times 3$ 小例子确认下标。
  \item 是否有日期、时间、进制、取整、小数？翻对应模板，不临场硬想语法。
\end{enumerate}

\begin{longtable}{p{4.6cm}p{5.4cm}p{4.0cm}}
\toprule
题面信号 & 做法 & 模板/章节 \\
\midrule
求和、计数、最大最小 & 一遍扫描 & A 题统计 \\
出现次数、直方图 & 计数数组 / map & T106 \\
小数、百分比、平均、方差 & 公式 + \Code{setprecision} & A 题公式 \\
分组、分页、向上取整、负数取模 & 整数公式，不用浮点硬凑 & T510 \\
排序、排名、中位数 & sort & T105 \\
字符串、字符、校验码 & string 扫描 & T501 \\
进制、数位、编码 & 进制转换或取模拆位 & T513 \\
日期、星期、红绿灯周期 & 闰年、月份、取模 & T514 \\
坐标距离、矩形面积 & 公式，注意边界 & A 题坐标 \\
二维矩阵变换 & 写原坐标到新坐标公式 & A 题二维表格 \\
\bottomrule
\end{longtable}

\subsection{A 题详细叶子表}

\begingroup
\small
\setlength{\tabcolsep}{3pt}
\begin{longtable}{p{2.8cm}p{3.0cm}p{3.0cm}p{3.0cm}p{2.6cm}}
\toprule
叶子 & 识别信号 & 先写什么 & 易错点 & 翻到哪里 \\
\midrule
A1 纯统计 & 求总和、数量、最大最小 & 一遍扫描 & 初值、long long & 第 5 章统计 \\
A2 频率统计 & 出现次数、直方图 & 数组计数或 map & 值域是否能开数组 & T106 \\
A3 排序输出 & 第几大、排名、中位数 & sort 后取值 & 下标从 0 还是 1 & T105 \\
A4 公式计算 & 平均、方差、加权、现值、分组 & 按公式翻译 & 整数除法、取整、小数位 & 第 5 章公式 / T510 \\
A5 字符处理 & 校验码、字符类别、字符串计数 & string 扫描 & getline 和 cin 混用 & T501 \\
A6 数位/进制 & 数位和、编码、进制分解 & 取模除法 & 0、负数、long long & T513 \\
A7 日期时间 & 星期、天数、周期灯 & 月份表 + 闰年 & 闭区间、取模 & T514 \\
A8 二维表格 & 图像、矩阵、邻居平均 & 画小矩阵找下标 & 行列反、边界 & 第 5 章二维 \\
A9 坐标几何 & 距离、矩形面积、点位置 & 套公式 & 交集为负取 0 & 第 5 章坐标 \\
A10 简单状态 & 连续得分、开关、当前状态 & 变量记录状态 & 状态重置时机 & 第 5 章模拟 \\
\bottomrule
\end{longtable}
\endgroup

\section{B 题快速判断树}

B 题目标是尽量拿满。它常常不是“高级算法”，而是把暴力换成预处理。

\begin{enumerate}
  \item 暴力做法是什么？每次是否重复算同一段、同一块、同一类？
  \item 重复算区间：前缀和。重复改区间：差分。
  \item 重复查存在：set/map/hash。
  \item 排序后规律明显：排序扫描、双指针、二分。
  \item 有最小时间/最大阈值：先问答案是否单调。
\end{enumerate}

\begin{longtable}{p{4.6cm}p{5.4cm}p{4.0cm}}
\toprule
题面信号 & 做法 & 模板 \\
\midrule
多次区间/矩阵查询 & 前缀和 & T101 / T102 \\
区间修改最后输出 & 差分 & T103 / T104 \\
频率统计、查重、点集匹配 & 哈希 / map / set & T106 \\
排序后扫描、阈值统计 & sort + 扫描 & T105 \\
两个有序序列匹配 & 双指针 & T108 \\
答案具有单调性 & 二分答案 & T107 \\
最短步数、扩散 & BFS & T201 \\
连通关系、合并集合 & 并查集 & T203 \\
简单依赖、任务计划 & 拓扑排序 & T205 \\
背包信号、至少达到某值 & DP & T301 / T302 \\
质因数、约数、幂 & 数学模板 & T511 / T512 \\
\bottomrule
\end{longtable}

\subsection{B 题详细叶子表}

\begingroup
\small
\setlength{\tabcolsep}{3pt}
\begin{longtable}{p{2.6cm}p{3.2cm}p{3.2cm}p{3.0cm}p{2.6cm}}
\toprule
叶子 & 暴力形态 & 优化判断 & 模板 & 保底 \\
\midrule
B1 一维前缀 & 每次扫区间 & 只查不改，多次区间和 & T101 & 暴力扫小数据 \\
B2 二维前缀 & 每次扫矩形 & 矩阵静态，多次子矩阵 & T102 & 暴力双循环 \\
B3 一维差分 & 每次区间逐个加 & 区间修改，最后输出/查询少 & T103 & 直接修改小数据 \\
B4 二维差分 & 每次矩形逐格加 & 矩形覆盖，最后统计 & T104 & 二维暴力 \\
B5 哈希统计 & 每次线性找存在 & 查重、频率、坐标点存在 & T106 & map 先写稳 \\
B6 排序扫描 & 枚举每个阈值再统计 & 排序后相邻段贡献可更新 & T105/T101 & 枚举阈值小数据 \\
B7 双指针 & 双重循环匹配 & 两序列有序或可排序 & T108 & map 匹配 \\
B8 二分答案 & 枚举答案 & check 单调，答案范围大 & T107 & 枚举小范围 \\
B9 简单 BFS & 枚举路径 & 无权最短步数/扩散 & T201 & DFS 小图 \\
B10 并查集 & 每次 DFS 判连通 & 只有合并和查连通 & T203 & 邻接表 DFS \\
B11 拓扑 & 反复扫描找可做任务 & 前置依赖固定 & T205 & 小 n 反复扫描 \\
B12 背包 & 枚举所有选择 & 选或不选，容量不太大 & T301 & DFS 子集 \\
B13 基础数学 & 枚举因子到 n & 质因数、幂、约数 & T511/T512 & 小范围试除 \\
B14 矩阵优化 & 按题意多重循环超时 & 调整乘法/扫描顺序 & 第 6 章矩阵 & 暴力部分分 \\
B15 时间区间 & 每分钟逐个标记 & 时间范围不大或可差分 & T103/T110 & 数组标记小范围 \\
\bottomrule
\end{longtable}
\endgroup

\section{C 题快速判断树}

C 题最怕一边读一边写。正确流程是先拆对象、状态、命令、规则。

\subsection{C 题读题分解}

\begin{longtable}{p{3.0cm}p{6.0cm}p{5.0cm}}
\toprule
要找什么 & 在题面中怎么标记 & 代码落点 \\
\midrule
对象 & 用户、文件、节点、任务、请求、设备、词条 & struct \\
状态 & 空闲/占用、有效/过期、允许/拒绝、在线/离线 & enum 或 int \\
命令 & create、delete、query、request、release & T602 命令分发 \\
时间 & 到期、当前时刻、持续时间、周期 & T110、T514 \\
规则优先级 & 先判断 A，再判断 B；相同时间谁优先 & 独立函数，写注释 \\
输出 & 每次命令输出、最后统一输出、格式化文本 & 先写输出样例函数 \\
\bottomrule
\end{longtable}

\subsection{C 题模板选择}

\begin{longtable}{p{4.6cm}p{5.4cm}p{4.0cm}}
\toprule
题面信号 & 做法 & 模板 \\
\midrule
长题面、多对象、多规则 & 大模拟拆解 & T601 \\
每行一个命令 & 命令分发 & T602 \\
对象状态变化 & 状态机 & T603 \\
按时间发生 & 事件排序 & T110 \\
每次取最优任务 & 优先队列 / 懒删除 & T604 \\
路径、URL、配置 & 字符串解析 & T501 / T606 / T607 \\
括号、表达式、方程式 & 栈 / 递归解析 / 扫描 & T608 / T609 \\
缓存淘汰 & list + map 或队列 & T605 \\
地图带状态 & 状态 BFS & C 题章节 \\
树或目录结构 & DFS、父子 map、路径查找 & T606 / T208 \\
\bottomrule
\end{longtable}

\subsection{C 题详细叶子表}

\begingroup
\small
\setlength{\tabcolsep}{3pt}
\begin{longtable}{p{2.7cm}p{3.2cm}p{3.2cm}p{3.0cm}p{2.6cm}}
\toprule
叶子 & 题面信号 & 代码结构 & 易错点 & 模板 \\
\midrule
C1 命令分发 & 每行操作名不同 & 主循环读 op，分函数处理 & 参数数量不同 & T602 \\
C2 状态机 & 对象有状态变化 & 状态字段 + 转移函数 & 同一操作多次转移 & T603 \\
C3 事件排序 & 时间、到期、归还、同时发生 & Event 结构体排序 & 同时事件优先级 & T110 \\
C4 优先队列 & 每次取最优任务/对象 & heap 或 set & 失效元素删除 & T604 \\
C5 路径目录 & 文件、目录、父子、配额 & 树节点 + path map & 创建失败回滚 & T606 \\
C6 URL/参数 & 路径、查询参数、匹配规则 & split + 参数 map & 空参数、类型检查 & T607 \\
C7 JSON/配置 & 嵌套键值、查询路径 & 栈/递归建结构 & 转义和层级 & T501 \\
C8 Markdown/文本 & 多行文本、块、行内格式 & 先分块再处理行 & 空行分块 & T601 \\
C9 表达式/括号 & 运算、括号、化学式 & 栈或递归解析 & 优先级、嵌套倍数 & T608/T609 \\
C10 编码解码 & 十六进制、压缩、校验 & 指针扫描字符串 & 长度字段越界 & T513 \\
C11 矩阵图像 & 固定矩阵、扫描顺序、变换 & 二维数组 + 顺序表 & 行列、取整 & T102 \\
C12 权限规则 & 用户、角色、资源、权限 & set/map 规则判断 & 通配符、最大等级 & T603 \\
C13 带状态 BFS & 地图、钥匙、方向、时间 & visited 包含状态 & 状态维度漏掉 & T201 \\
C14 树/目录模拟 & 子树、删除、查询、合并 & DFS + 父子关系 & 删除后引用失效 & T208 \\
C15 分层拿分 & 规则太多写不完 & 先实现核心命令 & 输出格式别乱 & T601 \\
\bottomrule
\end{longtable}
\endgroup

\section{逐句读题分支树}

考场读题时，建议拿笔在题面旁边标记下面六类词。标完后再决定模板。

\begin{longtable}{p{2.8cm}p{5.8cm}p{5.4cm}}
\toprule
标记 & 看到什么就圈出来 & 后续判断 \\
\midrule
对象 & 学生、用户、文件、节点、仓库、服务器、任务、词条 & 是否需要 struct？是否有编号到对象的 map？ \\
属性 & 权重、时间、状态、坐标、容量、权限、父节点 & 是否要排序、比较器、状态机？ \\
动作 & 查询、修改、创建、删除、移动、合并、授权、到期 & 是否是命令分发或数据结构？ \\
范围 & $n$、$m$、$q$、值域、字符串总长、矩阵大小 & 直接排除不可能复杂度。 \\
关系 & 相邻、包含、依赖、祖先、连通、路径、前缀 & 图、树、拓扑、字符串、区间？ \\
目标 & 最大、最小、数量、是否存在、方案数、输出过程 & 统计、二分、DP、搜索还是模拟？ \\
\bottomrule
\end{longtable}

\subsection{从句子到模板}

\begin{longtable}{p{4.8cm}p{5.2cm}p{4.2cm}}
\toprule
题面句子模式 & 立即追问 & 模板方向 \\
\midrule
“给定 $q$ 次询问” & 每次询问是否只读不改？ & 前缀和 / 数据结构 \\
“每次修改后回答” & 修改是点、区间还是对象状态？ & 树状数组 / 线段树 / 状态机 \\
“求满足条件的数量” & 条件能否拆成前缀、排序或哈希？ & 计数 / 排序扫描 / map \\
“求最小的 $x$ 使得……” & $x$ 越大是否越容易满足？ & 二分答案 \\
“选择若干个” & 每个物品选一次还是多次？ & 01 背包 / 完全背包 \\
“从 A 到 B 的最短……” & 边权是否全为 1？ & BFS / Dijkstra \\
“任务必须在前置任务后” & 是否有环？ & 拓扑排序 \\
“字符串由若干规则组成” & 是否有括号或嵌套？ & split / 栈 / 递归解析 \\
“路径包含 . 和 ..” & 是否需要规范化？ & T606 \\
“状态在某时刻过期” & 过期是在操作前还是操作后处理？ & 状态机 / 事件模拟 \\
\bottomrule
\end{longtable}

\section{B 题专项分支树}

B 题常见陷阱是“暴力能写但会超时”。先写出暴力，再问它重复算了什么。

\subsection{B 题从暴力到优化}

\begin{longtable}{p{4.4cm}p{5.2cm}p{4.4cm}}
\toprule
暴力形态 & 重复在哪里 & 优化方向 \\
\midrule
每次询问都扫一段数组 & 区间和/数量重复计算 & 一维前缀和 T101 \\
每次询问都扫一个矩形 & 子矩阵重复计算 & 二维前缀和 T102 \\
每次区间加都逐个改 & 区间修改重复 & 差分 T103/T104 \\
每次判断点是否存在都扫所有点 & 存在性查询重复 & set / unordered\_set \\
每次找匹配点都双重循环 & 坐标或编号匹配重复 & map / 双指针 / 排序 \\
每个候选答案都完整模拟 & 答案范围大 & 二分答案 T107 \\
每次重新排序 & 只需维护极值或有序集合 & priority\_queue / set \\
每次从头判断依赖任务 & 依赖结构固定 & 拓扑排序 T205 \\
\bottomrule
\end{longtable}

\subsection{B 题二分答案确认表}

\begin{longtable}{p{4.6cm}p{4.8cm}p{4.6cm}}
\toprule
问题 & 是 & 否 \\
\midrule
答案是一个整数或实数吗？ & 继续问单调性 & 多半不是二分答案。 \\
给定答案 $x$，能快速判断可不可行吗？ & 写 \Code{check(x)} & 不适合二分。 \\
$x$ 变大后，可行性只朝一个方向变化吗？ & 可以二分 & 不能二分。 \\
题目问最小时间/最大阈值/最少资源吗？ & 优先尝试二分 & 可能是 DP/贪心。 \\
\bottomrule
\end{longtable}

\section{C 题专项分支树}

C 题不是先想算法，而是先建模。建模错误比模板不会更致命。

\subsection{C 题对象建模表}

\begin{longtable}{p{3.6cm}p{5.2cm}p{5.0cm}}
\toprule
题面对象 & 常见字段 & 容器选择 \\
\midrule
用户 / 角色 / 权限 & 名称、权限集合、等级、资源类型 & \Code{map<string, Obj>}，\Code{set<string>} \\
文件 / 目录 & 名称、父节点、子节点、大小、配额 & map 路径到节点，或树节点数组 \\
服务器 / 节点 & 编号、算力、任务数、可用状态 & vector + set/priority\_queue \\
任务 / 请求 & 时间、优先级、资源需求、状态 & struct + sort/heap \\
词条 / token & 字符串、权重、出现位置 & map/unordered\_map + heap \\
地图状态 & 坐标、方向、钥匙、时间、能量 & queue + visited 多维数组/map \\
\bottomrule
\end{longtable}

\subsection{C 题规则优先级表}

\begin{longtable}{p{4.6cm}p{5.0cm}p{4.6cm}}
\toprule
题面出现 & 必须单独确认 & 常见错误 \\
\midrule
“若同时发生” & 同时事件的排序规则 & 先后顺序写反。 \\
“如果不存在则……” & 缺省值或失败输出 & map 自动创建误判存在。 \\
“到期” & 到期时刻是否还有效 & $<$ 和 $\le$ 写错。 \\
“优先选择” & 比较器所有关键字 & 忘记编号升序等最后规则。 \\
“忽略非法操作” & 是否改变状态、是否输出 & 非法时仍修改了数据。 \\
“回滚/失败” & 前面临时改动要撤销 & 配额、路径创建类常见。 \\
\bottomrule
\end{longtable}

\subsection{C 题字符串解析分支}

\begin{longtable}{p{4.6cm}p{5.2cm}p{4.2cm}}
\toprule
字符串形态 & 判断 & 模板 \\
\midrule
固定分隔符，无嵌套 & split 即可 & T501 \\
路径，含 \Code{.}、\Code{..}、多个斜杠 & 栈式规范化 & T606 \\
URL，含路径和查询参数 & 先按 \Code{?}，再按 \Code{\&} & T607 \\
括号嵌套 & 栈或递归 & T608 \\
有运算优先级 & 两栈或递归下降 & T609 \\
化学式/表达式含倍数 & 递归解析返回 map & C 题表达式扩展 \\
多行文本块 & 先分块，再行内解析 & T601 \\
\bottomrule
\end{longtable}

\section{D/E 题快速判断树}

D/E 题不要先问“我会不会满分”，先问“我能不能拿部分分”。

\begin{enumerate}
  \item 找子任务数据范围：$n \le 20$、链、树、无修改、边权 1、查询少。
  \item 写最朴素暴力，估算它能过哪档。
  \item 如果暴力只有 0 分，再找一个最接近的模板：BFS、Dijkstra、Kruskal、树状数组、线段树、DP。
  \item 优先实现你能调通的版本，不把时间全部压在复杂满分做法上。
\end{enumerate}

\begin{longtable}{p{4.6cm}p{5.2cm}p{4.2cm}}
\toprule
题面信号 & 满分方向 & 部分分方向 \\
\midrule
$n \le 20$ 子任务 & 状压 / DFS & 枚举子集 \\
树上路径 & LCA / 树上差分 & 每次 DFS 找路径 \\
区间修改查询 & 线段树 & 暴力 / 树状数组 / 差分 \\
最小生成代价 & Kruskal & 小图枚举或直接 Kruskal \\
字符串匹配 & KMP / Trie & find / 暴力匹配 \\
正权最短路 & Dijkstra & 小图 Floyd / 无权 BFS \\
最大化最小值 & 二分答案 & 枚举答案小范围 \\
复杂 DP & 优化 DP & DFS / 小状态 DP \\
动态集合取最值 & 平衡树 / 堆懒删除 & 每次排序或线性扫描 \\
\bottomrule
\end{longtable}

\section{D/E 部分分决策树}

下面这棵树只服务一件事：不会满分时，尽快产出能提交的代码。

\begin{enumerate}
  \item 先看子任务表，有没有 $n \le 20$、$n \le 100$、链、树、无修改、边权 1、查询少。
  \item 如果有小 $n$：优先 DFS、枚举、Floyd、$O(n^2)$ DP。
  \item 如果有特殊图：边权 1 用 BFS，树用 DFS，DAG 用拓扑。
  \item 如果有特殊操作：无修改用前缀，有修改少就每次重算，查询少就每次暴力。
  \item 如果只有大数据：选择你最熟的相近模板，写 30 到 50 分版本。
\end{enumerate}

\begin{longtable}{p{4.4cm}p{5.2cm}p{4.4cm}}
\toprule
子任务/特殊性质 & 可写版本 & 常用模板 \\
\midrule
$n \le 20$ & 枚举子集、DFS、状压 & 状压 DP、搜索 \\
$n \le 100$ & Floyd、$O(n^3)$ DP、暴力查询 & T206 \\
树退化成链 & 数组/前缀处理路径 & T101 / T103 \\
没有修改 & 预处理后回答查询 & 前缀、DFS 序 \\
查询次数很少 & 每次 DFS/BFS/扫描 & 暴力部分分 \\
边权全为 1 & BFS 替代 Dijkstra & T201 \\
图是 DAG & 拓扑 DP & T205 \\
值域小 & 计数数组/桶 & T106 \\
参数 $k$ 很小 & 状态压缩或按 $k$ 枚举 & 状压入口 \\
\bottomrule
\end{longtable}

\subsection{D/E 题详细叶子表}

\begingroup
\small
\setlength{\tabcolsep}{3pt}
\begin{longtable}{p{2.6cm}p{3.2cm}p{3.2cm}p{3.0cm}p{2.8cm}}
\toprule
叶子 & 识别信号 & 满分方向 & 部分分写法 & 翻到哪里 \\
\midrule
DE1 小 n & $n \le 20$ 或参数小 & 状压/搜索优化 & 枚举子集、DFS & 第 8 章状压 \\
DE2 小图 & $n \le 100$ & Floyd/DP & Floyd 或暴力建图 & T206 \\
DE3 树路径 & 树上多次路径 & LCA/树剖/树差分 & 每次 DFS 找路径 & T208 \\
DE4 子树查询 & 子树加、子树查 & DFS 序 + 树状数组/线段树 & 每次遍历子树 & T401/T402 \\
DE5 动态区间 & 区间改查交替 & 线段树 & 小范围数组模拟 & T402 \\
DE6 动态集合 & 插入删除查极值 & set/堆懒删除 & 每次排序/扫描 & T604 \\
DE7 最短路 & 正权图路径 & Dijkstra/建图优化 & Floyd 小图、BFS 边权1 & T204 \\
DE8 最小生成 & 连通最小代价 & Kruskal/Prim & 直接 Kruskal & T207 \\
DE9 字符串匹配 & 长串、多模式、前缀 & KMP/Trie/AC 自动机 & find 暴力 & T502/T503 \\
DE10 复杂 DP & 最优方案数、限制多 & 状态设计/优化 & 记忆化 DFS 小数据 & 第 8 章 DP \\
DE11 数学计数 & 组合、取模、约数 & 预处理/数论 & 枚举小范围 & T511/T512 \\
DE12 离线处理 & 操作可重排、查询很多 & 排序/分治/离线 DSU & 按输入暴力 & T110/T203 \\
DE13 几何 & 点线面、距离、交叉 & 计算几何 & 小数据枚举 & 后续补强 \\
DE14 网络流/匹配 & 分配、匹配、容量流 & 网络流/二分图 & 贪心或小图 DFS & 放弃满分 \\
\bottomrule
\end{longtable}
\endgroup

\section{容易误判的分支}

\begin{longtable}{p{4.5cm}p{4.8cm}p{4.8cm}}
\toprule
看起来像 & 但如果题面出现 & 应改判为 \\
\midrule
简单模拟 & $q$ 很大，每次都查区间 & 前缀和 / 差分 / 数据结构 \\
排序题 & 要动态插入删除后仍排名 & set、树状数组、线段树 \\
字符串题 & 有括号嵌套、表达式优先级 & 栈或递归下降 \\
图搜索题 & 边权不是 1 & Dijkstra，不是普通 BFS \\
BFS 题 & 状态包含钥匙、方向、时间、能量 & 状态 BFS \\
贪心题 & 选或不选、容量限制明显 & 背包 DP \\
二分题 & check 不单调 & 不能二分，改 DP/贪心/扫描 \\
树题 & 只有父子目录和路径字符串 & 可能是大模拟，不一定是 LCA \\
数学题 & 有多次询问和大范围 & 预处理、筛法、快速幂 \\
\bottomrule
\end{longtable}

\section{无法判断时的保底流程}

如果读完题仍不知道算法：

\begin{enumerate}
  \item 写出最暴力做法。
  \item 估算暴力复杂度。
  \item 看数据范围中有没有小数据能过。
  \item 找特殊性质：无修改、链、树、边权 1、参数小。
  \item 先实现能过样例和小数据的版本。
  \item 再尝试从本章关键词表找优化方向。
\end{enumerate}

\section{判断树反向测试表}

刷完一道题后，用这张表检查第 3 章是否真的帮你判断。如果某行填不上，就说明手册还要继续补。

\begin{longtable}{p{4.0cm}p{5.0cm}p{5.0cm}}
\toprule
问题 & 例子 & 应能定位到 \\
\midrule
题面主体是什么？ & 数组、矩阵、图、字符串、对象、事件 & 题面结构第一判断 \\
是否有多次操作？ & $q$ 次查询、每次修改、每次输出 & 操作模式判断树 \\
暴力复杂度是多少？ & $O(nq)$、$O(n^2)$、$O(nm)$ & 数据范围判断树 \\
关键词是否误导？ & 看起来像模拟但 $q$ 很大 & 容易误判分支 \\
满分不会时能拿什么？ & 小数据、特殊性质、暴力 & D/E 部分分决策树 \\
对应模板编号是什么？ & T101、T107、T603 等 & 关键词到模板总索引 \\
\bottomrule
\end{longtable}

\section{最后 5 分钟判断}

\begin{longtable}{p{4.2cm}p{9.8cm}}
\toprule
情况 & 处理 \\
\midrule
A/B 还没过样例 & 放下后面题，先修 A/B。 \\
C 写了一半但样例不过 & 注释掉高风险优化，保留最基础模拟。 \\
D/E 没思路 & 找最小子任务，写暴力或特殊性质。 \\
数组越界风险高 & 全部下标检查一遍，必要时数组开大 5 到 10。 \\
输出格式不确定 & 对照样例逐字符检查空格、换行、小数位。 \\
\bottomrule
\end{longtable}


\part{基础与题型}
\chapter{C++ 语法与 STL 速查}\label{ch:04}

\section{本章使用方式}\label{sec:04-1}

这一章放常用的 C++ 写法。忘了容器怎么声明、循环怎么写，或者拿不准整数除法和取模的结果时，可以从这里查：

\begin{itemize}
  \item 一个程序应该怎么开头。
  \item 输入输出怎么写。
  \item 什么时候必须用 \Code{long long}。
  \item 整数除法、取模、向上取整、向下取整怎么写。
  \item 数组、\Code{vector}、\Code{string}、结构体怎么用。
  \item \Code{sort}、\Code{map}、\Code{set}、\Code{queue}、\Code{priority\_queue} 怎么写。
  \item lambda 比较器怎么写，常见错误是什么。
\end{itemize}

写法熟悉就继续用，不必为了把代码写短而临时换一种。

\section{语言标准与代码拼装}\label{sec:04-2}
本书代码使用 \textbf{GNU C++17}。结构化绑定 \Code{auto [x,y]} 和 \Code{std::gcd} 需要 C++17，编译时可使用 \Code{-std=c++17}。

若只支持旧标准，把 \Code{auto [x,y]=p;} 改成 \Code{auto x=p.first; auto y=p.second;}；\Code{gcd} 可自行写欧几里得算法。\Code{bits/stdc++.h} 是 GNU 扩展；非 GNU 环境改用对应标准头文件。

\textbf{模板不是一整份提交代码。}先写下面的主函数骨架；结构体、普通函数放在 \Code{main} 外，输入、循环、输出放在 \Code{main} 或 \Code{solve} 内。不要把两个同名 \Code{main}、\Code{Node}、\Code{Edge} 或 \Code{INF} 一起粘贴。带 \Code{check}、\Code{handle...} 和省略号的代码是待按题意补全的框架。

\section{最小程序骨架}\label{sec:04-3}

大部分 CSP 题目都可以从这个骨架开始：

\begin{lstlisting}
#include <bits/stdc++.h>
using namespace std;

using ll = long long;

int main() {
    ios::sync_with_stdio(false);
    cin.tie(nullptr);

    // write code here

    return 0;
}
\end{lstlisting}

说明：

\begin{itemize}
  \item \Code{\#include <bits/stdc++.h>} 是 GNU C++ 的扩展头文件，包含大多数常用标准库。
  \item \Code{using namespace std;} 可以少写 \Code{std::}。
  \item \Code{using ll = long long;} 以后可以用 \Code{ll} 表示 \Code{long long}。
  \item \Code{ios::sync\_with\_stdio(false); cin.tie(nullptr);} 用来加速 C++ 输入输出。
  \item 使用上面两行加速后，不要混用 \Code{scanf/printf} 和 \Code{cin/cout}。
\end{itemize}

\section{输入输出}\label{sec:04-4}

\subsection{普通输入}

\begin{lstlisting}
int n;
cin >> n;

int a, b;
cin >> a >> b;

string s;
cin >> s; // reads until whitespace
\end{lstlisting}

\subsection{读入数组}

\begin{lstlisting}
int n;
cin >> n;
vector<int> a(n);
for (int i = 0; i < n; i++) {
    cin >> a[i];
}
\end{lstlisting}

如果题目下标从 1 开始，考场上可以直接开成 \Code{n + 1}：

\begin{lstlisting}
int n;
cin >> n;
vector<int> a(n + 1);
for (int i = 1; i <= n; i++) {
    cin >> a[i];
}
\end{lstlisting}

\subsection{多组数据}

题目明确给出 $T$：

\begin{lstlisting}
int T;
cin >> T;
while (T--) {
    solve();
}
\end{lstlisting}

题目没有给 $T$，读到文件结束：

\begin{lstlisting}
int a, b;
while (cin >> a >> b) {
    cout << a + b << '\n';
}
\end{lstlisting}

\Warn{多组数据常见坑：}每组开始前要清空数组、图、队列、\Code{map}、答案变量。

\subsection{整行输入}

\Code{cin >> s} 只能读到空格前。如果要读整行，用 \Code{getline}：

\begin{lstlisting}
string line;
getline(cin, line);
\end{lstlisting}

如果前面刚用了 \Code{cin >> n}，再用 \Code{getline}，要先吃掉换行：

\begin{lstlisting}
int n;
cin >> n;
cin.ignore(numeric_limits<streamsize>::max(), '\n');

string line;
getline(cin, line);
\end{lstlisting}

如果担心行尾有多个字符，可以写：

\begin{lstlisting}
cin.ignore(numeric_limits<streamsize>::max(), '\n');
\end{lstlisting}

\subsection{输出}

\begin{lstlisting}
cout << ans << '\n';
cout << a << ' ' << b << '\n';
\end{lstlisting}

\Warn{不要输出提示语。}例如不要写 \Code{cout << "answer=" << ans;}，除非题目明确要求。

\subsection{小数输出}

\begin{lstlisting}
double x = 3.1415926;
cout << fixed << setprecision(2) << x << '\n'; // 3.14
\end{lstlisting}

说明：

\begin{itemize}
  \item \Code{fixed} 表示按小数位输出。
  \item \Code{setprecision(2)} 表示保留 2 位小数。
  \item 需要包含 \Code{<iomanip>}，但 \Code{bits/stdc++.h} 已包含。
\end{itemize}

\section{数据类型}\label{sec:04-5}

\subsection{整数类型}

\begin{longtable}{p{3cm}p{4cm}p{7cm}}
\toprule
类型 & 大约范围 & 使用场景 \\
\midrule
\endfirsthead
\toprule
类型 & 大约范围 & 使用场景 \\
\midrule
\endhead
\Code{int} & $-2.1 \times 10^9$ 到 $2.1 \times 10^9$ & 下标、普通计数、小范围整数。 \\
\Code{long long} & $-9.2 \times 10^{18}$ 到 $9.2 \times 10^{18}$ & 求和、乘法、答案、距离、方案数。 \\
\Code{unsigned long long} & 0 到约 $1.8 \times 10^{19}$ & 很少必须用，新手慎用。 \\
\bottomrule
\end{longtable}

必须用 \Code{long long} 的常见情况：

\begin{itemize}
  \item $n$ 个数求和，每个数可达 $10^9$。
  \item 两个最大可达 $10^9$ 的数相乘。
  \item 最短路边权和可能很大。
  \item 方案数、组合数、累计时间、累计代价。
\end{itemize}

乘法防溢出：

\begin{lstlisting}
int a = 1000000000;
int b = 1000000000;
long long x = 1LL * a * b;
\end{lstlisting}

\Warn{错误写法：}

\begin{lstlisting}
long long x = a * b; // a*b already overflowed as int
\end{lstlisting}

\subsection{浮点类型}

\begin{itemize}
  \item \Code{double} 是常用浮点数。
  \item 判断浮点相等时不要直接用 \Code{==}，通常用误差。
\end{itemize}

\begin{lstlisting}
const double EPS = 1e-9;
if (abs(a - b) < EPS) {
    // a and b are almost equal
}
\end{lstlisting}

\subsection{字符和布尔}

\begin{lstlisting}
char c = 'A';
bool ok = true;
\end{lstlisting}

字符判断：

\begin{lstlisting}
isdigit((unsigned char)c); // is digit
isalpha((unsigned char)c); // is letter
islower((unsigned char)c);
isupper((unsigned char)c);
tolower((unsigned char)c);
toupper((unsigned char)c);
\end{lstlisting}

\Warn{注意：}\Code{isdigit((unsigned char)c)} 返回非零表示真，不一定返回 1。

\section{表达式、整除和取整}\label{sec:04-6}

\subsection{基本运算符}

C++ 中常用算术运算符如下：

\begin{longtable}{p{3cm}p{4cm}p{7cm}}
\toprule
运算符 & 含义 & 注意点 \\
\midrule
\endfirsthead
\toprule
运算符 & 含义 & 注意点 \\
\midrule
\endhead
\Code{+ - *} & 加、减、乘 & 两个 \Code{int} 相乘仍先按 \Code{int} 计算。 \\
\Code{/} & 除法 & 两边都是整数时做整数除法，结果向 0 截断。 \\
\Code{\%} & 取余 & 只能用于整数；负数取余的结果可能为负。 \\
\Code{+= -= *= /= \%=} & 复合赋值 & \Code{x += y} 等价于 \Code{x = x + y}。 \\
\Code{++ --} & 自增、自减 & 简单循环中优先写 \Code{i++} 或 \Code{++i}，不要嵌进复杂表达式。 \\
\bottomrule
\end{longtable}

\Warn{整数除法不是数学里的向下取整。}在 C++ 中，\Code{-3 / 2} 的结果是 \Code{-1}，因为整数除法向 0 截断；数学意义上的 $\lfloor -3/2 \rfloor$ 是 $-2$。

\subsection{整数除法和类型转换}

如果两边都是整数，\Code{/} 会丢掉小数部分：

\begin{lstlisting}
int a = 5, b = 2;
cout << a / b << '\n'; // 2
\end{lstlisting}

如果要小数结果，至少把一个操作数转成浮点数：

\begin{lstlisting}
cout << 1.0 * a / b << '\n';       // 2.5
cout << (double)a / b << '\n';     // 2.5
\end{lstlisting}

整数答案尽量用整数运算求。大整数转成 \Code{double} 后可能丢失低位，取整时尤其要留意。

\subsection{正整数向上取整}

最常见场景：$a$ 个物品，每组最多放 $b$ 个，问需要几组。若 $a \ge 0, b > 0$：

\begin{lstlisting}
long long ceil_div_positive(long long a, long long b) {
    return a / b + (a % b != 0);
}
\end{lstlisting}

常见特例：

\begin{lstlisting}
long long half_up = (n + 1) / 2;          // ceil(n / 2), n >= 0
long long blocks = (n + block - 1) / block;
\end{lstlisting}

\Warn{这个简式只适合 $a \ge 0, b > 0$。}如果 $a$ 可能为负，或者 $b$ 可能为负，不能直接套 \Code{a/b+(a\%b!=0)}。

\subsection{有符号整数向下和向上取整}

如果题目涉及负坐标、负偏移、时间差、数学分段，建议使用下面的通用函数。它们要求 \Code{b != 0}，并且遵循数学意义上的 floor 和 ceil。

\begin{lstlisting}
// Require b != 0; exclude LLONG_MIN / -1.
long long floor_div(long long a, long long b) {
    long long q = a / b;
    long long r = a % b;
    if (r != 0 && ((r > 0) != (b > 0))) q--;
    return q;
}

long long ceil_div(long long a, long long b) {
    long long q = a / b;
    long long r = a % b;
    if (r != 0 && ((r > 0) == (b > 0))) q++;
    return q;
}
\end{lstlisting}

快速自测：

\begin{longtable}{p{4cm}p{4cm}p{4cm}}
\toprule
表达式 & \Code{floor\_div} & \Code{ceil\_div} \\
\midrule
\endfirsthead
\toprule
表达式 & \Code{floor\_div} & \Code{ceil\_div} \\
\midrule
\endhead
$3 / 2$ & 1 & 2 \\
$-3 / 2$ & -2 & -1 \\
$3 / -2$ & -2 & -1 \\
$-3 / -2$ & 1 & 2 \\
\bottomrule
\end{longtable}

\subsection{负数取模}

C++ 中 \Code{x \% m} 的符号跟 \Code{x} 相同。例如 \Code{-1 \% 5} 是 \Code{-1}，不是 \Code{4}。如果题目需要 $0$ 到 $m-1$ 的标准余数，写：

\begin{lstlisting}
long long norm_mod(long long x, long long m) {
    long long r = x % m;
    if (r < 0) r += m;
    return r;
}
\end{lstlisting}

适用：\Code{m > 0}。若 \Code{x} 可能小于 \Code{-m} 很多，上面仍然正确，因为 C++ 的一次取余后 \Code{r} 已经在 \Code{(-m,m)} 内。

\subsection{运算优先级和括号}

常用的优先级关系如下。混合运算拿不准时，可以多加一层括号：

\begin{itemize}
  \item 乘、除、取模优先于加减。
  \item 比较运算优先级低于加减乘除。
  \item \Code{\&\&} 优先于 \Code{||}。
  \item 位运算和比较混在一起时必须加括号。
\end{itemize}

\begin{lstlisting}
if ((mask & (1 << k)) != 0) {
    // bit k is 1
}
\end{lstlisting}

\Warn{有疑问就加括号。}括号不会让程序变慢，却能显著减少读错表达式的风险。

\subsection{常量和初始化}

常用写法：

\begin{lstlisting}
const int N = 100000 + 5;
const long long INF = (long long)4e18;
const int MOD = 1000000007;

int cnt = 0;
long long sum = 0;
bool ok = true;
\end{lstlisting}

提醒：

\begin{itemize}
  \item 局部普通变量不会自动清零，定义后要赋初值。
  \item 全局数组会默认清零，但多组数据不会自动恢复初始状态。
  \item \Code{const} 常量不能被修改，适合写数组上限、模数、无穷大。
\end{itemize}


\subsection{四舍五入与银行家舍入不能混用}
\Code{llround} 的半整数规则是远离 0；题面要求“恰好一半时取偶数”时应自行实现。若输入为非负数、恰有一位小数，且整数部分不超过 $10^5$，可以完全避免浮点：
\begin{lstlisting}
pair<long long,long long> roundOneDecimal(const string& s) {
    size_t dot = s.find('.');
    long long a = stoll(s.substr(0, dot));
    int b = s[dot + 1] - '0';
    long long halfUp = a + (b >= 5);
    long long halfEven = a + (b > 5 || (b == 5 && a % 2 == 1));
    return {halfUp, halfEven};
}
\end{lstlisting}
例如 2.5 分别得到 3 与 2，3.5 两者均为 4。此代码仅针对上述输入格式；不能直接用于负数、多位小数或超大整数。\Code{setprecision} 负责显示格式，不替代题目定义的舍入算法。

\section{控制流}\label{sec:04-7}

\subsection{if}

\begin{lstlisting}
if (x > 0) {
    cout << "positive\n";
} else if (x == 0) {
    cout << "zero\n";
} else {
    cout << "negative\n";
}
\end{lstlisting}

\subsection{for}

\begin{lstlisting}
for (int i = 0; i < n; i++) {
    // 0, 1, ..., n-1
}

for (int i = 1; i <= n; i++) {
    // 1, 2, ..., n
}
\end{lstlisting}

倒序：

\begin{lstlisting}
for (int i = n - 1; i >= 0; i--) {
    // be careful when i is unsigned
}
\end{lstlisting}

\subsection{while}

\begin{lstlisting}
while (condition) {
    // loop
}
\end{lstlisting}

读到结束：

\begin{lstlisting}
int x;
while (cin >> x) {
    // process x
}
\end{lstlisting}

\subsection{break 和 continue}

\begin{lstlisting}
for (int i = 0; i < n; i++) {
    if (a[i] < 0) continue; // skip this round
    if (a[i] == target) break; // exit loop
}
\end{lstlisting}

\section{auto、引用和范围 for}\label{sec:04-8}

\subsection{auto}

\Code{auto} 让编译器根据右侧表达式推断类型，常用于迭代器、pair、复杂容器遍历。

\begin{lstlisting}
auto it = lower_bound(a.begin(), a.end(), x);
auto p = make_pair(3, 5);
\end{lstlisting}

遍历 \Code{map}：

\begin{lstlisting}
map<string, int> cnt;
for (auto [key, value] : cnt) {
    cout << key << ' ' << value << '\n';
}
\end{lstlisting}

\Warn{注意复制。}\Code{auto [key, value]} 会复制键和值；如果元素很大，写 \Code{const auto\&} 更合适。

\begin{lstlisting}
for (const auto& item : records) {
    // read item without copying
}
\end{lstlisting}

\subsection{引用}

引用相当于变量的别名。函数传大对象时常用引用避免复制。

\begin{lstlisting}
void printVector(const vector<int>& a) {
    for (int x : a) cout << x << ' ';
}

void addOne(vector<int>& a) {
    for (int& x : a) x++;
}
\end{lstlisting}

区别：

\begin{itemize}
  \item \Code{const vector<int>\& a}：不复制，也不允许函数内修改。
  \item \Code{vector<int>\& a}：不复制，允许函数内修改原数组。
  \item \Code{vector<int> a}：会复制整个数组，数据大时慢。
\end{itemize}

\subsection{范围 for}

只读遍历：

\begin{lstlisting}
for (int x : a) {
    cout << x << '\n';
}
\end{lstlisting}

修改原元素：

\begin{lstlisting}
for (int& x : a) {
    x++;
}
\end{lstlisting}

\Warn{如果忘了写引用，修改不会回到原数组。}

\begin{lstlisting}
for (int x : a) {
    x++; // only modifies the copy
}
\end{lstlisting}

\subsection{size 的类型}

\Code{a.size()} 的返回类型是无符号整数。考场上为了减少类型麻烦，常写成：

\begin{lstlisting}
int n = (int)a.size();
for (int i = 0; i < n; i++) {
    // use a[i]
}
\end{lstlisting}

倒序遍历不要写下面这种：

\begin{lstlisting}
for (int i = a.size() - 1; i >= 0; i--) {
    // dangerous when a is empty
}
\end{lstlisting}

更稳的写法：

\begin{lstlisting}
for (int i = (int)a.size() - 1; i >= 0; i--) {
    // ok even when a is empty: i starts at -1
}
\end{lstlisting}

\section{函数}\label{sec:04-9}

把重复逻辑写成函数，能减少大模拟和算法题的错误。

\begin{lstlisting}
int add(int a, int b) {
    return a + b;
}

bool isEven(int x) {
    return x % 2 == 0;
}
\end{lstlisting}

传 vector：

\begin{lstlisting}
long long sumVector(const vector<int>& a) {
    long long sum = 0;
    for (int x : a) sum += x;
    return sum;
}
\end{lstlisting}

说明：

\begin{itemize}
  \item \Code{const vector<int>\& a} 表示不复制整个数组，并且函数内部不修改它。
  \item 如果函数需要修改数组，去掉 \Code{const}。
\end{itemize}

\begin{lstlisting}
void addOne(vector<int>& a) {
    for (int& x : a) x++;
}
\end{lstlisting}

\section{数组}\label{sec:04-10}

\subsection{普通数组}

全局数组适合大数组：

\begin{lstlisting}
const int N = 100000 + 5;
int a[N];
long long pre[N];
\end{lstlisting}

说明：

\begin{itemize}
  \item 全局数组默认初始化为 0。
  \item 局部大数组可能爆栈，尽量放全局或用 \Code{vector}。
\end{itemize}

\subsection{二维数组}

\begin{lstlisting}
const int N = 1005;
int grid[N][N];
\end{lstlisting}

如果规模不确定：

\begin{lstlisting}
int n, m;
cin >> n >> m;
vector<vector<int>> grid(n, vector<int>(m, 0));
\end{lstlisting}

\Warn{内存估算：}\Code{int a[10000][10000]} 大约 400MB，通常危险。

\subsection{初始化}

全局普通数组默认是 0。如果要重置数组：

\begin{lstlisting}
int a[1005];
memset(a, 0, sizeof a);
memset(a, -1, sizeof a);
\end{lstlisting}

\Warn{注意：}\Code{memset} 适合把数组设为 0 或 -1，不适合把 \Code{int} 数组设成 1、2、无穷大。

设置成大数：

\begin{lstlisting}
const int INF = 0x3f3f3f3f;
int dist[1005];
memset(dist, 0x3f, sizeof dist);
\end{lstlisting}

\Code{vector} 初始化：

\begin{lstlisting}
vector<int> a(n, 0);
vector<long long> dist(n + 1, (long long)4e18);
vector<vector<int>> grid(n, vector<int>(m, 0));
\end{lstlisting}

\section{vector}\label{sec:04-11}

\subsection{创建}

\begin{lstlisting}
vector<int> a;          // empty
vector<int> b(10);      // 10 zeros
vector<int> c(10, -1);  // 10 values, all -1
\end{lstlisting}

\subsection{常用操作}

\begin{lstlisting}
a.push_back(5);
a.pop_back();
a.clear();
int n = (int)a.size();
bool empty = a.empty();
\end{lstlisting}

访问：

\begin{lstlisting}
cout << a[0] << '\n';
cout << a.back() << '\n';
\end{lstlisting}

\Warn{访问前检查非空。}\Code{a.back()}、\Code{a[0]} 在空数组上会出错。

\subsection{遍历}

\begin{lstlisting}
for (int i = 0; i < (int)a.size(); i++) {
    cout << a[i] << '\n';
}

for (int x : a) {
    cout << x << '\n';
}

for (int& x : a) {
    x++; // modify original element
}
\end{lstlisting}

\subsection{删除元素}

删除最后一个：

\begin{lstlisting}
a.pop_back();
\end{lstlisting}

删除某个位置：

\begin{lstlisting}
a.erase(a.begin() + pos);
\end{lstlisting}

\Warn{注意：}\Code{erase} 是 $O(n)$，大量删除时可能超时。

删除所有等于 x 的元素：

\begin{lstlisting}
a.erase(remove(a.begin(), a.end(), x), a.end());
\end{lstlisting}

\subsection{二维 vector}

二维网格：

\begin{lstlisting}
int n, m;
cin >> n >> m;
vector<vector<int>> a(n, vector<int>(m));
for (int i = 0; i < n; i++) {
    for (int j = 0; j < m; j++) {
        cin >> a[i][j];
    }
}
\end{lstlisting}

二维字符图：

\begin{lstlisting}
int n, m;
cin >> n >> m;
vector<string> g(n);
for (int i = 0; i < n; i++) {
    cin >> g[i];
}
cout << g[0][0] << '\n';
\end{lstlisting}

图的邻接表：

\begin{lstlisting}
int n, m;
cin >> n >> m;
vector<vector<int>> g(n + 1);
for (int i = 0; i < m; i++) {
    int u, v;
    cin >> u >> v;
    g[u].push_back(v);
    g[v].push_back(u);
}
\end{lstlisting}

\section{string}\label{sec:04-12}

\subsection{基本操作}

\begin{lstlisting}
string s = "abc";
int n = (int)s.size();
char c = s[0];
s.push_back('d');
s += "ef";
\end{lstlisting}

\subsection{子串}

\begin{lstlisting}
string s = "abcdef";
string t = s.substr(2, 3); // "cde"
\end{lstlisting}

\Code{substr(pos, len)} 从 \Code{pos} 开始取 \Code{len} 个字符，下标从 0 开始。

\subsection{查找}

\begin{lstlisting}
string s = "abcdef";
if (s.find("cd") != string::npos) {
    cout << "found\n";
}
\end{lstlisting}

\Warn{不要写：}

\begin{lstlisting}
if (s.find("cd")) {
    // wrong when position is 0
}
\end{lstlisting}

\subsection{字符串转数字}

\begin{lstlisting}
int x = stoi("123");
long long y = stoll("1234567890123");
double z = stod("3.14");
\end{lstlisting}

数字转字符串：

\begin{lstlisting}
string s = to_string(123);
\end{lstlisting}

\Warn{如果数字很长，}\Code{stoi/stoll} 会溢出或异常，要用字符串模拟。

\subsection{按空格分割}

\begin{lstlisting}
string line;
getline(cin, line);
stringstream ss(line);

string word;
while (ss >> word) {
    cout << word << '\n';
}
\end{lstlisting}

\section{pair 和 tuple}\label{sec:04-13}

\subsection{pair}

\begin{lstlisting}
pair<int, int> p = {3, 5};
cout << p.first << ' ' << p.second << '\n';
\end{lstlisting}

常用于坐标、边、距离和编号：

\begin{lstlisting}
vector<pair<int, int>> points;
points.push_back({x, y});
\end{lstlisting}

默认排序：

\begin{lstlisting}
sort(points.begin(), points.end());
\end{lstlisting}

默认先按 \Code{first} 升序，再按 \Code{second} 升序。

\subsection{tuple}

字段超过两个时可以用 \Code{tuple}，但新手更推荐结构体。

\begin{lstlisting}
tuple<int, int, int> t = {1, 2, 3};
auto [a, b, c] = t;
\end{lstlisting}

\section{结构体 struct}\label{sec:04-14}

结构体适合表示学生、任务、边、事件、状态。

\begin{lstlisting}
struct Student {
    int id;
    string name;
    int score;
};

Student s;
s.id = 1;
s.name = "Alice";
s.score = 90;
\end{lstlisting}

数组：

\begin{lstlisting}
vector<Student> students;
students.push_back({1, "Alice", 90});
\end{lstlisting}

图的边：

\begin{lstlisting}
struct Edge {
    int to;
    int w;
};

vector<vector<Edge>> g(n + 1);
g[u].push_back({v, w});
\end{lstlisting}

\section{排序 sort}\label{sec:04-15}

\subsection{普通排序}

\begin{lstlisting}
sort(a.begin(), a.end()); // ascending
sort(a.rbegin(), a.rend()); // descending
\end{lstlisting}

数组排序：

\begin{lstlisting}
sort(a, a + n);
\end{lstlisting}

\subsection{自定义排序：从大到小}

\begin{lstlisting}
sort(a.begin(), a.end(), [](int x, int y) {
    return x > y;
});
\end{lstlisting}

\subsection{结构体排序}

成绩高的在前；成绩相同，编号小的在前：

\begin{lstlisting}
struct Student {
    int id;
    int score;
};

vector<Student> v;

sort(v.begin(), v.end(), [](const Student& a, const Student& b) {
    if (a.score != b.score) return a.score > b.score;
    return a.id < b.id;
});
\end{lstlisting}

\subsection{按 vector 的某一列排序}

\begin{lstlisting}
vector<vector<int>> a;

sort(a.begin(), a.end(), [](const vector<int>& x, const vector<int>& y) {
    return x[0] < y[0]; // sort by column 0 ascending
});
\end{lstlisting}

先按第 0 列升序，再按第 1 列降序：

\begin{lstlisting}
sort(a.begin(), a.end(), [](const vector<int>& x, const vector<int>& y) {
    if (x[0] != y[0]) return x[0] < y[0];
    return x[1] > y[1];
});
\end{lstlisting}

\subsection{stable\_sort}

如果题目要求相同关键字保持原输入顺序：

\begin{lstlisting}
stable_sort(v.begin(), v.end(), [](const Student& a, const Student& b) {
    return a.score > b.score;
});
\end{lstlisting}

\subsection{比较器高危错误}

比较器必须表达“谁应该排在前面”。

\Good{正确：}

\begin{lstlisting}
sort(v.begin(), v.end(), [](const Student& a, const Student& b) {
    if (a.score != b.score) return a.score > b.score;
    return a.id < b.id;
});
\end{lstlisting}

\Warn{错误：}

\begin{lstlisting}
sort(v.begin(), v.end(), [](const Student& a, const Student& b) {
    return a.score >= b.score; // wrong: equality should return false
});
\end{lstlisting}

相等时返回 \Code{true} 可能导致排序行为错误。

\section{二分相关函数}\label{sec:04-16}

\subsection{binary\_search}

\begin{lstlisting}
sort(a.begin(), a.end());
bool ok = binary_search(a.begin(), a.end(), x);
\end{lstlisting}

\subsection{lower\_bound 和 upper\_bound}

\begin{lstlisting}
sort(a.begin(), a.end());

auto it1 = lower_bound(a.begin(), a.end(), x); // first >= x
auto it2 = upper_bound(a.begin(), a.end(), x); // first > x
\end{lstlisting}

下标：

\begin{lstlisting}
int pos = lower_bound(a.begin(), a.end(), x) - a.begin();
\end{lstlisting}

\Warn{解引用前检查：}

\begin{lstlisting}
auto it = lower_bound(a.begin(), a.end(), x);
if (it != a.end()) {
    cout << *it << '\n';
}
\end{lstlisting}

\section{map 和 unordered\_map}\label{sec:04-17}

\subsection{map}

\Code{map} 按 key 自动升序排列，常用于字符串计数、编号映射、需要有序遍历的场景。

\begin{lstlisting}
map<string, int> cnt;
cnt["apple"]++;
cnt["banana"] += 2;

cout << cnt["apple"] << '\n';
\end{lstlisting}

判断是否存在：

\begin{lstlisting}
if (cnt.count("apple")) {
    cout << "exists\n";
}

auto it = cnt.find("apple");
if (it != cnt.end()) {
    cout << it->second << '\n';
}
\end{lstlisting}

\Warn{注意：}\Code{cnt[key]} 会在 key 不存在时自动创建一个默认值。

遍历：

\begin{lstlisting}
for (auto [key, value] : cnt) {
    cout << key << ' ' << value << '\n';
}
\end{lstlisting}

\subsection{unordered\_map}

\Code{unordered\_map} 平均更快，但没有顺序。

\begin{lstlisting}
unordered_map<string, int> cnt;
cnt["apple"]++;
\end{lstlisting}

考场建议：

\begin{itemize}
  \item 需要按 key 排序输出，用 \Code{map}。
  \item 只做计数且数据量大，用 \Code{unordered\_map}。
  \item 自定义 key 时新手优先考虑转成字符串或 \Code{long long} 编码，不要现场写复杂 hash。
\end{itemize}

\section{unordered\_set}\label{sec:04-18}

\Code{unordered\_set} 用来快速判断一个值是否出现过，不关心顺序。

\begin{lstlisting}
unordered_set<int> seen;
seen.insert(5);

if (seen.count(5)) {
    cout << "yes\n";
}
\end{lstlisting}

字符串集合：

\begin{lstlisting}
unordered_set<string> words;
words.insert("abc");
if (words.count("abc")) {
    cout << "exists\n";
}
\end{lstlisting}

如果需要按升序遍历，用 \Code{set}，不要用 \Code{unordered\_set}。

\section{set 和 multiset}\label{sec:04-19}

\subsection{set}

\Code{set} 自动排序且不重复。

\begin{lstlisting}
set<int> s;
s.insert(3);
s.insert(1);
s.erase(3);

if (s.count(1)) {
    cout << "exists\n";
}
\end{lstlisting}

查找第一个大于等于 x 的元素：

\begin{lstlisting}
auto it = s.lower_bound(x);
if (it != s.end()) {
    cout << *it << '\n';
}
\end{lstlisting}

\subsection{multiset}

\Code{multiset} 允许重复元素。

\begin{lstlisting}
multiset<int> s;
s.insert(5);
s.insert(5);
cout << s.count(5) << '\n'; // 2
\end{lstlisting}

删除一个 5：

\begin{lstlisting}
auto it = s.find(5);
if (it != s.end()) s.erase(it);
\end{lstlisting}

删除所有 5：

\begin{lstlisting}
s.erase(5);
\end{lstlisting}

\section{queue、stack、deque}\label{sec:04-20}

\subsection{queue}

BFS 常用队列。

\begin{lstlisting}
queue<int> q;
q.push(1);
int x = q.front();
q.pop();
bool empty = q.empty();
\end{lstlisting}

\subsection{stack}

括号匹配、表达式、DFS 模拟可用栈。

\begin{lstlisting}
stack<char> st;
st.push('(');
char c = st.top();
st.pop();
\end{lstlisting}

\Warn{调用 \Code{front/top} 前必须检查非空。}

\subsection{deque}

双端队列，两边都能进出。

\begin{lstlisting}
deque<int> dq;
dq.push_back(1);
dq.push_front(2);
dq.pop_back();
dq.pop_front();
\end{lstlisting}

单调队列会用到 \Code{deque}。

\section{priority\_queue}\label{sec:04-21}

\subsection{大根堆}

默认是大根堆，最大值在顶部。

\begin{lstlisting}
priority_queue<int> pq;
pq.push(3);
pq.push(10);
cout << pq.top() << '\n'; // 10
pq.pop();
\end{lstlisting}

\subsection{小根堆}

\begin{lstlisting}
priority_queue<int, vector<int>, greater<int>> pq;
pq.push(3);
pq.push(10);
cout << pq.top() << '\n'; // 3
\end{lstlisting}

\subsection{pair 小根堆}

Dijkstra 常用：

\begin{lstlisting}
using P = pair<long long, int>; // distance, node
priority_queue<P, vector<P>, greater<P>> pq;
pq.push({0, 1});
\end{lstlisting}

pair 的默认比较先看 \Code{first}，再看 \Code{second}。

\section{常用算法函数}\label{sec:04-22}

\subsection{min、max、swap}

\begin{lstlisting}
int x = min(a, b);
int y = max(a, b);
swap(a, b);
\end{lstlisting}

\subsection{reverse}

\begin{lstlisting}
reverse(a.begin(), a.end());
reverse(s.begin(), s.end());
\end{lstlisting}

\subsection{unique 去重}

若要去掉所有重复值且允许改变原顺序，先排序。\Code{unique} 只消除相邻重复；要保留首次出现顺序，用集合配合一次扫描：

\begin{lstlisting}
sort(a.begin(), a.end());
a.erase(unique(a.begin(), a.end()), a.end());
\end{lstlisting}

\subsection{accumulate}

\begin{lstlisting}
long long sum = accumulate(a.begin(), a.end(), 0LL);
\end{lstlisting}

\Warn{第三个参数写 \Code{0LL}，否则可能按 \Code{int} 求和。}

\subsection{gcd}

\begin{lstlisting}
// Require a,b >= 0 and lcm(a,b) <= LLONG_MAX.
long long g = gcd(a, b);
long long l = (a == 0 || b == 0) ? 0 : a / g * b;
\end{lstlisting}

\Warn{最小公倍数先除后乘。}这样比 \Code{a * b / gcd(a,b)} 更不容易溢出。

\subsection{常用数学函数}

\begin{lstlisting}
double r = sqrt(x);      // square root
double y = pow(a, b);    // a^b, floating point
double f = floor(v);     // round down as floating point
double c = ceil(v);      // round up as floating point
long long z = llround(v); // nearest integer
\end{lstlisting}

提醒：

\begin{itemize}
  \item \Code{sqrt}、\Code{pow} 返回浮点数，不适合直接处理大整数精确幂。
  \item 整数平方尽量写 \Code{1LL * x * x}，不要写 \Code{pow(x, 2)}。
  \item 判断一个整数是否为完全平方数时，先取 \Code{sqrt}，再用整数乘法校验。
  \item 整数向上取整不要用 \Code{ceil(1.0*a/b)} 代替整数公式。
\end{itemize}

\begin{lstlisting}
bool isSquare(long long x) {
    if (x < 0) return false;
    long long r = sqrt((long double)x);
    while (r > 0 && r > x / r) r--;
    while ((r + 1) <= x / (r + 1)) r++;
    return r * r == x;
}
\end{lstlisting}

\section{lambda 入门}\label{sec:04-23}

lambda 是临时函数，排序时非常常用。

最基本写法：

\begin{lstlisting}
[](int a, int b) {
    return a < b;
}
\end{lstlisting}

用于排序：

\begin{lstlisting}
sort(a.begin(), a.end(), [](int x, int y) {
    return x > y;
});
\end{lstlisting}

使用外部变量：

\begin{lstlisting}
int mod = 10;
sort(a.begin(), a.end(), [mod](int x, int y) {
    return x % mod < y % mod;
});
\end{lstlisting}

引用捕获：

\begin{lstlisting}
int cnt = 0;
auto add = [&](int x) {
    cnt += x;
};
add(5);
\end{lstlisting}

考场建议：

\begin{itemize}
  \item 排序比较器中优先使用 \Code{const Type\&}，避免复制。
  \item 不熟悉捕获时，尽量不要写复杂 lambda。
  \item 复杂逻辑可以改成普通函数或先计算辅助字段。
\end{itemize}

\section{常见编译错误和原因}\label{sec:04-24}

\begin{longtable}{p{4cm}p{5cm}p{5cm}}
\toprule
报错或现象 & 常见原因 & 处理方法 \\
\midrule
\endfirsthead
\toprule
报错或现象 & 常见原因 & 处理方法 \\
\midrule
\endhead
\Code{expected ';'} & 少分号 & 看上一行结构体、变量定义、语句末尾。 \\
\Code{was not declared} & 变量未定义或作用域不对 & 检查拼写和定义位置。 \\
\Code{no matching function} & 函数参数类型不匹配 & 检查传参类型、引用、const。 \\
\Code{segmentation fault} & 越界、空容器取值、递归爆栈 & 查数组下标、\Code{front/top/back} 前是否非空。 \\
样例正确但大数据错 & 溢出或复杂度过高 & 检查 \Code{long long} 和复杂度。 \\
输出格式错误 & 多输出、少换行、空格不对 & 严格对照题目输出格式。 \\
\bottomrule
\end{longtable}

\section{C 和 C++ 混用提醒}\label{sec:04-25}

如果你使用 C++，建议全程使用 \Code{cin/cout} 或全程使用 \Code{scanf/printf}，不要混用。

如果确实使用 C 风格输入输出：

\begin{lstlisting}
int x;
scanf("%d", &x);

long long y;
scanf("%lld", &y);

printf("%d\n", x);
printf("%lld\n", y);
\end{lstlisting}

字符串：

\begin{lstlisting}
char s[1005];
scanf("%s", s);
\end{lstlisting}

初学时可以先熟悉 \Code{string}、\Code{vector} 和 \Code{map}，字符串和动态数组的存储会方便一些。

\section{考场 C++ 快速检查}\label{sec:04-26}

提交前快速看一遍：

\begin{itemize}
  \item 是否包含了头文件和 \Code{using namespace std;}。
  \item 是否开了快速输入输出。
  \item 是否混用了 \Code{cin/cout} 和 \Code{scanf/printf}。
  \item 求和、乘法、答案是否用了 \Code{long long}。
  \item vector 是否按正确大小初始化。
  \item 访问 \Code{front/top/back} 前是否判断非空。
  \item \Code{lower\_bound} 返回 \Code{end} 时是否处理。
  \item 排序比较器相等时是否返回 \Code{false} 或继续比较。
  \item 多组数据是否清空容器。
  \item 输出是否完全符合题意。
\end{itemize}

\chapter{A 题决策树与高频模板}

\section{A 题定位}

A 题通常不是为了考很深的算法，而是考你能否准确读题、正确实现、处理边界、按格式输出。对新手来说，A 题最重要的不是写得快，而是写得稳。

\begin{itemize}
  \item 建议目标时间：20--35 分钟。
  \item 最大容忍时间：45 分钟。超过后仍无法通过样例，应先转下一题。
  \item 目标分数：100 分。
  \item 核心原则：先正确，再简洁；先样例，再边界。
\end{itemize}

\Warn{A 题最大敌人不是算法，而是低级错误。}常见低级错误包括输出格式错、下标错、整数溢出、没读完整输入、把题意中的“包含/不包含”看反。

\section{A 题开题流程}

每道 A 题按下面顺序处理，不要直接看完第一段就开始写。

\begin{enumerate}
  \item 看输入格式：有几个数、几行、是否有字符串、是否有多组。
  \item 看输出格式：输出一个数、一行多个数、字符串、小数、是否有特殊提示。
  \item 看数据范围：是否需要 \Code{long long}，数组要开多大。
  \item 判断题型：统计、公式、排序、字符串、日期、进制、简单模拟。
  \item 在草稿纸上写 1 到 2 个极端样例。
  \item 写代码。
  \item 过样例后，测试边界。
\end{enumerate}

\section{A 题读题标注法}

新手做 A 题时，最容易出现“感觉看懂了，但代码写偏了”。建议在草稿纸上用固定格式标注。

\subsection{读题时必须圈出的信息}

\begin{longtable}{p{3.5cm}p{5cm}p{5.5cm}}
\toprule
要圈的信息 & 为什么重要 & 例子 \\
\midrule
输入数量 & 决定循环次数 & 读 $n$ 个数、读 $n$ 行、读到 EOF。 \\
数据范围 & 决定数组大小和复杂度 & $n \le 10^5$ 时不要双重循环。 \\
数值范围 & 决定是否用 \Code{long long} & 每个数 $10^9$，总和可能溢出。 \\
下标规则 & 决定 0 下标或 1 下标 & 第 1 个、第 $i$ 个、编号从 0 开始。 \\
区间端点 & 决定是否包含端点 & $[l,r]$、$(l,r)$、从 $l$ 到 $r$。 \\
输出格式 & 决定空格、换行、小数位 & 输出 YES，保留 2 位小数。 \\
特殊情况 & 决定是否特判 & 无解、相等、为空、越界、最后一天。 \\
\bottomrule
\end{longtable}

\subsection{草稿纸模板}

做题前写 6 行：

\begin{lstlisting}
Input:
Output:
n range:
value range:
Need long long? yes/no
Type: statistic / formula / sort / string / date / base / simulate
\end{lstlisting}

示例：

\begin{lstlisting}
Input: n, then n integers
Output: sum and max
n range: 1e5
value range: 1e9
Need long long? yes
Type: statistic
\end{lstlisting}

\subsection{样例手算方法}

样例不是只拿来跑程序的。写代码前先手算一遍：

\begin{enumerate}
  \item 把输入拆成变量。
  \item 按题目规则一步一步算。
  \item 写出中间结果。
  \item 对比样例输出。
\end{enumerate}

如果你手算不出样例，就不要急着写代码。A 题手算不清，代码大概率也会写偏。

\section{A 题总决策树}

\begin{longtable}{p{4cm}p{5cm}p{5cm}}
\toprule
题面信号 & 优先判断为 & 去看 \\
\midrule
求和、平均、最大、最小、计数 & 简单统计 & 本章“统计模板” \\
根据公式算一个值 & 公式题 & 本章“公式题检查” \\
要求排序、排名、第几名 & 排序题 & 第二章 sort；本章“排序排名” \\
出现字符、字符串、单词、大小写 & 字符串题 & 第二章 string；本章“字符串处理” \\
出现二进制、十六进制、进制转换 & 进制题 & 本章“进制处理” \\
出现日期、时间、星期、天数 & 日期时间题 & 本章“日期时间” \\
出现坐标、网格、方向移动 & 简单模拟或网格模拟 & 本章“简单模拟” \\
出现百分比、保留小数、四舍五入 & 小数输出题 & 第二章小数输出 \\
出现区间但数据很小 & 暴力枚举或前缀和 & 本章“区间入门” \\
出现重复统计、出现次数 & 计数数组或 map & 第二章 map；本章“频率统计” \\
\bottomrule
\end{longtable}

\section{A 题关键词到动作}

这张表用于读题时快速决定下一步。每一行都应该能从题面直接判断。

\begin{longtable}{p{3.5cm}p{4.3cm}p{6.2cm}}
\toprule
你看到 & 立刻做什么 & 小心什么 \\
\midrule
“共 $n$ 个” & 开数组或 \Code{vector}，循环 $n$ 次 & 循环是 \Code{i < n} 还是 \Code{i <= n}。 \\
“编号从 1 到 n” & 优先开 \Code{n+1}，按 1 下标写 & 不要访问 \Code{a[0]}。 \\
“第 k 个” & 注意 $k$ 是 1 下标 & \Code{a[k-1]} 或排序后 \Code{a[k-1]}。 \\
“最大/最小” & 一遍扫描或排序 & 全负数时最大值不能初始化 0。 \\
“总和/累计” & 答案用 \Code{long long} & \Code{int} 乘法先溢出。 \\
“平均” & 判断输出整数还是小数 & 整数除法会截断。 \\
“每 k 个一组”“分成若干页” & 正整数向上取整 & 只在 $a \ge 0,b>0$ 时用 \Code{(a+b-1)/b}。 \\
“保留 x 位” & 用 \Code{fixed << setprecision(x)} & 是否四舍五入由输出库自动处理。 \\
“升序/降序” & 使用 \Code{sort} & 降序可用 \Code{greater<int>()}。 \\
“相同则按...” & 写完整比较器 & 相等时不能返回 \Code{true}。 \\
“字符串可能含空格” & 使用 \Code{getline} & 前面用过 \Code{cin} 要 \Code{ignore}。 \\
“统计出现次数” & 值域小用数组，值域大用 \Code{map} & 是否要求有序输出。 \\
“判断是否存在” & 用 \Code{set/count} 或排序后二分 & 找不到时不要解引用迭代器。 \\
“日期” & 写闰年、每月天数 & 2 月和跨年。 \\
“进制” & 先判断是否超过 \Code{long long} & 大数不能 \Code{stoll}。 \\
“区间 [l,r]” & 看有多少次查询 & 多次查询考虑前缀和。 \\
\bottomrule
\end{longtable}

\section{A 题数据范围判断}

\begin{longtable}{p{4cm}p{5cm}p{5cm}}
\toprule
看到的数据范围 & 含义 & 建议 \\
\midrule
$n \le 100$ & 很小 & 直接模拟或双重循环通常可行。 \\
$n \le 1000$ & 小到中等 & $O(n^2)$ 通常可试。 \\
$n \le 10^5$ & 常见大数组 & 需要 $O(n)$ 或 $O(n\log n)$，不要双重循环。 \\
数值 $\le 10^9$ & 单个数可放 int & 但加法和乘法可能要 \Code{long long}。 \\
乘积、总和、答案可能很大 & 有溢出风险 & 默认用 \Code{long long}。 \\
字符串长度很大 & 不能频繁插入删除 & 用线性扫描。 \\
\bottomrule
\end{longtable}

\section{A 题常见输入形态}

\subsection{输入一个数}

\begin{lstlisting}
int x;
cin >> x;
cout << x << '\n';
\end{lstlisting}

\subsection{输入两个或多个数}

\begin{lstlisting}
int a, b, c;
cin >> a >> b >> c;
cout << a + b + c << '\n';
\end{lstlisting}

\subsection{输入 n 个数}

\begin{lstlisting}
int n;
cin >> n;
vector<int> a(n);
for (int i = 0; i < n; i++) {
    cin >> a[i];
}
\end{lstlisting}

\subsection{输入 n 行记录}

例如每行一个学生的编号和成绩：

\begin{lstlisting}
struct Student {
    int id;
    int score;
};

int n;
cin >> n;
vector<Student> a(n);
for (int i = 0; i < n; i++) {
    cin >> a[i].id >> a[i].score;
}
\end{lstlisting}

\subsection{输入 n 行字符串}

\begin{lstlisting}
int n;
cin >> n;
vector<string> s(n);
for (int i = 0; i < n; i++) {
    cin >> s[i];
}
\end{lstlisting}

\subsection{读到文件结束}

\begin{lstlisting}
int x;
while (cin >> x) {
    // process x
}
\end{lstlisting}

\Warn{题目没有说有 T 组数据时，不要自己加 T。}

\section{简单统计}

\subsection{识别信号}

题面出现：

\begin{itemize}
  \item 求总和、平均数。
  \item 统计满足条件的个数。
  \item 找最大值、最小值。
  \item 统计正数、负数、奇数、偶数。
  \item 统计超过某个阈值的元素。
\end{itemize}

\subsection{一遍扫描模板}

\begin{lstlisting}
int n;
cin >> n;

long long sum = 0;
int cnt = 0;
int mx = -1000000000;
int mn = 1000000000;

for (int i = 0; i < n; i++) {
    int x;
    cin >> x;
    sum += x;
    if (x > mx) mx = x;
    if (x < mn) mn = x;
    if (x % 2 == 0) cnt++;
}

cout << sum << ' ' << mx << ' ' << mn << ' ' << cnt << '\n';
\end{lstlisting}

\subsection{注意事项}

\begin{itemize}
  \item 求和变量用 \Code{long long}。
  \item 最大值、最小值初始值要足够小或足够大。
  \item 如果数据可能全是负数，最大值不能初始化为 0。
  \item 平均数如果要小数，使用 \Code{double} 或输出时转换。
\end{itemize}

\subsection{统计满足条件的个数}

\begin{lstlisting}
int n;
cin >> n;
int cnt = 0;
for (int i = 0; i < n; i++) {
    int x;
    cin >> x;
    if (x >= 60) {
        cnt++;
    }
}
cout << cnt << '\n';
\end{lstlisting}

\subsection{同时统计多个类别}

例如统计正数、负数、零：

\begin{lstlisting}
int n;
cin >> n;
int pos = 0, neg = 0, zero = 0;
for (int i = 0; i < n; i++) {
    int x;
    cin >> x;
    if (x > 0) pos++;
    else if (x < 0) neg++;
    else zero++;
}
cout << pos << ' ' << neg << ' ' << zero << '\n';
\end{lstlisting}

\subsection{找最大值的位置}

\begin{lstlisting}
int n;
cin >> n;
vector<int> a(n);
for (int i = 0; i < n; i++) cin >> a[i];

int bestPos = 0;
for (int i = 1; i < n; i++) {
    if (a[i] > a[bestPos]) {
        bestPos = i;
    }
}
cout << bestPos << ' ' << a[bestPos] << '\n';
\end{lstlisting}

如果题目说“第几个”，通常要输出 \Code{bestPos + 1}。

\subsection{计数数组}

当值域很小，例如分数在 $0$ 到 $100$：

\begin{lstlisting}
int n;
cin >> n;
vector<int> cnt(101, 0);
for (int i = 0; i < n; i++) {
    int score;
    cin >> score;
    cnt[score]++;
}

for (int s = 0; s <= 100; s++) {
    if (cnt[s] > 0) {
        cout << s << ' ' << cnt[s] << '\n';
    }
}
\end{lstlisting}

\Warn{计数数组前提：}数值不能是很大的 $10^9$，也不能是负数，除非你做了偏移。

\section{存在、全部、任意判断}

A 题常见问法：

\begin{itemize}
  \item 是否存在一个数满足条件。
  \item 是否所有数都满足条件。
  \item 是否至少有一个字符串包含某内容。
  \item 是否没有任何违规情况。
\end{itemize}

\subsection{判断是否存在}

\begin{lstlisting}
int n;
cin >> n;
bool found = false;
for (int i = 0; i < n; i++) {
    int x;
    cin >> x;
    if (x == 0) {
        found = true;
    }
}

cout << (found ? "YES" : "NO") << '\n';
\end{lstlisting}

\subsection{判断是否全部满足}

\begin{lstlisting}
int n;
cin >> n;
bool ok = true;
for (int i = 0; i < n; i++) {
    int x;
    cin >> x;
    if (x < 60) {
        ok = false;
    }
}

cout << (ok ? "YES" : "NO") << '\n';
\end{lstlisting}

\subsection{提前结束}

如果后面的输入还必须读完，不要随便 \Code{break}。除非你确定后续数据不需要再读，或者已经读完本组数据。

\begin{lstlisting}
bool found = false;
for (int i = 0; i < n; i++) {
    int x;
    cin >> x;
    if (x == target) {
        found = true;
    }
}
\end{lstlisting}

\Warn{新手建议：}A 题中先不要为了省时间提前 \Code{break}，完整读入更稳。

\section{公式题}

\subsection{识别信号}

题面给出明确计算规则，例如：

\begin{itemize}
  \item 根据若干输入计算费用、得分、距离、面积。
  \item 分段收费。
  \item 按比例、百分比、折扣计算。
  \item 对结果取整、向上取整、向下取整。
\end{itemize}

\subsection{公式题流程}

\begin{enumerate}
  \item 先在草稿纸上写出数学式。
  \item 判断每个中间变量是否可能溢出。
  \item 判断除法是整数除法还是小数除法。
  \item 判断是否需要四舍五入、向上取整、向下取整。
  \item 用样例手算一遍。
\end{enumerate}

\subsection{整数除法和取整}

公式题中最容易错的是“除法到底按什么规则取整”。先分清三种情况：

\begin{longtable}{p{4cm}p{4.8cm}p{5.2cm}}
\toprule
题目要求 & 常用写法 & 适用条件 \\
\midrule
普通整数除法，向 0 截断 & \Code{a / b} & C++ 默认规则，负数时不是向下取整。 \\
正整数向上取整 & \Code{(a + b - 1) / b} & 只适合 $a \ge 0, b > 0$。 \\
数学意义向下/向上取整 & 使用 \Code{floor\_div} / \Code{ceil\_div} & 适合可能出现负数的公式。 \\
\bottomrule
\end{longtable}

正整数向上取整最常见，例如 $a$ 个物品每组放 $b$ 个：

\begin{lstlisting}
long long ceil_div_positive(long long a, long long b) {
    return (a + b - 1) / b;
}
\end{lstlisting}

若 $a$ 可能为负，或题目明确要求数学意义的向下/向上取整，使用通用写法：

\begin{lstlisting}
long long floor_div(long long a, long long b) {
    long long q = a / b;
    long long r = a % b;
    if (r != 0 && ((r > 0) != (b > 0))) q--;
    return q;
}

long long ceil_div(long long a, long long b) {
    long long q = a / b;
    long long r = a % b;
    if (r != 0 && ((r > 0) == (b > 0))) q++;
    return q;
}
\end{lstlisting}

\Warn{不要随便用浮点数做整数取整题。}如果题目全是整数，通常可以用整数公式避免误差；特别是数值接近 $10^{18}$ 时，\Code{double} 可能无法精确表示每一个整数。

\subsection{分段计算模板}

\begin{lstlisting}
int x;
cin >> x;

int ans = 0;
if (x <= 10) {
    ans = x * 2;
} else if (x <= 20) {
    ans = 10 * 2 + (x - 10) * 3;
} else {
    ans = 10 * 2 + 10 * 3 + (x - 20) * 5;
}

cout << ans << '\n';
\end{lstlisting}

注意：

\begin{itemize}
  \item 边界是 \Code{<=} 还是 \Code{<}。
  \item 每一段是否包含端点。
  \item 前面几段的费用是否已经累计。
\end{itemize}

\subsection{分类讨论模板}

A 题经常给出若干规则，按条件输出不同结果。

\begin{lstlisting}
int x;
cin >> x;

if (x < 0) {
    cout << "negative\n";
} else if (x == 0) {
    cout << "zero\n";
} else {
    cout << "positive\n";
}
\end{lstlisting}

分类讨论检查：

\begin{itemize}
  \item 条件是否覆盖所有情况。
  \item 条件顺序是否正确。
  \item 边界值属于哪一类。
  \item 是否存在重叠条件。
\end{itemize}

\subsection{绝对值和距离}

\begin{lstlisting}
int a, b;
cin >> a >> b;
cout << abs(a - b) << '\n';
\end{lstlisting}

如果 $a,b$ 很大：

\begin{lstlisting}
long long a, b;
cin >> a >> b;
cout << llabs(a - b) << '\n';
\end{lstlisting}

\subsection{取模}

\begin{lstlisting}
long long ans = 0;
const long long MOD = 1000000007;
ans = (ans + x) % MOD;
ans = ans * y % MOD;
\end{lstlisting}

如果可能出现负数：

\begin{lstlisting}
long long r = x % MOD;
if (r < 0) r += MOD;
\end{lstlisting}

\subsection{小数和百分比}

保留小数：

\begin{lstlisting}
double x;
cin >> x;
cout << fixed << setprecision(2) << x << '\n';
\end{lstlisting}

整数转百分比：

\begin{lstlisting}
int good, total;
cin >> good >> total;
double rate = 100.0 * good / total;
cout << fixed << setprecision(2) << rate << "%\n";
\end{lstlisting}

平均数：

\begin{lstlisting}
long long sum = 0;
int n;
cin >> n;
for (int i = 0; i < n; i++) {
    int x;
    cin >> x;
    sum += x;
}
double avg = 1.0 * sum / n;
cout << fixed << setprecision(1) << avg << '\n';
\end{lstlisting}

\Warn{注意：}\Code{sum / n} 如果两边都是整数，会做整数除法。写成 \Code{1.0 * sum / n}。

\section{排序和排名}

\subsection{识别信号}

题面出现：

\begin{itemize}
  \item 按成绩排序、按编号排序、按时间排序。
  \item 输出第 $k$ 个。
  \item 排名、名次、优先级。
  \item 相同分数时按另一个条件排序。
\end{itemize}

\subsection{单个数组排序}

\begin{lstlisting}
int n;
cin >> n;
vector<int> a(n);
for (int i = 0; i < n; i++) cin >> a[i];

sort(a.begin(), a.end());

for (int x : a) {
    cout << x << ' ';
}
cout << '\n';
\end{lstlisting}

降序：

\begin{lstlisting}
sort(a.begin(), a.end(), greater<int>());
\end{lstlisting}

\subsection{结构体排序}

\begin{lstlisting}
struct Student {
    int id;
    int score;
};

int n;
cin >> n;
vector<Student> a(n);
for (int i = 0; i < n; i++) {
    cin >> a[i].id >> a[i].score;
}

sort(a.begin(), a.end(), [](const Student& x, const Student& y) {
    if (x.score != y.score) return x.score > y.score;
    return x.id < y.id;
});

for (const auto& s : a) {
    cout << s.id << ' ' << s.score << '\n';
}
\end{lstlisting}

\subsection{排名常见规则}

排名题一定要看清楚“并列名次”。

普通排序后第 $k$ 个：

\begin{lstlisting}
sort(a.begin(), a.end(), greater<int>());
cout << a[k - 1] << '\n';
\end{lstlisting}

并列排名示例：

\begin{lstlisting}
vector<int> score = {100, 90, 90, 80};
int rank = 1;
for (int i = 0; i < (int)score.size(); i++) {
    if (i > 0 && score[i] != score[i - 1]) {
        rank = i + 1;
    }
    cout << score[i] << ' ' << rank << '\n';
}
\end{lstlisting}

\Warn{排序题高危点：}比较器里相等时不能返回 \Code{true}，需要继续比较或返回 \Code{false}。

\subsection{按字符串字典序排序}

\begin{lstlisting}
int n;
cin >> n;
vector<string> s(n);
for (int i = 0; i < n; i++) cin >> s[i];

sort(s.begin(), s.end());

for (string x : s) {
    cout << x << '\n';
}
\end{lstlisting}

\subsection{先去重再排序}

\begin{lstlisting}
sort(a.begin(), a.end());
a.erase(unique(a.begin(), a.end()), a.end());
\end{lstlisting}

\subsection{找第 k 小和第 k 大}

\begin{lstlisting}
sort(a.begin(), a.end());
cout << a[k - 1] << '\n'; // kth smallest

sort(a.begin(), a.end(), greater<int>());
cout << a[k - 1] << '\n'; // kth largest
\end{lstlisting}

\Warn{第 k 个通常是 1 下标，所以用 \Code{k - 1}。}

\section{字符串处理}

\subsection{识别信号}

题面出现：

\begin{itemize}
  \item 字符、单词、大小写。
  \item 判断前缀、后缀、是否包含。
  \item 统计某字符出现次数。
  \item 按空格、逗号、冒号分割。
  \item 文件路径、命令、表达式。
\end{itemize}

\subsection{逐字符扫描}

\begin{lstlisting}
string s;
cin >> s;

int digits = 0;
int letters = 0;
for (char c : s) {
    if (isdigit(c)) digits++;
    if (isalpha(c)) letters++;
}

cout << digits << ' ' << letters << '\n';
\end{lstlisting}

\subsection{大小写转换}

\begin{lstlisting}
string s;
cin >> s;
for (char& c : s) {
    c = tolower(c);
}
cout << s << '\n';
\end{lstlisting}

\subsection{查找子串}

\begin{lstlisting}
string s, t;
cin >> s >> t;
if (s.find(t) != string::npos) {
    cout << "YES\n";
} else {
    cout << "NO\n";
}
\end{lstlisting}

\subsection{按空格分割一整行}

\begin{lstlisting}
string line;
getline(cin, line);

stringstream ss(line);
string word;
while (ss >> word) {
    cout << word << '\n';
}
\end{lstlisting}

如果前面读过整数：

\begin{lstlisting}
int n;
cin >> n;
cin.ignore();

string line;
getline(cin, line);
\end{lstlisting}

\subsection{按指定字符分割}

\begin{lstlisting}
vector<string> split(const string& s, char sep) {
    vector<string> res;
    string cur;
    for (char c : s) {
        if (c == sep) {
            res.push_back(cur);
            cur.clear();
        } else {
            cur.push_back(c);
        }
    }
    res.push_back(cur);
    return res;
}
\end{lstlisting}

\Warn{字符串题常见坑：}空字符串、连续分隔符、行尾换行、下标从 0 开始。

\subsection{判断前缀和后缀}

前缀：

\begin{lstlisting}
bool startsWith(const string& s, const string& p) {
    if (p.size() > s.size()) return false;
    for (int i = 0; i < (int)p.size(); i++) {
        if (s[i] != p[i]) return false;
    }
    return true;
}
\end{lstlisting}

后缀：

\begin{lstlisting}
bool endsWith(const string& s, const string& p) {
    if (p.size() > s.size()) return false;
    int offset = (int)s.size() - (int)p.size();
    for (int i = 0; i < (int)p.size(); i++) {
        if (s[offset + i] != p[i]) return false;
    }
    return true;
}
\end{lstlisting}

\subsection{统计每个小写字母}

\begin{lstlisting}
string s;
cin >> s;
vector<int> cnt(26, 0);
for (char c : s) {
    if ('a' <= c && c <= 'z') {
        cnt[c - 'a']++;
    }
}
\end{lstlisting}

\subsection{反转字符串}

\begin{lstlisting}
string s;
cin >> s;
reverse(s.begin(), s.end());
cout << s << '\n';
\end{lstlisting}

\subsection{回文判断}

\begin{lstlisting}
bool isPalindrome(const string& s) {
    int l = 0, r = (int)s.size() - 1;
    while (l < r) {
        if (s[l] != s[r]) return false;
        l++;
        r--;
    }
    return true;
}
\end{lstlisting}

\section{进制处理}

\subsection{识别信号}

题面出现：

\begin{itemize}
  \item 二进制、八进制、十六进制。
  \item 把某进制转为十进制。
  \item 把十进制转为某进制。
  \item 位、编码、补码、掩码。
\end{itemize}

\subsection{任意进制转十进制}

\begin{lstlisting}
long long toDecimal(const string& s, int base) {
    long long ans = 0;
    for (char c : s) {
        int v;
        if ('0' <= c && c <= '9') v = c - '0';
        else if ('A' <= c && c <= 'F') v = c - 'A' + 10;
        else v = c - 'a' + 10;
        ans = ans * base + v;
    }
    return ans;
}
\end{lstlisting}

\subsection{十进制转任意进制}

\begin{lstlisting}
string fromDecimal(long long x, int base) {
    if (x == 0) return "0";
    string digits = "0123456789ABCDEF";
    string res;
    while (x > 0) {
        res.push_back(digits[x % base]);
        x /= base;
    }
    reverse(res.begin(), res.end());
    return res;
}
\end{lstlisting}

\subsection{位运算入门}

\begin{lstlisting}
if (x & 1) {
    // x is odd
}

if (x & (1 << k)) {
    // bit k is 1
}

x |= (1 << k);   // set bit k to 1
x &= ~(1 << k);  // set bit k to 0
\end{lstlisting}

\Warn{如果 $k \ge 31$，写 \Code{1LL << k}，不要写 \Code{1 << k}。}

\subsection{进制题检查}

\begin{itemize}
  \item 输入数字是否可能很长。
  \item 转成十进制是否超过 \Code{long long}。
  \item 十六进制字母是大写还是小写。
  \item 输出是否要求补前导零。
  \item 二进制位编号从 0 还是 1 开始。
\end{itemize}

\section{日期时间}

\subsection{识别信号}

题面出现：

\begin{itemize}
  \item 年、月、日、星期。
  \item 给定日期求过几天。
  \item 两个日期之间的天数。
  \item 时间格式如 \Code{HH:MM:SS}。
\end{itemize}

\subsection{闰年}

\begin{lstlisting}
bool leap(int y) {
    return (y % 400 == 0) || (y % 4 == 0 && y % 100 != 0);
}
\end{lstlisting}

\subsection{每月天数}

\begin{lstlisting}
int mdays(int y, int m) {
    int d[] = {0,31,28,31,30,31,30,31,31,30,31,30,31};
    if (m == 2 && leap(y)) return 29;
    return d[m];
}
\end{lstlisting}

\subsection{下一天}

\begin{lstlisting}
void nextDay(int& y, int& m, int& d) {
    d++;
    if (d > mdays(y, m)) {
        d = 1;
        m++;
        if (m > 12) {
            m = 1;
            y++;
        }
    }
}
\end{lstlisting}

\subsection{时间转秒}

\begin{lstlisting}
int toSeconds(int h, int m, int s) {
    return h * 3600 + m * 60 + s;
}
\end{lstlisting}

\Warn{日期题先看范围。}如果只跨几年，可以逐日模拟；如果年份很大，要转为天数公式或按周期处理。

\section{简单模拟}

\subsection{识别信号}

题面出现：

\begin{itemize}
  \item 按规则一步一步执行。
  \item 有当前位置、方向、状态。
  \item 有若干操作命令。
  \item 有网格移动，但不要求最短路。
  \item 有游戏规则、计分规则。
\end{itemize}

\subsection{方向数组}

四方向：

\begin{lstlisting}
int dx[4] = {-1, 0, 1, 0};
int dy[4] = {0, 1, 0, -1};
\end{lstlisting}

八方向：

\begin{lstlisting}
int dx[8] = {-1,-1,-1,0,0,1,1,1};
int dy[8] = {-1,0,1,-1,1,-1,0,1};
\end{lstlisting}

判断是否在网格内：

\begin{lstlisting}
bool inside(int x, int y, int n, int m) {
    return 0 <= x && x < n && 0 <= y && y < m;
}
\end{lstlisting}

\subsection{按命令移动}

\begin{lstlisting}
int x = 0, y = 0;
string ops;
cin >> ops;

for (char op : ops) {
    if (op == 'U') x--;
    else if (op == 'D') x++;
    else if (op == 'L') y--;
    else if (op == 'R') y++;
}

cout << x << ' ' << y << '\n';
\end{lstlisting}

\subsection{模拟题写法建议}

\begin{itemize}
  \item 状态变量先列出来，例如位置、方向、分数、时间。
  \item 每个操作只改变少量状态。
  \item 复杂规则拆成函数。
  \item 样例过后自己造最小边界和最大边界。
\end{itemize}

\section{坐标和距离入门}

A 题中若出现坐标，通常是简单计算，不一定是图搜索。

\subsection{二维点}

\begin{lstlisting}
struct Point {
    long long x;
    long long y;
};

Point a, b;
cin >> a.x >> a.y >> b.x >> b.y;
\end{lstlisting}

\subsection{曼哈顿距离}

只能上下左右走时，常见距离：

\begin{lstlisting}
long long dist = llabs(a.x - b.x) + llabs(a.y - b.y);
cout << dist << '\n';
\end{lstlisting}

\subsection{欧几里得距离}

直线距离：

\begin{lstlisting}
double dx = a.x - b.x;
double dy = a.y - b.y;
double dist = sqrt(dx * dx + dy * dy);
cout << fixed << setprecision(6) << dist << '\n';
\end{lstlisting}

\subsection{坐标题检查}

\begin{itemize}
  \item 坐标可能为负数。
  \item 坐标差的平方可能溢出，必要时用 \Code{long long} 或 \Code{double}。
  \item 题目要求的是曼哈顿距离还是直线距离。
  \item 输出是否需要小数。
\end{itemize}

\section{二维表格和矩阵入门}

A 题可能出现很简单的二维表格，例如统计每行和、每列和、找最大值。

\subsection{读入 n 行 m 列}

\begin{lstlisting}
int n, m;
cin >> n >> m;
vector<vector<int>> a(n, vector<int>(m));
for (int i = 0; i < n; i++) {
    for (int j = 0; j < m; j++) {
        cin >> a[i][j];
    }
}
\end{lstlisting}

\subsection{每行求和}

\begin{lstlisting}
for (int i = 0; i < n; i++) {
    long long sum = 0;
    for (int j = 0; j < m; j++) {
        sum += a[i][j];
    }
    cout << sum << '\n';
}
\end{lstlisting}

\subsection{每列求和}

\begin{lstlisting}
for (int j = 0; j < m; j++) {
    long long sum = 0;
    for (int i = 0; i < n; i++) {
        sum += a[i][j];
    }
    cout << sum << '\n';
}
\end{lstlisting}

\subsection{找矩阵最大值}

\begin{lstlisting}
int bx = 0, by = 0;
for (int i = 0; i < n; i++) {
    for (int j = 0; j < m; j++) {
        if (a[i][j] > a[bx][by]) {
            bx = i;
            by = j;
        }
    }
}
cout << bx << ' ' << by << ' ' << a[bx][by] << '\n';
\end{lstlisting}

如果题目要求第几行第几列，通常输出 \Code{bx + 1} 和 \Code{by + 1}。

\section{频率统计}

\subsection{值域小：用数组}

如果数值范围是 $0$ 到 $1000$：

\begin{lstlisting}
int cnt[1001] = {};
int n;
cin >> n;
for (int i = 0; i < n; i++) {
    int x;
    cin >> x;
    cnt[x]++;
}
\end{lstlisting}

\subsection{值域大或 key 是字符串：用 map}

\begin{lstlisting}
int n;
cin >> n;
map<string, int> cnt;
for (int i = 0; i < n; i++) {
    string s;
    cin >> s;
    cnt[s]++;
}

for (auto [s, c] : cnt) {
    cout << s << ' ' << c << '\n';
}
\end{lstlisting}

\subsection{只需要快速统计：unordered\_map}

\begin{lstlisting}
unordered_map<int, int> cnt;
for (int x : a) {
    cnt[x]++;
}
\end{lstlisting}

\Warn{如果题目要求按字典序输出，不能用 \Code{unordered\_map} 直接遍历。}

\section{区间入门}

A 题中也可能出现简单区间。先看数据范围。

\subsection{数据小：直接循环}

\begin{lstlisting}
int l, r;
cin >> l >> r;
long long sum = 0;
for (int i = l; i <= r; i++) {
    sum += a[i];
}
\end{lstlisting}

\subsection{多次区间和：前缀和}

\begin{lstlisting}
int n, q;
cin >> n >> q;
vector<long long> pre(n + 1, 0);
for (int i = 1; i <= n; i++) {
    long long x;
    cin >> x;
    pre[i] = pre[i - 1] + x;
}

while (q--) {
    int l, r;
    cin >> l >> r;
    cout << pre[r] - pre[l - 1] << '\n';
}
\end{lstlisting}

注意：

\begin{itemize}
  \item 上面模板假设下标从 1 开始。
  \item 如果题目给的是 0 下标，要改公式。
  \item 前缀和数组用 \Code{long long}。
\end{itemize}

\section{A 题输出格式专项}

提交前逐项对照：

\begin{itemize}
  \item 题目要求输出 \Code{YES/NO} 还是 \Code{Yes/No}。
  \item 是否每个答案一行。
  \item 多个数之间是空格还是换行。
  \item 最后一项后面是否允许空格。一般允许，但不依赖。
  \item 小数保留几位。
  \item 是否要求百分号、冒号、逗号等符号。
\end{itemize}

输出数组推荐写法：

\begin{lstlisting}
for (int i = 0; i < n; i++) {
    if (i) cout << ' ';
    cout << a[i];
}
cout << '\n';
\end{lstlisting}

\section{A 题新手错误对照表}

\begin{longtable}{p{4cm}p{5cm}p{5cm}}
\toprule
现象 & 可能原因 & 快速检查 \\
\midrule
样例输出差一个数 & 循环次数错 & 检查 \Code{i < n} 和 \Code{i <= n}。 \\
样例数字对但格式错 & 空格、换行、大小写错误 & 对照输出格式逐字符看。 \\
小数据对，大数据错 & \Code{int} 溢出 & 把和、乘积、答案改为 \Code{long long}。 \\
排序后顺序怪 & 比较器错误 & 相等时返回 \Code{false} 或继续比较。 \\
读入字符串为空 & \Code{cin} 后直接 \Code{getline} & 加 \Code{cin.ignore()}。 \\
程序崩溃 & 数组越界或空容器取值 & 查 \Code{a[i]}、\Code{front/top/back}。 \\
答案少算最后一项 & 区间端点没包含 & 检查 \Code{<= r} 还是 \Code{< r}。 \\
答案多算一项 & 下标起点错 & 检查 0 下标和 1 下标是否混用。 \\
找子串时明明在开头却判断失败 & 错用 \Code{if (s.find(t))} & 改为和 \Code{string::npos} 比较。 \\
多组数据第二组错 & 没清空状态 & 每组重置数组、map、答案。 \\
\bottomrule
\end{longtable}

\section{A 题推荐代码组织}

简单题可以直接写在 \Code{main}，但如果题目有多组数据，推荐写 \Code{solve()}。

\begin{lstlisting}
#include <bits/stdc++.h>
using namespace std;
using ll = long long;

void solve() {
    int n;
    cin >> n;
    vector<int> a(n);
    for (int i = 0; i < n; i++) {
        cin >> a[i];
    }

    ll sum = 0;
    for (int x : a) {
        sum += x;
    }
    cout << sum << '\n';
}

int main() {
    ios::sync_with_stdio(false);
    cin.tie(nullptr);

    solve();
    return 0;
}
\end{lstlisting}

如果题目有 $T$ 组：

\begin{lstlisting}
int main() {
    ios::sync_with_stdio(false);
    cin.tie(nullptr);

    int T;
    cin >> T;
    while (T--) {
        solve();
    }
    return 0;
}
\end{lstlisting}

\Warn{如果题目没有多组数据，不要自己读一个 T。}

\section{A 题自造边界样例}

样例通过后，至少测试：

\begin{itemize}
  \item 最小输入：$n=0$ 或 $n=1$，如果题目允许。
  \item 最大输入附近：最大数值、最大长度。
  \item 全相同数据。
  \item 全为 0。
  \item 全为负数，如果题目允许。
  \item 排序中有相同关键字。
  \item 字符串为空或只有一个字符，如果题目允许。
  \item 日期为 2 月 28 日、2 月 29 日、12 月 31 日。
\end{itemize}

\section{A 题本地调试流程}

如果样例不过，不要盲目改代码。按下面顺序查。

\subsection{第一步：确认读入}

临时输出读到的变量：

\begin{lstlisting}
cerr << "n=" << n << '\n';
for (int i = 0; i < n; i++) {
    cerr << "a[" << i << "]=" << a[i] << '\n';
}
\end{lstlisting}

\Warn{提交前必须删除或注释调试输出。}虽然 \Code{cerr} 通常不进入标准输出，但考场上不要依赖这一点，最终代码保持干净。

\subsection{第二步：确认中间结果}

例如统计题：

\begin{lstlisting}
cerr << "sum=" << sum << " cnt=" << cnt << '\n';
\end{lstlisting}

排序题：

\begin{lstlisting}
for (int x : a) cerr << x << ' ';
cerr << '\n';
\end{lstlisting}

\subsection{第三步：用最小样例}

把输入改成最小情况：

\begin{lstlisting}
1
5
\end{lstlisting}

如果最小样例都错，通常是读入、初始化或输出格式错。

\subsection{第四步：检查边界}

\begin{itemize}
  \item 循环是否少跑一次。
  \item 数组是否越界。
  \item 是否把 $n$ 当成最大下标。
  \item 是否忘记处理最后一个元素。
  \item 是否忘记初始化答案变量。
\end{itemize}

\section{A 题常见边界库}

做完样例后，根据题型选 2 到 3 个边界测。

\begin{longtable}{p{4cm}p{5cm}p{5cm}}
\toprule
题型 & 建议边界 & 目的 \\
\midrule
统计 & $n=1$，全 0，全负数，全相同 & 检查初始化和条件判断。 \\
最大最小 & 全负数、最大值在第一项、最大值在最后一项 & 检查初始值和循环。 \\
排序 & 已有序、逆序、全相等、有重复关键字 & 检查比较器和稳定性。 \\
字符串 & 长度 1、全数字、全字母、含空格、子串在开头 & 检查读取和 \Code{find}。 \\
日期 & 2 月 28 日、闰年 2 月 29 日、12 月 31 日 & 检查闰年和跨年。 \\
进制 & 0、最大位、含 A-F、小写字母 & 检查特殊值和字符转换。 \\
区间 & $l=r$，$l=1$，$r=n$ & 检查前缀和端点。 \\
二维表 & $1 \times 1$，单行，单列 & 检查双重循环。 \\
\bottomrule
\end{longtable}

\section{A 题手写模板清单}

这些模板建议单独练到不用看也能写。

\begin{longtable}{p{4cm}p{4cm}p{6cm}}
\toprule
模板 & 必须熟练程度 & 关联章节 \\
\midrule
读入 $n$ 个数 & 必须会默写 & 输入形态、vector \\
一遍求和、最大、最小 & 必须会默写 & 简单统计 \\
统计满足条件的个数 & 必须会默写 & 简单统计 \\
结构体排序 & 必须照着能写 & 排序和排名 \\
字符串逐字符扫描 & 必须会默写 & 字符串处理 \\
\Code{getline} + \Code{stringstream} & 必须照着能写 & 字符串处理 \\
进制转十进制 & 照着能写 & 进制处理 \\
闰年和每月天数 & 照着能写 & 日期时间 \\
方向数组 & 照着能写 & 简单模拟 \\
二维矩阵读入 & 必须会默写 & 二维表格 \\
前缀和查询 & 照着能写 & 区间入门 \\
\bottomrule
\end{longtable}

\section{A 题临场优先级}

如果你在 A 题上紧张，按这个优先级保证分数：

\begin{enumerate}
  \item 先保证读入正确。
  \item 再保证样例能手算清楚。
  \item 再写最直接的做法，不要优化。
  \item 再处理边界。
  \item 最后才考虑代码好不好看。
\end{enumerate}

如果有两个写法：

\begin{itemize}
  \item 一个写法短但你不熟。
  \item 一个写法长但你确定。
\end{itemize}

A 题选第二个。

\section{A 题提交前总检查}

\begin{enumerate}
  \item 输入是否完整读入。
  \item 输出是否完全符合题目。
  \item 是否有调试输出。
  \item 是否需要 \Code{long long}。
  \item 数组是否越界。
  \item 下标从 0 还是 1 开始是否统一。
  \item 排序比较器是否正确。
  \item 字符串读取是否受空格影响。
  \item 小数精度是否正确。
  \item 边界样例是否测试过。
\end{enumerate}

\section{A 题卡住时怎么办}

\begin{itemize}
  \item 读不懂题：重看输入输出样例，手算样例。
  \item 样例不过：打印中间变量在本地查错，提交前删掉。
  \item 怀疑溢出：把答案和中间量改成 \Code{long long}。
  \item 怀疑下标：画出长度为 3 的小数组，手动走一遍。
  \item 超过 45 分钟仍未解决：保存代码，先做 B 或 C，之后回来。
\end{itemize}

\chapter{B 题决策树与高频算法入口}

\section{B 题定位}

B 题通常是从“会写代码”进入“会选算法”的第一道题。它一般不会要求特别深的算法证明，但常常要求你识别出比暴力更快的做法。

\begin{itemize}
  \item 建议目标时间：40--70 分钟。
  \item 目标分数：尽量 100 分。
  \item 高频主题：排序、哈希统计、前缀和、差分、二分、简单 BFS、简单图、简单 DP。
  \item 核心风险：样例能过的暴力在大数据超时。
\end{itemize}

\Warn{B 题第一原则：}读完题后先看数据范围，再决定能不能暴力。不要因为 A 题刚做完很顺，就直接写双重循环。

\section{B 题开题流程}

\begin{enumerate}
  \item 圈出数据范围：$n,m,q$ 分别多大。
  \item 圈出操作次数：是否有很多查询、很多修改。
  \item 判断暴力复杂度：如果每次操作都扫描一遍，总次数是多少。
  \item 找题面信号：区间、排序、统计、最短、连通、依赖、答案范围。
  \item 查 B 题决策树，定位高频模板。
  \item 如果满分做法明确，直接写；如果不明确，先写可拿部分分的暴力。
\end{enumerate}

\section{B 题总决策树}

\begin{longtable}{p{4cm}p{5cm}p{5cm}}
\toprule
题面信号 & 优先怀疑 & 首选模板 \\
\midrule
多次询问区间和、区间最大统计 & 前缀和或预处理 & 一维/二维前缀和 \\
多次区间加，最后查询结果 & 差分 & 一维/二维差分 \\
区间修改和区间查询交替出现 & 树状数组或线段树 & 数据结构章节，B 题先看是否可简化 \\
统计出现次数、找重复、频率 & 哈希统计 & 数组计数、map、unordered\_map \\
排序后选择、排名、贪心安排 & 排序与扫描 & sort + 自定义比较器 \\
前后两边信息、删除一个元素后的最优值 & 前后缀预处理 & 前缀最大/最小、后缀最大/最小 \\
余数、整除、模 $k$ 相同 & 余数统计 & cnt[mod]、map 统计 \\
第一个满足条件、答案具有单调性 & 二分或二分答案 & lower\_bound / check \\
最少步数、网格移动、迷宫 & BFS & 队列、距离数组 \\
合并集合、判断是否同组 & 并查集 & DSU \\
依赖关系、先后顺序 & 拓扑排序 & 入度 + 队列 \\
小规模任意两点距离 & Floyd & 三重循环 \\
正权图最短路 & Dijkstra & 堆优化 Dijkstra \\
选择若干、容量限制、最优值 & 简单 DP & 背包或线性 DP \\
\bottomrule
\end{longtable}

\section{B 题数据范围判断}

\begin{longtable}{p{4cm}p{5cm}p{5cm}}
\toprule
数据范围 & 暴力是否可行 & 常见替代 \\
\midrule
$n \le 1000$ & $O(n^2)$ 常常可行 & 双重循环、简单 DP \\
$n \le 10^5$ & $O(n^2)$ 通常不可行 & 排序、哈希、前缀和、二分 \\
$q \le 10^5$ 且每次查区间 & 每次扫描不可行 & 前缀和、树状数组、线段树 \\
二维 $n,m \le 1000$ & $O(nm)$ 可行，$O(n^2m^2)$ 不可行 & 二维前缀和 \\
图点数 $\le 500$ & Floyd 可能可行 & Floyd \\
图点数、边数 $\le 10^5$ & Floyd 不可行 & BFS、Dijkstra、并查集 \\
值域 $\le 10^6$ & 可用计数数组 & 数组计数 \\
值域很大或字符串 key & 不能开值域数组 & map / unordered\_map \\
\bottomrule
\end{longtable}

\section{B 题关键词到动作}

\begin{longtable}{p{3.8cm}p{4.8cm}p{5.4cm}}
\toprule
看到 & 立刻想 & 小心 \\
\midrule
“多次询问” & 预处理，不能每次从头算 & 查询次数 $q$ 是否很大。 \\
“区间 [l,r] 的和” & 前缀和 & 下标和 \Code{pre[r]-pre[l-1]}。 \\
“子矩阵和” & 二维前缀和 & 公式四项别写错。 \\
“区间加” & 差分 & 右端点后一位是否存在。 \\
“出现次数” & 计数数组或 map & 值域是否能开数组。 \\
“去重后数量” & set 或 sort+unique & 是否需要保持顺序。 \\
“第一个不小于” & lower\_bound & 容器必须有序。 \\
“最小的最大值” & 二分答案 & 必须能写 check。 \\
“删除一个后”“左边右边” & 前后缀预处理 & 注意边界位置。 \\
“除以 k 的余数”“能否整除” & 余数统计 & 负数取模要修正。 \\
“尽量多选”“最早结束”“最少次数覆盖” & 贪心 & 必须按正确关键字排序。 \\
“最短步数” & BFS & BFS 适合每步代价相同。 \\
“连通、合并” & 并查集 & 初始化每个点的父亲。 \\
“前置条件、依赖” & 拓扑排序 & 入度更新。 \\
\bottomrule
\end{longtable}

\section{从暴力到优化的转换表}

B 题经常不是完全不会，而是暴力会写但会超时。下面这张表用于把暴力做法改成常见优化。

\begin{longtable}{p{4.2cm}p{4.8cm}p{5cm}}
\toprule
暴力想法 & 为什么危险 & 常见优化 \\
\midrule
每次查询都遍历 $[l,r]$ & $O(nq)$，$q$ 大时超时 & 前缀和、树状数组 \\
每次区间加都逐个元素修改 & $O(nq)$ & 差分、线段树 \\
双重循环找两数关系 & $O(n^2)$ & 排序、哈希、双指针 \\
每次判断是否出现都扫描数组 & $O(nq)$ & set、unordered\_set、排序后二分 \\
每次找最大值都扫描 & $O(nq)$ & priority\_queue、set、预处理 \\
枚举所有子矩阵 & $O(n^2m^2)$ 或更高 & 二维前缀和 \\
图上每对点都 BFS & 点数大时超时 & 看是否只需单源；正权用 Dijkstra \\
字符串反复拼接或插入 & 长字符串会很慢 & 线性扫描、vector 存片段 \\
\bottomrule
\end{longtable}

\section{前缀和}

\subsection{识别信号}

题面出现：

\begin{itemize}
  \item 多次查询区间和。
  \item 查询前 $k$ 项和。
  \item 子数组和、某段累计值。
  \item 数组不修改，或者所有查询都在修改之前/之后。
\end{itemize}

\subsection{一维前缀和模板}

使用 1 下标最稳：

\begin{lstlisting}
int n, q;
cin >> n >> q;
vector<long long> pre(n + 1, 0);

for (int i = 1; i <= n; i++) {
    long long x;
    cin >> x;
    pre[i] = pre[i - 1] + x;
}

while (q--) {
    int l, r;
    cin >> l >> r;
    cout << pre[r] - pre[l - 1] << '\n';
}
\end{lstlisting}

\subsection{0 下标前缀和}

如果原数组是 0 下标，也可以让 \Code{pre[i+1]} 表示前 $i+1$ 个数之和：

\begin{lstlisting}
vector<int> a(n);
vector<long long> pre(n + 1, 0);
for (int i = 0; i < n; i++) {
    cin >> a[i];
    pre[i + 1] = pre[i] + a[i];
}

// sum of a[l..r], 0-indexed
long long ans = pre[r + 1] - pre[l];
\end{lstlisting}

\subsection{二维前缀和}

适合多次查询子矩阵和。

\begin{lstlisting}
int n, m, q;
cin >> n >> m >> q;
vector<vector<long long>> pre(n + 1, vector<long long>(m + 1, 0));

for (int i = 1; i <= n; i++) {
    for (int j = 1; j <= m; j++) {
        long long x;
        cin >> x;
        pre[i][j] = pre[i - 1][j] + pre[i][j - 1]
                  - pre[i - 1][j - 1] + x;
    }
}

while (q--) {
    int x1, y1, x2, y2;
    cin >> x1 >> y1 >> x2 >> y2;
    long long ans = pre[x2][y2] - pre[x1 - 1][y2]
                  - pre[x2][y1 - 1] + pre[x1 - 1][y1 - 1];
    cout << ans << '\n';
}
\end{lstlisting}

\subsection{前缀和易错点}

\begin{itemize}
  \item 查询端点是否从 1 开始。
  \item \Code{pre[0]} 必须是 0。
  \item 区间公式是一维两项、二维四项。
  \item 和可能溢出，使用 \Code{long long}。
  \item 如果中途有修改，普通前缀和通常不适用。
\end{itemize}

\section{前后缀预处理}

\subsection{识别信号}

题面出现：

\begin{itemize}
  \item 删除一个元素后，问剩余部分的最大、最小、和。
  \item 对每个位置，问它左边和右边的信息。
  \item 需要快速知道前 $i$ 个的最大值、后 $i$ 个的最大值。
\end{itemize}

\subsection{前缀最大和后缀最大}

\begin{lstlisting}
int n;
cin >> n;
vector<int> a(n + 1);
for (int i = 1; i <= n; i++) cin >> a[i];

vector<int> preMax(n + 2, -1000000000);
vector<int> sufMax(n + 2, -1000000000);

for (int i = 1; i <= n; i++) {
    preMax[i] = max(preMax[i - 1], a[i]);
}
for (int i = n; i >= 1; i--) {
    sufMax[i] = max(sufMax[i + 1], a[i]);
}

// max value except position i
for (int i = 1; i <= n; i++) {
    int best = max(preMax[i - 1], sufMax[i + 1]);
    cout << best << '\n';
}
\end{lstlisting}

\subsection{前后缀易错点}

\begin{itemize}
  \item 数组多开两位，方便处理 \Code{i-1} 和 \Code{i+1}。
  \item 初始值要根据题意设置，最大值用很小值，最小值用很大值。
  \item 如果是求和，使用 \Code{long long}。
  \item 删除第一个或最后一个元素时，一边为空，要提前定义空边的值。
\end{itemize}

\section{差分}

\subsection{识别信号}

题面出现：

\begin{itemize}
  \item 很多次区间加。
  \item 最后统一输出每个位置的值。
  \item 修改多，查询少或只在最后查询。
\end{itemize}

\subsection{一维差分模板}

区间 $[l,r]$ 加 $v$：

\begin{lstlisting}
int n, q;
cin >> n >> q;
vector<long long> diff(n + 2, 0);

while (q--) {
    int l, r;
    long long v;
    cin >> l >> r >> v;
    diff[l] += v;
    diff[r + 1] -= v;
}

long long cur = 0;
for (int i = 1; i <= n; i++) {
    cur += diff[i];
    cout << cur << '\n';
}
\end{lstlisting}

如果原数组已有初值：

\begin{lstlisting}
vector<long long> a(n + 1), diff(n + 2, 0);
for (int i = 1; i <= n; i++) cin >> a[i];
for (int i = 1; i <= n; i++) {
    diff[i] += a[i] - a[i - 1];
}
\end{lstlisting}

\subsection{差分易错点}

\begin{itemize}
  \item 数组要开到 \Code{n+2}，因为会访问 \Code{r+1}。
  \item 差分适合“多次修改后统一求结果”。
  \item 如果每次修改后立刻查询区间和，普通差分不够。
  \item 修改值可能很大，用 \Code{long long}。
\end{itemize}

\subsection{二维差分入口}

适合很多次对子矩阵加值，最后输出整个矩阵。

对矩形 $(x1,y1)$ 到 $(x2,y2)$ 加 $v$：

\begin{lstlisting}
int n, m, q;
cin >> n >> m >> q;
vector<vector<long long>> diff(n + 2, vector<long long>(m + 2, 0));

while (q--) {
    int x1, y1, x2, y2;
    long long v;
    cin >> x1 >> y1 >> x2 >> y2 >> v;
    diff[x1][y1] += v;
    diff[x2 + 1][y1] -= v;
    diff[x1][y2 + 1] -= v;
    diff[x2 + 1][y2 + 1] += v;
}

for (int i = 1; i <= n; i++) {
    for (int j = 1; j <= m; j++) {
        diff[i][j] += diff[i - 1][j] + diff[i][j - 1]
                    - diff[i - 1][j - 1];
        cout << diff[i][j] << (j == m ? '\n' : ' ');
    }
}
\end{lstlisting}

\Warn{二维差分只适合修改后统一求结果。}如果修改和查询交替出现，要考虑二维树状数组或其他方法，B 题中通常会有更简单性质。

\section{哈希统计}

\subsection{识别信号}

题面出现：

\begin{itemize}
  \item 统计每个数或字符串出现次数。
  \item 判断重复。
  \item 找出现最多的元素。
  \item 判断两个集合是否有交集。
\end{itemize}

\subsection{值域小：计数数组}

\begin{lstlisting}
int n;
cin >> n;
vector<int> cnt(100001, 0);
for (int i = 0; i < n; i++) {
    int x;
    cin >> x;
    cnt[x]++;
}
\end{lstlisting}

\subsection{值域大：map}

\begin{lstlisting}
int n;
cin >> n;
map<long long, int> cnt;
for (int i = 0; i < n; i++) {
    long long x;
    cin >> x;
    cnt[x]++;
}

for (auto [x, c] : cnt) {
    cout << x << ' ' << c << '\n';
}
\end{lstlisting}

\subsection{字符串统计}

\begin{lstlisting}
int n;
cin >> n;
unordered_map<string, int> cnt;
for (int i = 0; i < n; i++) {
    string s;
    cin >> s;
    cnt[s]++;
}
\end{lstlisting}

\subsection{找出现次数最多}

\begin{lstlisting}
map<string, int> cnt;
// fill cnt

string best;
int bestCnt = -1;
for (auto [s, c] : cnt) {
    if (c > bestCnt) {
        bestCnt = c;
        best = s;
    }
}
cout << best << ' ' << bestCnt << '\n';
\end{lstlisting}

如果有并列规则，例如字典序最小，由 \Code{map} 遍历天然按字典序升序。

\section{余数和同余统计}

\subsection{识别信号}

题面出现：

\begin{itemize}
  \item 能否被 $k$ 整除。
  \item 两个数相加后是 $k$ 的倍数。
  \item 前缀和模 $k$ 相同。
  \item 按余数分类统计。
\end{itemize}

\subsection{统计每个余数}

\begin{lstlisting}
int n, k;
cin >> n >> k;
vector<long long> cnt(k, 0);

for (int i = 0; i < n; i++) {
    long long x;
    cin >> x;
    int r = (int)(x % k);
    if (r < 0) r += k;
    cnt[r]++;
}
\end{lstlisting}

\subsection{两数和能被 k 整除}

\begin{lstlisting}
long long ans = 0;
for (int r = 0; r < k; r++) {
    int need = (k - r) % k;
    if (r < need) {
        ans += cnt[r] * cnt[need];
    } else if (r == need) {
        ans += cnt[r] * (cnt[r] - 1) / 2;
    }
}
cout << ans << '\n';
\end{lstlisting}

\subsection{前缀和模 k}

如果两个前缀和模 $k$ 相同，它们之间的区间和能被 $k$ 整除。

\begin{lstlisting}
int n, k;
cin >> n >> k;
vector<long long> cnt(k, 0);
cnt[0] = 1;

long long sum = 0;
long long ans = 0;
for (int i = 0; i < n; i++) {
    long long x;
    cin >> x;
    sum += x;
    int r = (int)(sum % k);
    if (r < 0) r += k;
    ans += cnt[r];
    cnt[r]++;
}
cout << ans << '\n';
\end{lstlisting}

\Warn{余数题注意：}如果 $k$ 很大，不能开 \Code{vector cnt(k)}，改用 \Code{map} 或 \Code{unordered\_map}。

\section{排序加扫描}

\subsection{识别信号}

题面出现：

\begin{itemize}
  \item 排序后相邻元素有关系。
  \item 需要合并区间。
  \item 需要按时间顺序处理事件。
  \item 需要按优先级选择。
\end{itemize}

\subsection{排序后统计相同元素段}

\begin{lstlisting}
sort(a.begin(), a.end());
int n = (int)a.size();
for (int i = 0; i < n; ) {
    int j = i;
    while (j < n && a[j] == a[i]) {
        j++;
    }
    int count = j - i;
    cout << a[i] << ' ' << count << '\n';
    i = j;
}
\end{lstlisting}

\subsection{合并区间}

输入若干区间，合并重叠部分：

\begin{lstlisting}
vector<pair<int, int>> segs;
sort(segs.begin(), segs.end());

vector<pair<int, int>> merged;
for (auto [l, r] : segs) {
    if (merged.empty() || l > merged.back().second) {
        merged.push_back({l, r});
    } else {
        merged.back().second = max(merged.back().second, r);
    }
}
\end{lstlisting}

如果题目认为相邻区间也可合并，例如 $[1,2]$ 和 $[3,4]$，条件要改成 \Code{l > merged.back().second + 1}。

\subsection{事件排序}

当题目有“按时间发生”的操作，先把事件读入结构体，再排序处理。

\begin{lstlisting}
struct Event {
    int time;
    int type;
    int id;
};

vector<Event> events;

sort(events.begin(), events.end(), [](const Event& a, const Event& b) {
    if (a.time != b.time) return a.time < b.time;
    return a.type < b.type; // tie rule depends on statement
});

for (const Event& e : events) {
    // process event
}
\end{lstlisting}

事件排序检查：

\begin{itemize}
  \item 同一时间多个事件的处理顺序。
  \item 开始事件和结束事件谁先处理。
  \item 是否需要稳定排序。
  \item 时间和答案是否要 \Code{long long}。
\end{itemize}

\subsection{坐标离散化入口}

当坐标很大，但只关心出现过的坐标，可以离散化。

识别信号：

\begin{itemize}
  \item 坐标可达 $10^9$，不能直接开数组。
  \item 点的数量只有 $10^5$。
  \item 只需要比较大小、排序、相对位置。
\end{itemize}

模板：

\begin{lstlisting}
vector<int> xs;
for (int x : original) {
    xs.push_back(x);
}

sort(xs.begin(), xs.end());
xs.erase(unique(xs.begin(), xs.end()), xs.end());

auto getId = [&](int x) {
    return (int)(lower_bound(xs.begin(), xs.end(), x) - xs.begin()) + 1;
};

int id = getId(original[0]);
\end{lstlisting}

\Warn{区间离散化要小心。}如果题目处理连续区间长度，不能只保存端点编号，还要考虑相邻坐标之间的真实长度。

\section{堆和动态维护}

\subsection{priority\_queue 识别信号}

题面出现：

\begin{itemize}
  \item 每次取当前最大或最小。
  \item 动态加入元素。
  \item 模拟优先级队列。
  \item 求前 $k$ 大、前 $k$ 小。
\end{itemize}

\subsection{小根堆}

\begin{lstlisting}
priority_queue<int, vector<int>, greater<int>> pq;
pq.push(5);
pq.push(2);
cout << pq.top() << '\n'; // 2
pq.pop();
\end{lstlisting}

\subsection{保留最大的 k 个数}

用小根堆保存当前最大的 $k$ 个数。

\begin{lstlisting}
int k;
cin >> k;
priority_queue<int, vector<int>, greater<int>> pq;

for (int x : a) {
    pq.push(x);
    if ((int)pq.size() > k) {
        pq.pop();
    }
}

cout << pq.top() << '\n'; // kth largest
\end{lstlisting}

\subsection{set 动态维护}

\Code{set} 适合动态插入、删除、找前驱后继。

\begin{lstlisting}
set<int> s;
s.insert(10);
s.insert(20);
s.insert(30);

int x = 25;
auto it = s.lower_bound(x); // first >= x
if (it != s.end()) {
    cout << "next=" << *it << '\n';
}
if (it != s.begin()) {
    --it;
    cout << "prev=" << *it << '\n';
}
\end{lstlisting}

\Warn{迭代器小心：}移动迭代器前先判断是不是 \Code{begin()} 或 \Code{end()}。

\subsection{map 动态维护}

\Code{map} 适合动态统计并按 key 有序遍历。

\begin{lstlisting}
map<int, int> cnt;
cnt[x]++;

cnt[x]--;
if (cnt[x] == 0) {
    cnt.erase(x);
}
\end{lstlisting}

取最小 key 和最大 key：

\begin{lstlisting}
int mn = cnt.begin()->first;
int mx = cnt.rbegin()->first;
\end{lstlisting}

\Warn{取 \Code{begin/rbegin} 前确认 map 非空。}

\section{基础贪心}

\subsection{识别信号}

B 题中的贪心通常比较直接，常见信号：

\begin{itemize}
  \item 尽量多选一些区间、任务、物品。
  \item 每次选择当前最早结束、最小代价、最大收益。
  \item 排序后从前往后扫描。
  \item 题目问最少次数覆盖、最多不冲突任务。
\end{itemize}

\subsection{最多不重叠区间}

按右端点从小到大排序，每次选最早结束的区间。

\begin{lstlisting}
vector<pair<int, int>> segs; // l, r
sort(segs.begin(), segs.end(), [](auto a, auto b) {
    if (a.second != b.second) return a.second < b.second;
    return a.first < b.first;
});

int ans = 0;
int lastEnd = -1000000000;
for (auto [l, r] : segs) {
    if (l > lastEnd) { // use l >= lastEnd if endpoints can touch
        ans++;
        lastEnd = r;
    }
}
cout << ans << '\n';
\end{lstlisting}

\subsection{贪心易错点}

\begin{itemize}
  \item 先确认排序关键字，是按左端点、右端点、代价还是收益。
  \item 区间相邻是否算冲突：\Code{l > lastEnd} 还是 \Code{l >= lastEnd}。
  \item 贪心不是看起来合理就一定对。B 题若想不出证明，至少用小样例反驳自己。
\end{itemize}

\section{单调栈和单调队列入口}

\subsection{单调栈识别信号}

题面出现：

\begin{itemize}
  \item 找每个元素左边或右边第一个更大/更小的元素。
  \item 柱状图、温度、最近更高。
  \item 暴力向左/向右扫描会超时。
\end{itemize}

\subsection{右边第一个更大元素}

\begin{lstlisting}
int n;
cin >> n;
vector<int> a(n);
for (int i = 0; i < n; i++) cin >> a[i];

vector<int> ans(n, -1);
stack<int> st; // indices, a[st] decreasing

for (int i = 0; i < n; i++) {
    while (!st.empty() && a[i] > a[st.top()]) {
        ans[st.top()] = i;
        st.pop();
    }
    st.push(i);
}
\end{lstlisting}

\subsection{单调队列识别信号}

题面出现固定长度滑动窗口最大值/最小值。

\begin{lstlisting}
int n, k;
cin >> n >> k;
vector<int> a(n);
for (int i = 0; i < n; i++) cin >> a[i];

deque<int> dq; // indices, values decreasing
for (int i = 0; i < n; i++) {
    while (!dq.empty() && dq.front() <= i - k) dq.pop_front();
    while (!dq.empty() && a[dq.back()] <= a[i]) dq.pop_back();
    dq.push_back(i);

    if (i >= k - 1) {
        cout << a[dq.front()] << '\n'; // max of window
    }
}
\end{lstlisting}

\Warn{单调结构新手建议：}如果没练过，不要在考场第一次使用。可以作为 B/C/D 题冲分模板，但必须提前练。

\section{双指针入门}

\subsection{识别信号}

题面出现：

\begin{itemize}
  \item 数组已排序，找一对数。
  \item 连续区间满足某个条件。
  \item 左右端点移动，区间和或区间长度变化。
\end{itemize}

\subsection{排序后找两数和}

\begin{lstlisting}
sort(a.begin(), a.end());
int l = 0, r = n - 1;
bool ok = false;

while (l < r) {
    long long sum = 1LL * a[l] + a[r];
    if (sum == target) {
        ok = true;
        break;
    } else if (sum < target) {
        l++;
    } else {
        r--;
    }
}

cout << (ok ? "YES" : "NO") << '\n';
\end{lstlisting}

\subsection{正数数组的连续区间}

如果数组元素都是非负数，可以用滑动窗口。

\begin{lstlisting}
int l = 0;
long long sum = 0;
int best = 0;

for (int r = 0; r < n; r++) {
    sum += a[r];
    while (sum > limit && l <= r) {
        sum -= a[l];
        l++;
    }
    best = max(best, r - l + 1);
}
\end{lstlisting}

\Warn{滑动窗口前提：}通常要求元素非负或条件具有单调性。若有负数，窗口移动可能失效。

\section{二分与二分答案}

\subsection{普通二分查找}

数组必须有序。

\begin{lstlisting}
sort(a.begin(), a.end());
auto it = lower_bound(a.begin(), a.end(), x);
if (it != a.end()) {
    cout << *it << '\n';
}
\end{lstlisting}

\subsection{二分答案识别}

题面出现：

\begin{itemize}
  \item 最大化最小值。
  \item 最小化最大值。
  \item 求最小的可行时间、容量、代价。
  \item 给定一个答案 $mid$，容易判断是否可行。
\end{itemize}

\subsection{二分答案模板}

找最小可行值：

\begin{lstlisting}
bool check(long long mid) {
    // return true if mid is feasible
    return true;
}

long long l = 0, r = 1000000000000000000LL;
while (l < r) {
    long long mid = l + (r - l) / 2;
    if (check(mid)) {
        r = mid;
    } else {
        l = mid + 1;
    }
}
cout << l << '\n';
\end{lstlisting}

找最大可行值：

\begin{lstlisting}
long long l = 0, r = 1000000000000000000LL;
while (l < r) {
    long long mid = l + (r - l + 1) / 2;
    if (check(mid)) {
        l = mid;
    } else {
        r = mid - 1;
    }
}
cout << l << '\n';
\end{lstlisting}

\Warn{二分答案关键：}不是所有题都能二分。必须满足可行性具有单调性。

\section{BFS 入门}

\subsection{识别信号}

题面出现：

\begin{itemize}
  \item 最少步数。
  \item 迷宫、网格移动。
  \item 每次操作代价相同。
  \item 从起点扩展到终点。
\end{itemize}

\subsection{网格 BFS 模板}

\begin{lstlisting}
int n, m;
cin >> n >> m;
vector<string> g(n);
for (int i = 0; i < n; i++) cin >> g[i];

vector<vector<int>> dist(n, vector<int>(m, -1));
queue<pair<int, int>> q;

int sx, sy, tx, ty;
cin >> sx >> sy >> tx >> ty;

dist[sx][sy] = 0;
q.push({sx, sy});

int dx[4] = {-1, 0, 1, 0};
int dy[4] = {0, 1, 0, -1};

while (!q.empty()) {
    auto [x, y] = q.front();
    q.pop();

    for (int d = 0; d < 4; d++) {
        int nx = x + dx[d];
        int ny = y + dy[d];
        if (nx < 0 || nx >= n || ny < 0 || ny >= m) continue;
        if (g[nx][ny] == '#') continue;
        if (dist[nx][ny] != -1) continue;

        dist[nx][ny] = dist[x][y] + 1;
        q.push({nx, ny});
    }
}

cout << dist[tx][ty] << '\n';
\end{lstlisting}

\subsection{BFS 易错点}

\begin{itemize}
  \item 坐标输入是 0 下标还是 1 下标。
  \item 障碍字符是什么。
  \item 起点是否可能等于终点。
  \item 应在入队时标记访问，避免重复入队。
  \item BFS 适合每步代价相同；边权不同要考虑 Dijkstra。
\end{itemize}

\subsection{多源 BFS}

如果有多个起点，且要求每个点到最近起点的距离，把所有起点同时入队。

\begin{lstlisting}
vector<vector<int>> dist(n, vector<int>(m, -1));
queue<pair<int, int>> q;

for (auto [x, y] : starts) {
    dist[x][y] = 0;
    q.push({x, y});
}

while (!q.empty()) {
    auto [x, y] = q.front();
    q.pop();
    for (int d = 0; d < 4; d++) {
        int nx = x + dx[d];
        int ny = y + dy[d];
        if (nx < 0 || nx >= n || ny < 0 || ny >= m) continue;
        if (dist[nx][ny] != -1) continue;
        dist[nx][ny] = dist[x][y] + 1;
        q.push({nx, ny});
    }
}
\end{lstlisting}

识别信号：

\begin{itemize}
  \item 多个火源、多个出口、多个起点。
  \item 求离最近特殊点的距离。
  \item 每一步代价相同。
\end{itemize}

\section{DFS 和连通块}

\subsection{识别信号}

题面出现：

\begin{itemize}
  \item 统计连通块数量。
  \item 找一片区域大小。
  \item 网格中相邻的同类格子。
  \item 图中从某点能到哪些点。
\end{itemize}

\subsection{网格 DFS}

\begin{lstlisting}
int n, m;
vector<string> g;
vector<vector<int>> vis;
int dx[4] = {-1, 0, 1, 0};
int dy[4] = {0, 1, 0, -1};

int dfs(int x, int y) {
    vis[x][y] = 1;
    int area = 1;
    for (int d = 0; d < 4; d++) {
        int nx = x + dx[d];
        int ny = y + dy[d];
        if (nx < 0 || nx >= n || ny < 0 || ny >= m) continue;
        if (vis[nx][ny]) continue;
        if (g[nx][ny] == '#') continue;
        area += dfs(nx, ny);
    }
    return area;
}
\end{lstlisting}

统计连通块：

\begin{lstlisting}
int components = 0;
for (int i = 0; i < n; i++) {
    for (int j = 0; j < m; j++) {
        if (!vis[i][j] && g[i][j] != '#') {
            components++;
            int area = dfs(i, j);
        }
    }
}
\end{lstlisting}

\Warn{递归深度风险：}如果网格很大，递归 DFS 可能爆栈。B 题中可改用 BFS 更稳。

\section{Floyd 入门}

\subsection{识别信号}

\begin{itemize}
  \item 问任意两点之间最短路。
  \item 点数较小，例如 $n \le 300$ 或 $n \le 500$。
  \item 图可能是稠密图。
\end{itemize}

\subsection{Floyd 模板}

\begin{lstlisting}
const long long INF = (long long)4e18;

int n, m;
cin >> n >> m;
vector<vector<long long>> dist(n + 1, vector<long long>(n + 1, INF));

for (int i = 1; i <= n; i++) {
    dist[i][i] = 0;
}

for (int i = 0; i < m; i++) {
    int u, v;
    long long w;
    cin >> u >> v >> w;
    dist[u][v] = min(dist[u][v], w);
    dist[v][u] = min(dist[v][u], w); // remove if directed
}

for (int k = 1; k <= n; k++) {
    for (int i = 1; i <= n; i++) {
        for (int j = 1; j <= n; j++) {
            if (dist[i][k] == INF || dist[k][j] == INF) continue;
            dist[i][j] = min(dist[i][j], dist[i][k] + dist[k][j]);
        }
    }
}
\end{lstlisting}

\subsection{Floyd 易错点}

\begin{itemize}
  \item 三层循环顺序必须是 \Code{k,i,j}。
  \item 无向图要双向赋值，有向图不要。
  \item 多条边取最小边权。
  \item 点数大时不能用 Floyd。
\end{itemize}

\section{Dijkstra 入门}

\subsection{识别信号}

\begin{itemize}
  \item 正权图最短路。
  \item 边权不是全 1。
  \item 点和边较多，不能 Floyd。
  \item 单源最短路：从一个起点到其他点。
\end{itemize}

\subsection{Dijkstra 模板}

\begin{lstlisting}
struct Edge {
    int to;
    long long w;
};

const long long INF = (long long)4e18;

int n, m, s;
cin >> n >> m >> s;
vector<vector<Edge>> g(n + 1);

for (int i = 0; i < m; i++) {
    int u, v;
    long long w;
    cin >> u >> v >> w;
    g[u].push_back({v, w});
    g[v].push_back({u, w}); // remove if directed
}

vector<long long> dist(n + 1, INF);
priority_queue<pair<long long, int>,
               vector<pair<long long, int>>,
               greater<pair<long long, int>>> pq;

dist[s] = 0;
pq.push({0, s});

while (!pq.empty()) {
    auto [d, u] = pq.top();
    pq.pop();
    if (d != dist[u]) continue;

    for (auto e : g[u]) {
        if (dist[e.to] > dist[u] + e.w) {
            dist[e.to] = dist[u] + e.w;
            pq.push({dist[e.to], e.to});
        }
    }
}
\end{lstlisting}

\subsection{Dijkstra 易错点}

\begin{itemize}
  \item 只适用于非负边权。存在负边权时不能直接用。
  \item 距离用 \Code{long long}。
  \item 无向图双向加边，有向图单向加边。
  \item 堆中可能有过期状态，要跳过。
  \item 如果所有边权都是 1，用 BFS 更简单。
\end{itemize}

\section{并查集入门}

\subsection{识别信号}

题面出现：

\begin{itemize}
  \item 合并两个集合。
  \item 判断两个元素是否属于同一组。
  \item 连通块数量。
  \item 朋友关系、同类关系、网络连通。
\end{itemize}

\subsection{并查集模板}

\begin{lstlisting}
struct DSU {
    vector<int> parent, sz;

    DSU(int n) {
        parent.resize(n + 1);
        sz.assign(n + 1, 1);
        for (int i = 1; i <= n; i++) parent[i] = i;
    }

    int find(int x) {
        if (parent[x] == x) return x;
        return parent[x] = find(parent[x]);
    }

    bool unite(int a, int b) {
        int ra = find(a);
        int rb = find(b);
        if (ra == rb) return false;
        if (sz[ra] < sz[rb]) swap(ra, rb);
        parent[rb] = ra;
        sz[ra] += sz[rb];
        return true;
    }

    bool same(int a, int b) {
        return find(a) == find(b);
    }
};
\end{lstlisting}

\section{拓扑排序入口}

\subsection{识别信号}

题面出现：

\begin{itemize}
  \item 任务必须按先后顺序完成。
  \item A 必须在 B 前面。
  \item 课程依赖、工程依赖。
  \item 判断是否存在合法顺序。
\end{itemize}

\subsection{拓扑排序模板}

\begin{lstlisting}
int n, m;
cin >> n >> m;
vector<vector<int>> g(n + 1);
vector<int> indeg(n + 1, 0);

for (int i = 0; i < m; i++) {
    int u, v;
    cin >> u >> v; // u before v
    g[u].push_back(v);
    indeg[v]++;
}

queue<int> q;
for (int i = 1; i <= n; i++) {
    if (indeg[i] == 0) q.push(i);
}

vector<int> order;
while (!q.empty()) {
    int u = q.front();
    q.pop();
    order.push_back(u);

    for (int v : g[u]) {
        indeg[v]--;
        if (indeg[v] == 0) q.push(v);
    }
}

if ((int)order.size() == n) {
    cout << "YES\n";
} else {
    cout << "NO\n"; // has cycle
}
\end{lstlisting}

\subsection{拓扑排序易错点}

\begin{itemize}
  \item 边的方向：如果 A 在 B 前，应连 A 到 B。
  \item 入度为 0 的点一开始全部入队。
  \item 输出顺序是否要求字典序最小。若要求，队列换成小根堆。
  \item 若排序结果少于 $n$ 个点，说明有环。
\end{itemize}

\section{简单 DP 入口}

\subsection{识别信号}

题面出现：

\begin{itemize}
  \item 求最优值：最大、最小、方案数。
  \item 当前结果依赖前一个或前几个位置。
  \item 选择或不选择某个物品。
  \item 容量、价值、花费、恰好装满。
\end{itemize}

\subsection{线性 DP 示例}

例如每一步只能从前一步或前两步转移：

\begin{lstlisting}
int n;
cin >> n;
vector<long long> dp(n + 1, 0);
dp[0] = 1;
dp[1] = 1;
for (int i = 2; i <= n; i++) {
    dp[i] = dp[i - 1] + dp[i - 2];
}
cout << dp[n] << '\n';
\end{lstlisting}

\subsection{01 背包入口}

每个物品只能选一次：

\begin{lstlisting}
int n, V;
cin >> n >> V;
vector<int> w(n + 1), val(n + 1);
for (int i = 1; i <= n; i++) {
    cin >> w[i] >> val[i];
}

vector<int> dp(V + 1, 0);
for (int i = 1; i <= n; i++) {
    for (int j = V; j >= w[i]; j--) {
        dp[j] = max(dp[j], dp[j - w[i]] + val[i]);
    }
}
cout << dp[V] << '\n';
\end{lstlisting}

\subsection{完全背包入口}

每个物品可以选无限次，容量正序。

\begin{lstlisting}
int n, V;
cin >> n >> V;
vector<int> w(n + 1), val(n + 1);
for (int i = 1; i <= n; i++) {
    cin >> w[i] >> val[i];
}

vector<int> dp(V + 1, 0);
for (int i = 1; i <= n; i++) {
    for (int j = w[i]; j <= V; j++) {
        dp[j] = max(dp[j], dp[j - w[i]] + val[i]);
    }
}
cout << dp[V] << '\n';
\end{lstlisting}

\subsection{DP 易错点}

\begin{itemize}
  \item 先定义 \Code{dp[i]} 表示什么。
  \item 写出初始状态。
  \item 写出转移方程。
  \item 01 背包容量倒序，完全背包容量正序。
  \item 方案数可能很大，要看是否取模。
\end{itemize}

\section{基础数学入口}

\subsection{gcd 和 lcm}

\begin{lstlisting}
long long g = gcd(a, b);
long long l = a / g * b;
\end{lstlisting}

\Warn{先除后乘，降低溢出风险。}

\subsection{快速幂}

如果题目要求 $a^b \bmod mod$：

\begin{lstlisting}
long long modPow(long long a, long long b, long long mod) {
    long long res = 1 % mod;
    a %= mod;
    while (b > 0) {
        if (b & 1) res = res * a % mod;
        a = a * a % mod;
        b >>= 1;
    }
    return res;
}
\end{lstlisting}

\subsection{判断素数}

单个数判断：

\begin{lstlisting}
bool isPrime(long long x) {
    if (x < 2) return false;
    for (long long d = 2; d * d <= x; d++) {
        if (x % d == 0) return false;
    }
    return true;
}
\end{lstlisting}

\subsection{枚举约数}

\begin{lstlisting}
vector<long long> divisors;
for (long long d = 1; d * d <= x; d++) {
    if (x % d == 0) {
        divisors.push_back(d);
        if (d != x / d) divisors.push_back(x / d);
    }
}
sort(divisors.begin(), divisors.end());
\end{lstlisting}

\Warn{数学题先看范围。}如果要对很多个数判断素数，可能需要筛法；如果只是少量数，试除通常够用。

使用：

\begin{lstlisting}
int n, m;
cin >> n >> m;
DSU dsu(n);

while (m--) {
    int op, a, b;
    cin >> op >> a >> b;
    if (op == 1) dsu.unite(a, b);
    else cout << (dsu.same(a, b) ? "YES" : "NO") << '\n';
}
\end{lstlisting}

\section{B 题提交前检查}

\begin{itemize}
  \item 是否估算过复杂度。
  \item 是否把暴力写到了大数据上。
  \item 前缀和、差分的下标是否统一。
  \item 求和、答案、距离是否用 \Code{long long}。
  \item 排序比较器是否正确。
  \item \Code{lower\_bound} 前是否排序。
  \item BFS 是否在入队时标记。
  \item 并查集是否初始化所有点。
  \item 多组数据是否清空容器。
\end{itemize}

\section{B 题冲分判断表}

当你已经有一个暴力思路，但担心超时时，用这张表判断下一步。

\begin{longtable}{p{4cm}p{5cm}p{5cm}}
\toprule
当前情况 & 优先尝试 & 原因 \\
\midrule
暴力是两层循环找一对数 & 排序、哈希、双指针 & 二元关系常能降到 $O(n\log n)$ 或 $O(n)$。 \\
暴力是每次求区间和 & 前缀和 & 静态区间查询标准做法。 \\
暴力是每次区间加 & 差分 & 最后统一查询时最简单。 \\
暴力是每次找最大/最小 & 堆或 set & 动态维护极值。 \\
坐标太大但点数不多 & 离散化 & 把大坐标压缩成小编号。 \\
规则按时间发生 & 事件排序 & 统一按时间顺序处理。 \\
左边右边都要快速知道信息 & 前后缀预处理 & 删除一个元素或分割数组常见。 \\
涉及整除、余数相同 & 余数统计 & 同余关系常能转成计数。 \\
最多不冲突任务或区间 & 贪心排序 & 常见按右端点排序。 \\
每个点找旁边第一个更大/更小 & 单调栈 & 避免每次向旁边扫。 \\
滑动窗口最大/最小 & 单调队列 & 固定窗口极值。 \\
网格最短步数 & BFS & 每步代价相同。 \\
网格连通区域 & DFS/BFS & 统计区域或连通块。 \\
图边权为 1 & BFS & 比 Dijkstra 简单。 \\
图边权为正 & Dijkstra & 单源正权最短路。 \\
点数小且任意两点 & Floyd & 写法短，适合小图。 \\
合并和查询同组 & 并查集 & 连通性标准做法。 \\
依赖顺序 & 拓扑排序 & 入度为 0 的点先处理。 \\
\bottomrule
\end{longtable}

\section{B 题常见边界库}

\begin{longtable}{p{4cm}p{5cm}p{5cm}}
\toprule
题型 & 边界样例 & 检查目标 \\
\midrule
前缀和 & $l=1$，$r=n$，$l=r$ & 端点公式。 \\
二维前缀和 & 单行、单列、整个矩阵、单个格子 & 四项公式。 \\
差分 & 修改到 $r=n$，只改一个点 & \Code{r+1} 和数组大小。 \\
排序 & 全相等、逆序、相同主关键字 & 比较器。 \\
哈希统计 & 没有重复、全部重复、key 很大 & 容器选择。 \\
双指针 & 无解、两个数刚好在两端、有重复 & 指针移动。 \\
二分答案 & 最小值可行、最大值不可行、边界答案 & \Code{check} 单调性。 \\
BFS & 起点等于终点、无路、全障碍附近 & 距离初始化。 \\
DFS 连通块 & 全障碍、全可走、单个格子 & 访问标记。 \\
并查集 & 自己和自己合并、重复合并 & find/unite。 \\
Dijkstra & 不连通、多条边、起点到自己 & INF 和取最小边。 \\
余数统计 & 负数、$k=1$、所有余数相同 & 取模修正和计数公式。 \\
贪心 & 区间端点相同、相邻是否冲突 & 排序关键字和边界条件。 \\
单调栈 & 全递增、全递减、全相等 & 比较符号是 \Code{<} 还是 \Code{<=}。 \\
\bottomrule
\end{longtable}

\section{B 题本地调试流程}

\subsection{先用暴力对拍思路}

如果你写了优化算法，但不确定是否正确，可以在本地用小数据和暴力对比。考场时间紧时不一定完整写对拍，但可以用这个思想手工造样例。

\begin{itemize}
  \item $n$ 取 3 到 6。
  \item 手算所有操作结果。
  \item 用暴力程序或草稿验证优化结果。
  \item 专门造边界：重复、空、单点、全相等。
\end{itemize}

\subsection{复杂度复查}

提交前问：

\begin{itemize}
  \item 有没有隐藏的双重循环。
  \item \Code{map} 操作是否在循环里被调用很多次，是否仍可接受。
  \item 排序是否放在循环内部。若是，通常危险。
  \item BFS/Dijkstra 是否会重复跑很多次。
  \item 字符串拼接是否在大循环里反复发生。
\end{itemize}

\subsection{状态清空}

多组数据时，尤其注意：

\begin{lstlisting}
g.assign(n + 1, vector<int>());
dist.assign(n + 1, INF);
vis.assign(n + 1, 0);
cnt.clear();
while (!q.empty()) q.pop();
\end{lstlisting}

\Warn{队列、优先队列没有 clear。}最简单的方法是在每组数据内部重新定义它。

\section{B 题卡住时的部分分策略}

\begin{enumerate}
  \item 如果不会优化，先写 $O(n^2)$ 暴力，确认题意。
  \item 如果有多档数据，先保证小数据正确。
  \item 如果是区间题，先写单次查询或无修改版本。
  \item 如果是图题，先写 BFS/DFS 处理无权或小图。
  \item 如果是复杂排序规则，先保证主关键字正确，再处理并列规则。
\end{enumerate}

\section{B 题手写模板清单}

这些模板建议在考前单独练习到“看着能快速写对”。

\begin{longtable}{p{4cm}p{4cm}p{6cm}}
\toprule
模板 & 熟练要求 & 主要用途 \\
\midrule
一维前缀和 & 必须会默写 & 静态区间和。 \\
二维前缀和 & 照着能写 & 子矩阵和。 \\
一维差分 & 必须会默写 & 区间加，最后查询。 \\
二维差分 & 照着能写 & 矩形加，最后输出。 \\
map/unordered\_map 统计 & 必须会默写 & 频率统计。 \\
排序加结构体比较器 & 必须会默写 & 排名、事件、贪心。 \\
双指针 & 照着能写 & 排序后找对、滑动窗口。 \\
二分答案 & 照着能写 & 最大化最小值、最小化最大值。 \\
前后缀预处理 & 照着能写 & 删除一个元素、左右信息。 \\
余数统计 & 照着能写 & 整除、同余、区间和取模。 \\
基础贪心 & 照着能写 & 区间选择、任务安排。 \\
单调栈/队列 & 了解并练过再用 & 最近更大、滑动窗口极值。 \\
网格 BFS & 必须照着能写 & 最短步数。 \\
并查集 & 照着能写 & 连通性。 \\
拓扑排序 & 照着能写 & 依赖顺序。 \\
Dijkstra & 照着能写 & 正权最短路。 \\
\bottomrule
\end{longtable}

\chapter{C 题决策树与大模拟方法}\label{ch:07}

\section{C 题定位}\label{sec:07-1}

C 题的长题面常见多个对象、几种命令和一些特殊规则。读题和组织代码会占不少时间，先把这些内容分清楚，后面才容易查错。

\begin{itemize}
  \item 建议目标时间：60--100 分钟。
  \item 目标：能满分则满分；若细节太多，先完成题目明确给出的子任务。
  \item 高频主题：大模拟、字符串解析、状态维护、网格搜索、图论基础、数据结构基础、简单 DP。
  \item 容易出错的地方：漏规则、状态没记全、不同命令混在同一段代码里。
\end{itemize}

可以先列一张对象和命令表，确认每种操作会改哪些字段。

\section{C 题开题流程}\label{sec:07-2}

\begin{enumerate}
  \item 快速浏览题面，判断是大模拟还是算法题。
  \item 圈出所有对象：人、文件、任务、节点、格子、命令、事件。
  \item 圈出所有状态：位置、时间、分数、开关、权限、颜色、方向。
  \item 圈出所有操作：新增、删除、查询、移动、更新、比较。
  \item 圈出优先级和并列规则。
  \item 设计结构体和函数。
  \item 先实现最基础规则，通过样例的一部分。
  \item 再逐条补特殊规则。
\end{enumerate}

\section{C 题读题三遍法}\label{sec:07-3}

题面较长时可以分几遍读：先看要模拟什么，再记状态和命令，最后核对特殊规则。

\subsection{第一遍：只判断题型}

第一遍快速读，不抠细节，只回答：

\begin{itemize}
  \item 是大模拟，还是算法题？
  \item 输入是一次性数据，还是很多命令？
  \item 输出是每次查询输出，还是最后统一输出？
  \item 数据范围是否明显要求优化？
\end{itemize}

第一遍还没理清时，可以结合样例看看一次操作前后发生了什么。

\subsection{第二遍：圈对象、状态、操作}

第二遍要动笔。把题目中的名词和动词分开。

\begin{longtable}{p{3cm}p{5cm}p{6cm}}
\toprule
类别 & 怎么找 & 例子 \\
\midrule
\endfirsthead
\toprule
类别 & 怎么找 & 例子 \\
\midrule
\endhead
对象 & 被操作的名词 & 用户、文件、进程、任务、车辆、节点、格子。 \\
属性 & 对象身上的数据 & 编号、名字、位置、时间、状态、得分。 \\
操作 & 题目中的动词或命令 & 新增、删除、移动、查询、合并、修改。 \\
规则 & 操作发生时的限制 & 如果不存在则忽略；如果冲突按优先级。 \\
输出 & 什么时候打印 & 每次 query、最后统计、发生错误时。 \\
\bottomrule
\end{longtable}

\subsection{第三遍：找坑点}

第三遍只找容易错的细节：

\begin{itemize}
  \item 相同时间、相同分数、相同优先级怎么处理。
  \item 编号从 0 还是 1 开始。
  \item 找不到对象时怎么处理。
  \item 重复添加时怎么处理。
  \item 删除后是否还能查询。
  \item 最后一条命令或最后一个事件是否要特殊处理。
  \item 输出顺序按输入顺序、编号顺序、字典序还是时间顺序。
\end{itemize}

\section{从题面到代码的转换}\label{sec:07-4}

\subsection{对象变成 struct}

题目中的每类对象，优先考虑一个结构体。

\begin{lstlisting}
struct Task {
    int id;
    int startTime;
    int endTime;
    int priority;
    int status;
};
\end{lstlisting}

\subsection{编号查找用 map 或 vector}

如果编号是 $1$ 到 $n$：

\begin{lstlisting}
vector<Task> tasks(n + 1);
\end{lstlisting}

如果编号很大或是字符串：

\begin{lstlisting}
map<string, Task> tasks;
unordered_map<int, Task> tasksById;
\end{lstlisting}

\subsection{命令变成函数}

每个命令单独写函数：

\begin{lstlisting}
void addTask(int id, int priority) {
    // add or update
}

void deleteTask(int id) {
    // delete if exists
}

void queryTask(int id) {
    // print answer
}
\end{lstlisting}

\subsection{规则变成 if}

题目中的条件要落实到代码里。例如，查询前先判断对象是否存在，再判断它的状态：

\begin{lstlisting}
if (!tasks.count(id)) {
    cout << "Not Found\n";
    return;
}

if (tasks[id].status == FINISHED) {
    cout << "Finished\n";
    return;
}
\end{lstlisting}

把几条不同规则分开写，样例不对时比较容易找到是哪一步。

\section{状态维护和容器选择}\label{sec:07-5}

\subsection{容器选择表}

\begin{longtable}{p{4cm}p{5cm}p{5cm}}
\toprule
需求 & 推荐容器 & 原因 \\
\midrule
\endfirsthead
\toprule
需求 & 推荐容器 & 原因 \\
\midrule
\endhead
编号是 $1$ 到 $n$ & \Code{vector<Obj>} & 直接下标访问，最快。 \\
编号很大但唯一 & \Code{unordered\_map<int,Obj>} & 平均快速查找。 \\
编号是字符串 & \Code{map<string,Obj>} 或 \Code{unordered\_map} & 字符串到对象映射。 \\
需要按 key 有序输出 & \Code{map} & 自动按 key 排序。 \\
需要判断是否存在 & \Code{set} / \Code{unordered\_set} & 查重。 \\
需要动态最小/最大 & \Code{set} / \Code{priority\_queue} & 维护极值。 \\
需要前驱后继 & \Code{set} & \Code{lower\_bound}。 \\
需要按进入顺序处理 & \Code{queue} & FIFO。 \\
需要最近使用顺序 & \Code{list + map} & LRU。 \\
需要父子层级 & \Code{vector<Node>} + child & 树结构。 \\
\bottomrule
\end{longtable}

\subsection{操作复杂度估算}

\begin{longtable}{p{4cm}p{4cm}p{6cm}}
\toprule
容器 & 常见复杂度 & 注意 \\
\midrule
\endfirsthead
\toprule
容器 & 常见复杂度 & 注意 \\
\midrule
\endhead
\Code{vector} 下标访问 & $O(1)$ & 中间插入删除慢。 \\
\Code{map} 查找插入 & $O(\log n)$ & 有序但常数较大。 \\
\Code{unordered\_map} 查找插入 & 平均 $O(1)$ & 无序，最坏情况不稳定。 \\
\Code{set} 查找插入删除 & $O(\log n)$ & 可找前驱后继。 \\
\Code{priority\_queue} 取最优 & $O(\log n)$ 插入删除堆顶 & 不能直接删除中间元素。 \\
\Code{queue} 进出队 & $O(1)$ & 只能访问队首。 \\
\Code{list} 插入删除 & 已有迭代器时 $O(1)$ & 不能随机访问。 \\
\bottomrule
\end{longtable}

\subsection{C 题容器选择错误信号}

\begin{itemize}
  \item 每次命令都全表扫描，且命令数很大：应该用 map/set/hash。
  \item 每次取最大值都排序：应该用 set 或堆。
  \item 每次删除 vector 中间元素：可能需要标记删除、list 或 set。
  \item 需要有序输出却用了 unordered\_map：输出顺序会错。
  \item priority\_queue 里元素更新了但旧值还在：需要懒删除。
\end{itemize}

\section{C 题总决策树}\label{sec:07-6}

\begin{longtable}{p{4cm}p{5cm}p{5cm}}
\toprule
题面信号 & 优先怀疑 & 首选方法 \\
\midrule
\endfirsthead
\toprule
题面信号 & 优先怀疑 & 首选方法 \\
\midrule
\endhead
题面很长，规则很多 & 大模拟 & 结构体 + 函数拆解 \\
输入是一串命令 & 命令解析模拟 & stringstream / split / 状态表 \\
按时间发生事件 & 事件模拟 & 事件排序 / 优先队列 \\
对象有多种状态 & 状态机模拟 & enum/int 表示状态 \\
文件路径、配置、表达式 & 字符串解析 & 逐字符扫描 / 栈 \\
网格、地图、移动 & 网格模拟或 BFS/DFS & 方向数组、vis、dist \\
图、连通、最短 & 图论 & BFS / Dijkstra / 并查集 \\
区间修改查询 & 数据结构 & 前缀和 / 差分 / 树状数组 / 线段树 \\
求最优值且有阶段 & DP & 先定义状态 \\
\bottomrule
\end{longtable}

\section{C 题常见题型}\label{sec:07-7}

这张表用于快速判断 C 题该从哪个方向切入。

\begin{longtable}{p{4cm}p{5cm}p{5cm}}
\toprule
题目长相 & 本质 & 先翻哪里 \\
\midrule
\endfirsthead
\toprule
题目长相 & 本质 & 先翻哪里 \\
\midrule
\endhead
很多命令，每条命令改变系统状态 & 命令模拟 & 命令解析、状态设计、函数化 \\
按时间发生、开始结束、排队 & 事件模拟 & 事件排序、优先级、队列模拟 \\
字符串里有路径、参数、配置 & 字符串解析 & split、key=value、路径规范化 \\
输入像 URL、路由、参数 & 路由匹配 & URL 路由和参数解析 \\
输入像 JSON、配置、嵌套括号 & 结构解析 & 递归解析、栈、逐字符扫描 \\
输入像 Markdown、HTML、轻量标记 & 标记语言模拟 & 行扫描、状态切换 \\
有用户、角色、权限、规则 & 规则匹配 & 权限和规则匹配 \\
有表格、矩阵、坐标变换 & 表格模拟 & 矩阵操作、行列映射 \\
有缓存、队列、窗口、淘汰 & 数据结构模拟 & 队列、set、map、懒删除 \\
有目录、树、层级关系 & 树结构模拟 & 目录树、父子关系、DFS \\
有地图并带钥匙/状态 & 状态搜索 & 带状态 BFS \\
\bottomrule
\end{longtable}

\section{C 题细分决策树}\label{sec:07-8}

\subsection{第一问：这是大模拟还是算法题}

\begin{longtable}{p{5cm}p{5cm}p{4cm}}
\toprule
题面特征 & 判断 & 下一步 \\
\midrule
\endfirsthead
\toprule
题面特征 & 判断 & 下一步 \\
\midrule
\endhead
规则很多、命令很多、对象很多 & 大模拟 & 拆对象、状态、命令 \\
题面不长但数据范围大 & 算法题 & 回到 B/D 算法决策树 \\
有地图和最短步数 & 搜索题 & BFS / 状态 BFS \\
有依赖关系 & 图论题 & 拓扑 / 最短路 / 并查集 \\
有复杂字符串格式 & 解析题 & 字符串解析决策树 \\
有大量修改查询 & 数据结构题 & 前缀/差分/树状数组/线段树 \\
\bottomrule
\end{longtable}

\subsection{大模拟细分树}

\begin{longtable}{p{5cm}p{5cm}p{4cm}}
\toprule
题面信号 & 类型 & 首选结构 \\
\midrule
\endfirsthead
\toprule
题面信号 & 类型 & 首选结构 \\
\midrule
\endhead
每行一个操作名 & 命令系统 & \Code{handleCommand} 分发 \\
对象有生命周期 & 增删改查模拟 & \Code{map<int,Obj>} 或 \Code{vector<Obj>} \\
对象有多个状态 & 状态机 & 状态常量 + 转移函数 \\
按时间发生 & 事件模拟 & \Code{Event} + 排序 \\
每次取最优对象 & 优先级模拟 & 堆 / set / 懒删除堆 \\
行列、表格、棋盘 & 表格模拟 & 二维数组 / 行列映射 \\
路径、目录、父子层级 & 树结构模拟 & 节点数组 + child map \\
规则匹配、权限判断 & 规则系统 & \Code{Rule} + \Code{match} \\
\bottomrule
\end{longtable}

\subsection{字符串解析细分树}

\begin{longtable}{p{5cm}p{5cm}p{4cm}}
\toprule
字符串形态 & 优先方法 & 小心 \\
\midrule
\endfirsthead
\toprule
字符串形态 & 优先方法 & 小心 \\
\midrule
\endhead
按空格分隔 & \Code{stringstream} & 前面是否要 \Code{ignore} \\
按固定字符分隔 & \Code{split} & 连续分隔符 \\
\Code{key=value} & \Code{find + substr} & 值是否含空格 \\
路径 \Code{/a/b/c} & 路径 split / normalize & 根目录、\Code{..} \\
括号匹配 & 栈 & 空栈和最后剩余 \\
表达式求值 & 扫描 / 递归解析 & 优先级、括号 \\
URL 参数 & 先按 \Code{?}，再按 \Code{\&} & 参数为空 \\
嵌套 JSON 类 & 逐字符扫描 / 栈 & 先利用题目限制 \\
\bottomrule
\end{longtable}

\subsection{图网格细分树}

\begin{longtable}{p{5cm}p{5cm}p{4cm}}
\toprule
题面信号 & 类型 & 首选模板 \\
\midrule
\endfirsthead
\toprule
题面信号 & 类型 & 首选模板 \\
\midrule
\endhead
最少步数、每步代价相同 & 无权最短路 & BFS \\
多个起点最近距离 & 多源 BFS & 所有起点入队 \\
有钥匙、方向、模式等额外状态 & 状态 BFS & \Code{dist[x][y][state]} \\
统计区域大小 & 连通块 & DFS / BFS \\
正权边最短路 & 正权图 & Dijkstra \\
只问是否连通 & 连通性 & 并查集 / DFS \\
任务依赖顺序 & DAG & 拓扑排序 \\
\bottomrule
\end{longtable}

\section{大模拟总方法}\label{sec:07-9}

\subsection{大模拟什么题会用到}

\begin{itemize}
  \item 题面长度明显长。
  \item 有很多“如果……否则……”规则。
  \item 输入包含多条命令。
  \item 有多个对象，每个对象有状态。
  \item 输出要求根据最终状态或每次查询得到。
\end{itemize}

\subsection{大模拟代码骨架}

\begin{lstlisting}
#include <bits/stdc++.h>
using namespace std;
using ll = long long;

struct Item {
    int id;
    string name;
    int state;
};

int n, m;
vector<Item> items;

void readInput() {
    cin >> n >> m;
    items.resize(n);
    for (int i = 0; i < n; i++) {
        cin >> items[i].id >> items[i].name;
        items[i].state = 0;
    }
}

void handleCommand(const string& cmd) {
    if (cmd == "ADD") {
        // handle add
    } else if (cmd == "DEL") {
        // handle delete
    } else if (cmd == "QUERY") {
        // handle query
    }
}

void solve() {
    for (int i = 0; i < m; i++) {
        string cmd;
        cin >> cmd;
        handleCommand(cmd);
    }
}

int main() {
    ios::sync_with_stdio(false);
    cin.tie(nullptr);

    readInput();
    solve();
    return 0;
}
\end{lstlisting}

\subsection{为什么要拆函数}

\begin{itemize}
  \item 读入、处理、输出分开，查错更容易。
  \item 每个命令一个函数，避免主函数过长。
  \item 特殊规则可以单独补，不容易破坏已有逻辑。
  \item 样例错时能定位是哪条规则错。
\end{itemize}

\section{长题面拆解表}\label{sec:07-10}

读 C 题时，在草稿纸上画这个表：

\begin{longtable}{p{3cm}p{5cm}p{6cm}}
\toprule
项目 & 要记录什么 & 示例 \\
\midrule
\endfirsthead
\toprule
项目 & 要记录什么 & 示例 \\
\midrule
\endhead
对象 & 题目中被操作的东西 & 用户、文件、任务、格子、节点。 \\
状态 & 对象身上的变量 & 位置、分数、权限、颜色、是否激活。 \\
命令 & 输入中每种操作 & add、delete、query、move。 \\
规则 & 每条命令如何改变状态 & 如果存在则更新，否则新增。 \\
优先级 & 相同条件时谁先 & 时间小优先、编号小优先。 \\
边界 & 特殊或非法情况 & 找不到、重复、越界、空集合。 \\
输出 & 什么时候输出什么 & 每次 query 输出，最后统一输出。 \\
\bottomrule
\end{longtable}

\section{命令解析}\label{sec:07-11}

\subsection{固定格式命令}

如果每条命令格式固定：

\begin{lstlisting}
string op;
cin >> op;
if (op == "add") {
    int id, x;
    cin >> id >> x;
} else if (op == "query") {
    int id;
    cin >> id;
}
\end{lstlisting}

\subsection{整行命令}

如果命令中可能有空格，先读整行。

\begin{lstlisting}
int q;
cin >> q;
cin.ignore(numeric_limits<streamsize>::max(), '\n');

while (q--) {
    string line;
    getline(cin, line);
    stringstream ss(line);
    string op;
    ss >> op;

    if (op == "add") {
        string name;
        int x;
        ss >> name >> x;
    }
}
\end{lstlisting}

\subsection{手写 split}

按指定字符分割路径、日期、配置项：

\begin{lstlisting}
vector<string> split(const string& s, char sep) {
    vector<string> res;
    string cur;
    for (char c : s) {
        if (c == sep) {
            res.push_back(cur);
            cur.clear();
        } else {
            cur.push_back(c);
        }
    }
    res.push_back(cur);
    return res;
}
\end{lstlisting}

\subsection{命令分发模板}

命令较多时，可以让每种命令对应一个函数，主循环只负责读入和调用。

\begin{lstlisting}
void handleAdd() {
    int id, x;
    cin >> id >> x;
    // add logic
}

void handleDelete() {
    int id;
    cin >> id;
    // delete logic
}

void handleQuery() {
    int id;
    cin >> id;
    // query logic
}

void handleCommand() {
    string op;
    cin >> op;
    if (op == "add") handleAdd();
    else if (op == "delete") handleDelete();
    else if (op == "query") handleQuery();
}
\end{lstlisting}

\subsection{命令解析易错点}

\begin{itemize}
  \item 命令参数个数是否固定。
  \item 参数中是否可能包含空格。
  \item 命令大小写是否敏感。
  \item 非法命令是否会出现。
  \item 每条命令是否都可能输出。
  \item 查询不存在对象时输出什么。
\end{itemize}

\section{状态设计}\label{sec:07-12}

\subsection{用结构体保存对象}

\begin{lstlisting}
struct User {
    int id;
    string name;
    int score;
    bool active;
};

map<int, User> users;
\end{lstlisting}

\subsection{状态编号}

状态不多时可以用整数常量：

\begin{lstlisting}
const int WAITING = 0;
const int RUNNING = 1;
const int FINISHED = 2;
const int FAILED = 3;
\end{lstlisting}

也可以用 \Code{enum}：

\begin{lstlisting}
enum Status {
    WAITING,
    RUNNING,
    FINISHED,
    FAILED
};
\end{lstlisting}

\subsection{状态转移表}

复杂状态题建议先画表：

\begin{longtable}{p{3cm}p{3cm}p{4cm}p{4cm}}
\toprule
当前状态 & 操作 & 新状态 & 输出/副作用 \\
\midrule
\endfirsthead
\toprule
当前状态 & 操作 & 新状态 & 输出/副作用 \\
\midrule
\endhead
WAITING & start & RUNNING & 记录开始时间 \\
RUNNING & finish & FINISHED & 计算得分 \\
RUNNING & fail & FAILED & 输出错误信息 \\
\bottomrule
\end{longtable}

\section{状态机模拟}\label{sec:07-13}

\subsection{什么题会用到}

题面出现：

\begin{itemize}
  \item 一个对象有多个状态。
  \item 不同状态下，同一个命令效果不同。
  \item 状态按规则切换。
  \item 非法状态转换需要忽略或报错。
\end{itemize}

\subsection{状态机模板}

\begin{lstlisting}
const int WAITING = 0;
const int RUNNING = 1;
const int FINISHED = 2;

struct Job {
    int id;
    int status;
    int score;
};

void start(Job& job) {
    if (job.status != WAITING) {
        return;
    }
    job.status = RUNNING;
}

void finish(Job& job) {
    if (job.status != RUNNING) {
        return;
    }
    job.status = FINISHED;
    job.score += 10;
}
\end{lstlisting}

\subsection{状态机检查}

\begin{itemize}
  \item 初始状态是什么。
  \item 每个操作允许在哪些状态下发生。
  \item 不允许的操作是忽略、报错，还是改变成特殊状态。
  \item 终止状态是否还能被修改。
  \item 状态改变时是否还要更新其他字段。
\end{itemize}

\section{事件模拟}\label{sec:07-14}

\subsection{什么题会用到}

题面出现：

\begin{itemize}
  \item 按时间顺序发生事件。
  \item 事件有开始、结束、到达、离开。
  \item 同一时间有多个事件。
  \item 需要维护当前正在发生的对象集合。
\end{itemize}

\subsection{事件结构体}

\begin{lstlisting}
struct Event {
    int time;
    int type;
    int id;
};

vector<Event> events;
\end{lstlisting}

\subsection{事件排序}

\begin{lstlisting}
sort(events.begin(), events.end(), [](const Event& a, const Event& b) {
    if (a.time != b.time) return a.time < b.time;
    return a.type < b.type; // depends on statement
});
\end{lstlisting}

\subsection{事件处理模板}

\begin{lstlisting}
for (const Event& e : events) {
    if (e.type == 0) {
        // start event
    } else if (e.type == 1) {
        // end event
    } else {
        // query event
    }
}
\end{lstlisting}

\subsection{同一时间事件顺序}

这是事件模拟最容易错的地方。一定要从题面确定：

\begin{itemize}
  \item 同一时间，结束事件先处理还是开始事件先处理。
  \item 查询是在事件前还是事件后。
  \item 如果没有说明，要从样例推断，但不要过度猜。
  \item 排序比较器中必须体现这个顺序。
\end{itemize}

\section{优先级模拟}\label{sec:07-15}

\subsection{什么题会用到}

题面出现：

\begin{itemize}
  \item 每次选择优先级最高的任务。
  \item 优先级相同按时间、编号、字典序。
  \item 动态加入任务。
  \item 当前最大、当前最小。
\end{itemize}

\subsection{priority\_queue 模板}

\begin{lstlisting}
struct Task {
    int priority;
    int time;
    int id;
};

struct Cmp {
    bool operator()(const Task& a, const Task& b) const {
        if (a.priority != b.priority) return a.priority < b.priority;
        if (a.time != b.time) return a.time > b.time;
        return a.id > b.id;
    }
};

priority_queue<Task, vector<Task>, Cmp> pq;
\end{lstlisting}

这个比较器表示：

\begin{itemize}
  \item 优先级大的先出。
  \item 优先级相同，时间小的先出。
  \item 时间相同，编号小的先出。
\end{itemize}

\Warn{priority\_queue 的比较器和 sort 相反，容易写错。}不确定时，用几个任务手动推一下 \Code{top()} 是谁。

\section{复杂排序和排名}\label{sec:07-16}

\subsection{什么题会用到}

题面出现：

\begin{itemize}
  \item 多个排序关键字。
  \item 排名、积分榜、成绩表。
  \item 相同分数按编号、时间、字典序排序。
  \item 并列排名或去重排名。
\end{itemize}

\subsection{复杂排序模板}

\begin{lstlisting}
struct Player {
    int id;
    int score;
    int penalty;
    string name;
};

sort(a.begin(), a.end(), [](const Player& x, const Player& y) {
    if (x.score != y.score) return x.score > y.score;
    if (x.penalty != y.penalty) return x.penalty < y.penalty;
    if (x.name != y.name) return x.name < y.name;
    return x.id < y.id;
});
\end{lstlisting}

\subsection{并列排名}

\begin{lstlisting}
int rank = 1;
for (int i = 0; i < (int)a.size(); i++) {
    if (i > 0 && a[i].score != a[i - 1].score) {
        rank = i + 1;
    }
    cout << a[i].id << ' ' << rank << '\n';
}
\end{lstlisting}

\subsection{复杂排序检查}

\begin{itemize}
  \item 每个关键字是升序还是降序。
  \item 相同关键字时是否继续比较。
  \item 是否要求稳定排序。
  \item 字符串按字典序还是输入顺序。
  \item 比较器相等时不能返回 \Code{true}。
\end{itemize}

\section{表格和矩阵模拟}\label{sec:07-17}

\subsection{什么题会用到}

题面出现：

\begin{itemize}
  \item 表格、座位、棋盘、网格。
  \item 行列操作。
  \item 旋转、翻转、交换行列。
  \item 查询某个单元格。
\end{itemize}

\subsection{矩阵读入和输出}

\begin{lstlisting}
int n, m;
cin >> n >> m;
vector<vector<int>> a(n, vector<int>(m));
for (int i = 0; i < n; i++) {
    for (int j = 0; j < m; j++) {
        cin >> a[i][j];
    }
}
\end{lstlisting}

\subsection{交换两行}

\begin{lstlisting}
int r1, r2;
cin >> r1 >> r2;
--r1; --r2;
swap(a[r1], a[r2]);
\end{lstlisting}

\subsection{交换两列}

\begin{lstlisting}
int c1, c2;
cin >> c1 >> c2;
--c1; --c2;
for (int i = 0; i < n; i++) {
    swap(a[i][c1], a[i][c2]);
}
\end{lstlisting}

\subsection{顺时针旋转矩阵}

\begin{lstlisting}
vector<vector<int>> rotateClockwise(const vector<vector<int>>& a) {
    int n = (int)a.size();
    if (n == 0) return {};
    int m = (int)a[0].size();
    vector<vector<int>> b(m, vector<int>(n));
    for (int i = 0; i < n; i++) {
        for (int j = 0; j < m; j++) {
            b[j][n - 1 - i] = a[i][j];
        }
    }
    return b;
}
\end{lstlisting}

\subsection{表格模拟易错点}

\begin{itemize}
  \item 题目行列通常从 1 开始，代码常从 0 开始。
  \item 旋转后行数和列数会交换。
  \item 如果操作很多，真实移动整个矩阵可能超时，要考虑只维护行列映射。
  \item 输出每行空格格式。
\end{itemize}

\section{目录树和路径模拟}\label{sec:07-18}

\subsection{什么题会用到}

题面出现：

\begin{itemize}
  \item 文件、目录、路径。
  \item \Code{/a/b/c}、相对路径、绝对路径。
  \item 创建、删除、查询文件。
  \item 父目录、子目录。
\end{itemize}

\subsection{路径规范化}

下面处理绝对路径的 \Code{.} 和 \Code{..}；相对路径需先拼接当前目录。依赖前文 split：

\begin{lstlisting}
vector<string> normalizePath(const string& path) {
    vector<string> st;
    vector<string> parts = split(path, '/');
    for (string part : parts) {
        if (part.empty() || part == ".") {
            continue;
        } else if (part == "..") {
            if (!st.empty()) st.pop_back();
        } else {
            st.push_back(part);
        }
    }
    return st;
}
\end{lstlisting}

\subsection{路径转字符串}

\begin{lstlisting}
string joinPath(const vector<string>& parts) {
    if (parts.empty()) return "/";
    string res;
    for (const string& p : parts) {
        res += "/";
        res += p;
    }
    return res;
}
\end{lstlisting}

\subsection{目录树节点}

\begin{lstlisting}
struct Node {
    string name;
    map<string, int> child;
    bool isFile = false;
};

vector<Node> tree;
\end{lstlisting}

新建节点：

\begin{lstlisting}
int newNode(const string& name, bool isFile) {
    Node node;
    node.name = name;
    node.isFile = isFile;
    tree.push_back(node);
    return (int)tree.size() - 1;
}
\end{lstlisting}

\subsection{路径题易错点}

\begin{itemize}
  \item 根目录怎么表示。
  \item 连续斜杠是否可能出现。
  \item 路径末尾斜杠是否影响结果。
  \item 查询不存在路径时输出什么。
  \item 文件和目录是否允许同名。
\end{itemize}

\section{URL 路由和参数解析}\label{sec:07-19}

\subsection{什么题会用到}

题面出现：

\begin{itemize}
  \item URL、路径、路由。
  \item \Code{/user/123/profile}。
  \item 查询参数，例如 \Code{?a=1\&b=2}。
  \item 通配符、变量匹配。
\end{itemize}

\subsection{拆分 URL}

\begin{lstlisting}
struct Url {
    string path;
    map<string, string> query;
};

Url parseUrl(const string& s) {
    Url res;
    size_t qpos = s.find('?');
    if (qpos == string::npos) {
        res.path = s;
        return res;
    }

    res.path = s.substr(0, qpos);
    string qs = s.substr(qpos + 1);
    vector<string> pairs = split(qs, '&');
    for (string p : pairs) {
        if (p.empty()) continue;
        size_t eq = p.find('=');
        if (eq == string::npos) {
            res.query[p] = "";
        } else {
            string key = p.substr(0, eq);
            string val = p.substr(eq + 1);
            res.query[key] = val;
        }
    }
    return res;
}
\end{lstlisting}

\subsection{简单路由匹配}

模式 \Code{/user/<id>} 匹配 \Code{/user/123}。先定义后文的 \Code{parsePath}，或在函数前声明它。此简化版忽略连续/末尾斜杠，仅支持整段参数，不支持 URL 映射真题的类型参数和剩余路径规则：

\begin{lstlisting}
bool matchRoute(const string& pattern, const string& path,
                map<string, string>& params) {
    params.clear();
    map<string, string> candidate;
    vector<string> a = parsePath(pattern);
    vector<string> b = parsePath(path);
    if (a.size() != b.size()) return false;

    for (int i = 0; i < (int)a.size(); i++) {
        if (!a[i].empty() && a[i][0] == '<' && a[i].back() == '>') {
            string name = a[i].substr(1, (int)a[i].size() - 2);
            candidate[name] = b[i];
        } else if (a[i] != b[i]) {
            return false;
        }
    }
    params = move(candidate);
    return true;
}
\end{lstlisting}

\subsection{URL 题易错点}

\begin{itemize}
  \item 有没有查询参数。
  \item 参数值是否可能为空。
  \item 路径末尾斜杠是否等价。
  \item 多个路由都匹配时，第一条生效还是更具体的生效。
  \item URL 是否需要解码，题目不要求就不要自己加。
\end{itemize}

\section{配置和 JSON 类解析入口}\label{sec:07-20}

\subsection{什么题会用到}

题面出现：

\begin{itemize}
  \item 键值对。
  \item 花括号、方括号、冒号、逗号。
  \item 嵌套对象。
  \item 查询某个字段。
\end{itemize}

\subsection{简单键值配置}

每行形如 \Code{key=value}：

\begin{lstlisting}
int n;
cin >> n;
map<string, string> conf;
for (int i = 0; i < n; i++) {
    string line;
    cin >> line;
    int pos = (int)line.find('=');
    string key = line.substr(0, pos);
    string value = line.substr(pos + 1);
    conf[key] = value;
}
\end{lstlisting}

\subsection{读取带空格的配置值}

\begin{lstlisting}
string line;
getline(cin, line);
int pos = (int)line.find('=');
string key = line.substr(0, pos);
string value = line.substr(pos + 1);
\end{lstlisting}

\subsection{嵌套结构处理思路}

JSON 类题先看题目允许哪些数据类型和格式，只实现题目实际需要的部分：

\begin{itemize}
  \item 是否只有字符串和对象，没有数组。
  \item 是否保证格式合法。
  \item 查询是否只按点号路径。
  \item 是否只需要保存完整路径到值的映射。
\end{itemize}

如果只需要查询路径，可以把嵌套字段拍平成 key：

\begin{lstlisting}
// Example flattened keys:
// user.name -> Alice
// user.age -> 18
map<string, string> value;
\end{lstlisting}

\Warn{配置/JSON 题先利用限制。}CSP 题通常不会要求你写工业级 JSON 解析器，重点是按题目给定范围实现。

\section{字符串解析}\label{sec:07-21}

\subsection{什么题会用到}

遇到下面这些格式时，需要先想清楚字符串怎样拆开：

\begin{itemize}
  \item 路径，例如 \Code{/a/b/c}。
  \item 配置项，例如 \Code{key=value}。
  \item 表达式，例如括号、加减乘除。
  \item 带引号的字符串。
  \item 嵌套结构。
  \item 命令参数中可能包含空格。
\end{itemize}

\subsection{逐字符扫描模板}

\begin{lstlisting}
string s;
cin >> s;

for (int i = 0; i < (int)s.size(); i++) {
    char c = s[i];
    if (isdigit((unsigned char)c)) {
        int j = i;
        while (j < (int)s.size() && isdigit((unsigned char)s[j])) j++;
        string number = s.substr(i, j - i);
        i = j - 1;
    } else if (isalpha((unsigned char)c)) {
        int j = i;
        while (j < (int)s.size() && isalpha((unsigned char)s[j])) j++;
        string word = s.substr(i, j - i);
        i = j - 1;
    } else {
        // symbol
    }
}
\end{lstlisting}

\subsection{解析 key=value}

\begin{lstlisting}
string s;
cin >> s;
int pos = (int)s.find('=');
string key = s.substr(0, pos);
string value = s.substr(pos + 1);
\end{lstlisting}

\subsection{路径解析}

\begin{lstlisting}
vector<string> parsePath(const string& path) {
    vector<string> parts;
    string cur;
    for (char c : path) {
        if (c == '/') {
            if (!cur.empty()) {
                parts.push_back(cur);
                cur.clear();
            }
        } else {
            cur.push_back(c);
        }
    }
    if (!cur.empty()) parts.push_back(cur);
    return parts;
}
\end{lstlisting}

\subsection{括号匹配}

\begin{lstlisting}
bool checkBrackets(const string& s) {
    stack<char> st;
    for (char c : s) {
        if (c == '(' || c == '[' || c == '{') {
            st.push(c);
        } else if (c == ')' || c == ']' || c == '}') {
            if (st.empty()) return false;
            char t = st.top();
            st.pop();
            if (c == ')' && t != '(') return false;
            if (c == ']' && t != '[') return false;
            if (c == '}' && t != '{') return false;
        }
    }
    return st.empty();
}
\end{lstlisting}

\subsection{字符串解析易错点}

\begin{itemize}
  \item \Code{cin >> s} 不能读取空格。
  \item \Code{getline} 前可能需要 \Code{cin.ignore()}。
  \item \Code{find} 找不到时返回 \Code{string::npos}。
  \item \Code{substr(pos, len)} 的第二个参数是长度，不是结束位置。
  \item 连续分隔符和开头/结尾分隔符要单独考虑。
  \item 字符串数字可能超过 \Code{long long}。
\end{itemize}

\section{表达式解析入口}\label{sec:07-22}

\subsection{什么题会用到}

题面出现：

\begin{itemize}
  \item 表达式、括号、运算符。
  \item 加减乘除优先级。
  \item 逻辑表达式。
  \item 嵌套结构。
\end{itemize}

\subsection{只含加减的表达式}

\begin{lstlisting}
long long evalPlusMinus(const string& s) {
    long long ans = 0;
    long long num = 0;
    int sign = 1;

    for (int i = 0; i <= (int)s.size(); i++) {
        if (i < (int)s.size() && isspace((unsigned char)s[i])) continue;
        if (i < (int)s.size() && isdigit((unsigned char)s[i])) {
            num = num * 10 + (s[i] - '0');
        } else {
            ans += sign * num;
            num = 0;
            if (i < (int)s.size()) {
                if (s[i] == '+') sign = 1;
                else if (s[i] == '-') sign = -1;
            }
        }
    }
    return ans;
}
\end{lstlisting}

完整整数四则运算见 \Tref{610}；下面仅展示读取一个数字的基础片段。

\subsection{递归解析思想}

复杂表达式通常需要递归下降。新手先掌握思想：

\begin{lstlisting}
string s;
int pos = 0;

long long parseNumber() {
    long long x = 0;
    while (pos < (int)s.size() && isdigit((unsigned char)s[pos])) {
        x = x * 10 + (s[pos] - '0');
        pos++;
    }
    return x;
}
\end{lstlisting}

\Warn{表达式题如果没练过，不要强行写复杂递归下降。}先拿没有括号、没有优先级的部分分。

\section{缓存和队列模拟}\label{sec:07-23}

\subsection{什么题会用到}

题面出现：

\begin{itemize}
  \item 缓存、页面置换、最近使用。
  \item 队列、排队、服务窗口。
  \item 先进先出、最近最少使用。
  \item 容量限制，满了要删除一个。
\end{itemize}

\subsection{FIFO 缓存}

\begin{lstlisting}
int cap;
cin >> cap;
queue<int> q;
set<int> inCache;

for (int x : requests) {
    if (cap <= 0) continue;
    if (inCache.count(x)) {
        continue;
    }
    if ((int)q.size() == cap) {
        int old = q.front();
        q.pop();
        inCache.erase(old);
    }
    q.push(x);
    inCache.insert(x);
}
\end{lstlisting}

\subsection{LRU 缓存入口}

LRU 可用 \Code{list + unordered\_map}：链表记访问顺序，映射保存每个对象的位置。

\begin{lstlisting}
int cap;
list<int> order; // front is most recent
unordered_map<int, list<int>::iterator> pos;

void access(int x) {
    if (cap <= 0) return;
    if (pos.count(x)) {
        order.erase(pos[x]);
    } else if ((int)order.size() == cap) {
        int old = order.back();
        order.pop_back();
        pos.erase(old);
    }
    order.push_front(x);
    pos[x] = order.begin();
}
\end{lstlisting}

\subsection{队列模拟易错点}

\begin{itemize}
  \item 队列为空时不能取 \Code{front}。
  \item 缓存容量可能为 0，要看题目是否允许。
  \item 命中缓存时是否需要更新时间。
  \item 删除元素时，辅助集合或 map 也要同步删除。
\end{itemize}

\section{懒删除优先队列}\label{sec:07-24}

\subsection{什么题会用到}

题面出现：

\begin{itemize}
  \item 需要动态删除任务。
  \item 还要每次取最大/最小。
  \item \Code{priority\_queue} 中间删除不方便。
\end{itemize}

\subsection{懒删除思想}

不直接从堆里删除旧元素，而是在取堆顶时检查它是否还有效。

\begin{lstlisting}
struct Task {
    int priority;
    int id;
};

struct Cmp {
    bool operator()(const Task& a, const Task& b) const {
        if (a.priority != b.priority) return a.priority < b.priority;
        return a.id > b.id;
    }
};

priority_queue<Task, vector<Task>, Cmp> pq;
map<int, int> currentPriority;

void addOrUpdate(int id, int priority) {
    currentPriority[id] = priority;
    pq.push({priority, id});
}

void removeTask(int id) {
    currentPriority.erase(id);
}

Task getTop() {
    while (!pq.empty()) {
        Task t = pq.top();
        if (currentPriority.count(t.id) &&
            currentPriority[t.id] == t.priority) {
            return t;
        }
        pq.pop();
    }
    return {-1, -1};
}
\end{lstlisting}

\Warn{懒删除会让堆里有废数据。}每次取 \Code{top} 前必须清理过期元素。

\section{时间和日程模拟}\label{sec:07-25}

\subsection{什么题会用到}

题面出现：

\begin{itemize}
  \item 时间段、会议、日程。
  \item 判断冲突。
  \item 开始时间、结束时间。
  \item 按时间排序处理。
\end{itemize}

\subsection{时间转分钟}

\begin{lstlisting}
int toMinute(const string& s) {
    int h = stoi(s.substr(0, 2));
    int m = stoi(s.substr(3, 2));
    return h * 60 + m;
}
\end{lstlisting}

\subsection{判断区间冲突}

两个半开区间 $[l1,r1)$ 和 $[l2,r2)$ 冲突：

\begin{lstlisting}
bool overlap(int l1, int r1, int l2, int r2) {
    return max(l1, l2) < min(r1, r2);
}
\end{lstlisting}

如果题目是闭区间，条件要改。

\subsection{扫描事件统计同时进行数量}

\begin{lstlisting}
vector<pair<int, int>> events;
for (auto [l, r] : intervals) {
    if (l >= r) continue; // empty half-open interval
    events.push_back({l, 1});
    events.push_back({r, -1});
}

sort(events.begin(), events.end(), [](auto a, auto b) {
    if (a.first != b.first) return a.first < b.first;
    return a.second < b.second; // end before start for [l,r)
});

int cur = 0, best = 0;
for (auto [time, delta] : events) {
    cur += delta;
    best = max(best, cur);
}
\end{lstlisting}

\subsection{时间题易错点}

\begin{itemize}
  \item 时间段是闭区间还是半开区间。
  \item 结束时间和开始时间相同是否冲突。
  \item 跨天如何处理。
  \item 同一时间开始和结束谁先处理。
\end{itemize}

\section{树形结构模拟}\label{sec:07-26}

\subsection{什么题会用到}

题面出现：

\begin{itemize}
  \item 层级、父子、目录、组织结构。
  \item 每个节点有父节点。
  \item 查询祖先、子树、路径。
\end{itemize}

\subsection{父子关系存储}

\begin{lstlisting}
int n;
cin >> n;
vector<vector<int>> children(n + 1);
vector<int> parent(n + 1, 0);

for (int i = 2; i <= n; i++) {
    int p;
    cin >> p;
    parent[i] = p;
    children[p].push_back(i);
}
\end{lstlisting}

\subsection{DFS 遍历树}

\begin{lstlisting}
vector<int> depth(n + 1, 0);

void dfsTree(int u) {
    for (int v : children[u]) {
        depth[v] = depth[u] + 1;
        dfsTree(v);
    }
}
\end{lstlisting}

\subsection{子树大小}

\begin{lstlisting}
vector<int> subtree(n + 1, 1);

void calcSubtree(int u) {
    subtree[u] = 1;
    for (int v : children[u]) {
        calcSubtree(v);
        subtree[u] += subtree[v];
    }
}
\end{lstlisting}

\subsection{树题易错点}

\begin{itemize}
  \item 根节点是谁。
  \item 输入给的是父节点，还是边。
  \item 节点编号从 0 还是 1 开始。
  \item 树很深时递归可能爆栈。
  \item 如果需要频繁祖先查询，简单爬父亲可能超时。
\end{itemize}

\section{权限和规则匹配}\label{sec:07-27}

\subsection{什么题会用到}

题面出现：

\begin{itemize}
  \item 用户、角色、权限。
  \item 规则匹配。
  \item 黑名单、白名单。
  \item 多条规则按顺序匹配，第一条生效。
  \item 默认允许或默认拒绝。
\end{itemize}

\subsection{规则结构体}

\begin{lstlisting}
struct Rule {
    string user;
    string action;
    string resource;
    bool allow;
};

vector<Rule> rules;
\end{lstlisting}

\subsection{规则匹配函数}

\begin{lstlisting}
bool match(const Rule& r, const string& user,
           const string& action, const string& resource) {
    if (r.user != "*" && r.user != user) return false;
    if (r.action != "*" && r.action != action) return false;
    if (r.resource != "*" && r.resource != resource) return false;
    return true;
}
\end{lstlisting}

\subsection{按顺序匹配}

\begin{lstlisting}
bool checkPermission(const string& user,
                     const string& action,
                     const string& resource) {
    for (const Rule& r : rules) {
        if (match(r, user, action, resource)) {
            return r.allow;
        }
    }
    return false; // default deny
}
\end{lstlisting}

\subsection{规则题易错点}

\begin{itemize}
  \item 多条规则是第一条生效，还是最后一条生效。
  \item 默认值是允许还是拒绝。
  \item 通配符如何匹配。
  \item 大小写是否敏感。
  \item 资源路径是否需要前缀匹配。
\end{itemize}

\section{网格和图混合题}\label{sec:07-28}

\subsection{什么题会用到}

题面出现：

\begin{itemize}
  \item 地图、障碍、传送门、钥匙、门。
  \item 不只是简单最短路，还有状态。
  \item 每个格子可能有属性。
  \item 移动过程中会改变状态。
\end{itemize}

\subsection{普通网格 BFS}

如果只有障碍和空地，用 B 题 BFS 模板即可。

\subsection{带状态 BFS 入口}

如果移动时还带着状态，例如是否拿到钥匙，可以把状态加入距离数组。

\begin{lstlisting}
struct Node {
    int x;
    int y;
    int state;
};

vector<vector<vector<int>>> dist(n,
    vector<vector<int>>(m, vector<int>(S, -1)));

queue<Node> q;
dist[sx][sy][0] = 0;
q.push({sx, sy, 0});

while (!q.empty()) {
    Node cur = q.front();
    q.pop();

    for (int d = 0; d < 4; d++) {
        int nx = cur.x + dx[d];
        int ny = cur.y + dy[d];
        int ns = cur.state;
        if (nx < 0 || nx >= n || ny < 0 || ny >= m) continue;
        if (g[nx][ny] == '#') continue;
        // Check door permission, then update ns from g[nx][ny].
        // Do not index the grid before the bounds check.
        if (dist[nx][ny][ns] != -1) continue;
        dist[nx][ny][ns] = dist[cur.x][cur.y][cur.state] + 1;
        q.push({nx, ny, ns});
    }
}
\end{lstlisting}

本段仅是状态扩展框架，必须按题意补全门禁检查、状态转移和起点状态。空间约为 $4nmS$ 字节，先算总量。

\Warn{状态数不能太大。}如果状态有 $2^k$，先看 $k$ 是否很小。$k \le 10$ 常见可做，$k=30$ 通常不行。

\subsection{图模拟策略}

如果题目是图但规则很多：

\begin{itemize}
  \item 先把图建对：有向/无向、权值、编号。
  \item 再处理规则：哪些边能走，什么时候能走。
  \item 如果边权全为 1，用 BFS。
  \item 如果边权为正，用 Dijkstra。
  \item 如果只是连通关系，用并查集或 DFS。
\end{itemize}

\section{大模拟易错点}\label{sec:07-29}

\begin{itemize}
  \item 同一时间多个事件的顺序。
  \item 相同优先级的并列规则。
  \item 删除元素后迭代器失效。
  \item 查询不存在对象。
  \item 重复添加同一对象。
  \item 空集合时取最大、最小。
  \item 最后一轮事件是否处理。
  \item 题目中的编号是 0 开始还是 1 开始。
  \item 输出顺序是否要求按编号、时间或字典序。
\end{itemize}

\section{C 题分阶段实现}\label{sec:07-30}

\subsection{不要一次写完}

每写完一种命令就运行一次，检查前后状态。可以按下面的顺序逐步完成：

\begin{enumerate}
  \item 只写读入，打印读到的数据检查。
  \item 写数据结构，不写复杂规则。
  \item 实现第一种命令或最基础操作。
  \item 跑样例，看是否能得到部分中间结果。
  \item 再补第二种命令。
  \item 最后处理特殊规则、并列规则、异常情况。
\end{enumerate}

\subsection{命令题实现顺序}

\begin{longtable}{p{3cm}p{6cm}p{5cm}}
\toprule
阶段 & 做什么 & 目标 \\
\midrule
\endfirsthead
\toprule
阶段 & 做什么 & 目标 \\
\midrule
\endhead
1 & 读入所有命令，只识别命令名 & 确认输入没错。 \\
2 & 实现新增/初始化类命令 & 建立基本状态。 \\
3 & 实现查询类命令 & 能输出最简单答案。 \\
4 & 实现修改/删除类命令 & 状态开始完整。 \\
5 & 实现特殊规则 & 冲满分。 \\
\bottomrule
\end{longtable}

\subsection{模拟题实现顺序}

\begin{enumerate}
  \item 先实现没有特殊情况的主流程。
  \item 再实现边界检查。
  \item 再实现冲突和优先级。
  \item 最后实现异常输入或题目特别说明的细节。
\end{enumerate}

\section{拆分函数时的检查}\label{sec:07-31}

代码较长时，把解析、查询和修改拆开，定位问题会方便一些。

\subsection{建议函数类型}

\begin{longtable}{p{4cm}p{5cm}p{5cm}}
\toprule
函数类型 & 示例 & 作用 \\
\midrule
\endfirsthead
\toprule
函数类型 & 示例 & 作用 \\
\midrule
\endhead
读入函数 & \Code{readInput()} & 避免读入和处理混在一起。 \\
解析函数 & \Code{parsePath(s)} & 字符串处理单独测试。 \\
查询函数 & \Code{query(id)} & 查询逻辑独立。 \\
修改函数 & \Code{update(id,x)} & 状态改变集中处理。 \\
检查函数 & \Code{canMove(x,y)} & 规则判断集中处理。 \\
输出函数 & \Code{printAnswer()} & 控制格式。 \\
\bottomrule
\end{longtable}

\subsection{函数设计原则}

\begin{itemize}
  \item 一个函数只做一类事情。
  \item 函数名直接反映题目操作。
  \item 一个函数尽量只处理一类操作；几种互不相关的规则可以拆开。
  \item 修改状态的函数要特别小心副作用。
  \item 查询函数尽量不要修改状态，除非题目要求。
\end{itemize}

\subsection{推荐代码顺序}

\begin{lstlisting}
// constants and types
struct ...

// global variables
int n, m;
vector<...> ...

// helper functions
bool inside(...);
vector<string> split(...);

// operation functions
void handleAdd(...);
void handleQuery(...);

// main solve
void solve() { ... }
\end{lstlisting}

\section{分层测试方法}\label{sec:07-32}

\subsection{第一层：解析测试}

只测试输入解析，不测算法。

\begin{itemize}
  \item 路径是否被正确分割。
  \item 命令参数是否读对。
  \item 数字字符串是否转换正确。
  \item 空格和换行是否处理正确。
\end{itemize}

\subsection{第二层：单操作测试}

每种操作单独测：

\begin{itemize}
  \item 只 add 一次。
  \item 只 query 一次。
  \item add 后 query。
  \item delete 后 query。
\end{itemize}

\subsection{第三层：组合测试}

测试多操作交替：

\begin{itemize}
  \item add、delete、add 同一个对象。
  \item 多个对象按优先级排序。
  \item 同一时间多个事件。
  \item 状态反复切换。
\end{itemize}

\subsection{第四层：边界测试}

测试异常和边界：

\begin{itemize}
  \item 空集合查询。
  \item 不存在对象。
  \item 重复对象。
  \item 最大数据附近。
  \item 最后一条命令。
\end{itemize}

\section{C 题常见失分原因}\label{sec:07-33}

\begin{longtable}{p{4cm}p{5cm}p{5cm}}
\toprule
失分原因 & 表现 & 预防 \\
\midrule
\endfirsthead
\toprule
失分原因 & 表现 & 预防 \\
\midrule
\endhead
读题漏规则 & 样例某一步对不上 & 三遍读题，列规则表。 \\
状态没初始化 & 第一次查询就错 & 结构体构造或统一 init。 \\
命令读错 & 后续全错 & 先调试输出 op 和参数。 \\
比较器错 & 排名/优先级错 & 手写 3 个对象测试 top/order。 \\
删除不同步 & 查到已删除对象 & 容器、set、map 同步维护。 \\
最后事件漏处理 & 前面都对，最后错 & 检查循环结束后的收尾。 \\
输出顺序错 & 数值对但顺序错 & 看题目要求按什么顺序。 \\
复杂度过高 & 小样例过，大数据超时 & 估算每条命令复杂度。 \\
\bottomrule
\end{longtable}

\section{C 题本地调试}\label{sec:07-34}

\subsection{调试输出建议}

调试时输出状态：

\begin{lstlisting}
void debugTask(const Task& t) {
    cerr << "id=" << t.id
         << " status=" << t.status
         << " priority=" << t.priority << '\n';
}
\end{lstlisting}

处理命令时：

\begin{lstlisting}
cerr << "op=" << op << '\n';
\end{lstlisting}

\Warn{提交前删除所有调试输出。}C 题代码长，尤其容易漏掉 \Code{cerr} 或临时 \Code{cout}。

\subsection{缩小样例}

如果官方样例很长，自己造小样例：

\begin{itemize}
  \item 1 个对象，1 条命令。
  \item 1 个对象，2 条命令。
  \item 2 个对象，测试排序或优先级。
  \item 不存在对象的查询。
  \item 重复添加和删除。
\end{itemize}

\subsection{定位 bug 的方法}

\begin{enumerate}
  \item 先确认读入正确。
  \item 再确认每条命令被识别正确。
  \item 再确认命令前状态正确。
  \item 再确认命令后状态正确。
  \item 最后确认输出格式。
\end{enumerate}

\section{C 题常见边界库}\label{sec:07-35}

\begin{longtable}{p{4cm}p{5cm}p{5cm}}
\toprule
题型 & 边界 & 检查目标 \\
\midrule
\endfirsthead
\toprule
题型 & 边界 & 检查目标 \\
\midrule
\endhead
命令模拟 & 只有 1 条命令、没有查询、连续查询 & 初始化和输出。 \\
新增删除 & 重复新增、删除不存在、删除后查询 & 对象生命周期。 \\
状态机 & 非法状态转换、终止状态继续操作 & 状态规则。 \\
事件模拟 & 同一时间多个事件、开始和结束同一时刻 & 事件排序。 \\
优先级 & 优先级相同、时间相同、编号相同附近 & 比较器。 \\
字符串解析 & 空字段、连续分隔符、路径以分隔符结尾 & split 和 substr。 \\
网格 & 起点等于终点、无路、只有一个格子 & BFS 初始化。 \\
带状态 BFS & 拿不到状态、重复状态、状态数最大 & dist 维度。 \\
图 & 自环、重边、不连通、有向/无向 & 建图。 \\
\bottomrule
\end{longtable}

\section{C 题部分分策略}\label{sec:07-36}
先看子任务限制。比如“没有删除命令”“没有嵌套结构”或“对象数量很少”，都可能允许更短的实现。选定一档后，把这一档会出现的规则写完整；只实现几条命令，却忽略它们会遇到的并列和边界规则，未必能得分。

\section{按子任务逐步完成}\label{sec:07-37}
可以按下面的顺序写，每完成一步就用小样例检查一次。

\subsection{读入与基本操作}
确认所有输入都能读完，每种命令的参数个数正确。实现初始状态和最简单的操作，检查一次新增或修改之后，查询能否得到预期结果。

\subsection{完成选定子任务}
补齐这个子任务会用到的所有命令、并列规则和异常处理。再试重复添加、删除不存在对象、空集合查询等边界。测试点有多次输出时，所有输出都要符合题意。

\subsection{扩大适用范围}
时间允许时，再加入嵌套、更复杂的规则或更快的数据结构。修改前保留已经检查过的版本，新版本通过对照后再替换。若题目采用捆绑测试，还要满足该子任务的全部要求。

\section{来不及写完时}\label{sec:07-38}

\subsection{写不完大模拟}

\begin{itemize}
  \item 保留已完成命令，不要为了补复杂规则改崩基础逻辑。
  \item 先完成子任务所需的状态修改，再检查相应查询输出。
  \item 特殊规则用清晰的 \Code{if} 补，不要重构大段代码。
\end{itemize}

\subsection{字符串解析写不出来}

\begin{itemize}
  \item 若有不含空格、不含嵌套的子任务，可以先处理这一档。
  \item 先用 \Code{split} 处理简单分隔符。
  \item 没有嵌套的子任务仍要正确处理其中所有 token 和输出。
\end{itemize}

\subsection{状态 BFS 写不出来}

\begin{itemize}
  \item 若子任务不含特殊状态，可以先写普通 BFS。
  \item 如果状态只有 2 种，尝试 \Code{dist[x][y][2]}。
  \item 如果有不含特殊状态的子任务，可以先完成这一档。
\end{itemize}

\subsection{复杂度可能超时}

\begin{itemize}
  \item 先估每条命令复杂度。
  \item 能用 map 查找就不要全表扫描。
  \item 能预处理就不要每次重复计算。
  \item 时间不够时，保留暴力拿小数据分。
\end{itemize}

\section{不同题型的实现顺序}\label{sec:07-39}

\subsection{如果是大模拟}

\begin{enumerate}
  \item 先读懂输入输出。
  \item 建结构体保存对象。
  \item 实现最简单命令。
  \item 实现查询输出。
  \item 再补删除、修改、异常情况。
  \item 最后补优先级、并列规则。
\end{enumerate}

\subsection{如果是字符串解析}

\begin{enumerate}
  \item 先确认用 \Code{cin} 还是 \Code{getline}。
  \item 先把字符串拆成 token。
  \item 能用 \Code{split} 就不要写复杂解析。
  \item 如果有括号，先写括号匹配。
  \item 如果有嵌套，先拿无嵌套部分分。
\end{enumerate}

\subsection{如果是图或网格}

\begin{enumerate}
  \item 先把图或网格读入正确。
  \item 判断是最短路、连通块还是模拟移动。
  \item 无权最短路用 BFS。
  \item 带状态时先估算状态数量。
  \item 只有子任务确实不含特殊状态时，才能用普通 BFS 简化。
\end{enumerate}

\subsection{如果是规则匹配}

\begin{enumerate}
  \item 把每条规则保存为结构体。
  \item 单独写 \Code{match} 函数。
  \item 确认第一条生效还是最后一条生效。
  \item 确认默认值。
  \item 用 2 到 3 条规则手测。
\end{enumerate}

\section{C 题模板索引}\label{sec:07-40}

\begin{longtable}{p{4cm}p{5cm}p{5cm}}
\toprule
题面关键词 & 使用模板 & 本章位置 \\
\midrule
\endfirsthead
\toprule
题面关键词 & 使用模板 & 本章位置 \\
\midrule
\endhead
多命令 & 命令分发 & 命令解析 \\
对象状态 & struct + 状态机 & 状态设计、状态机模拟 \\
按时间发生 & Event + sort & 事件模拟 \\
优先级 & priority\_queue & 优先级模拟 \\
排名 & 多关键字 sort & 复杂排序和排名 \\
表格、矩阵 & vector<vector<int>> & 表格和矩阵模拟 \\
路径、目录 & split + normalizePath & 目录树和路径模拟 \\
URL、参数 & parseUrl & URL 路由和参数解析 \\
配置、键值 & key=value & 配置和 JSON 类解析 \\
括号、表达式 & stack / 递归解析 & 字符串解析、表达式解析 \\
缓存 & queue / list + map & 缓存和队列模拟 \\
动态删除最大值 & 懒删除堆 & 懒删除优先队列 \\
时间段冲突 & overlap / events & 时间和日程模拟 \\
父子层级 & children + DFS & 树形结构模拟 \\
地图状态 & state BFS & 网格和图混合题 \\
\bottomrule
\end{longtable}

\section{C 题考场速查页}\label{sec:07-41}

\subsection{读题 5 分钟内要确定}

\begin{itemize}
  \item 是大模拟、字符串解析、图网格、数据结构，还是 DP。
  \item 有哪些对象。
  \item 每个对象有哪些状态。
  \item 有哪些命令或事件。
  \item 输出发生在每次查询还是最后。
  \item 数据范围是否允许全表扫描。
\end{itemize}

\subsection{写代码前先想清楚}

\begin{itemize}
  \item \Code{struct} 里有哪些字段。
  \item 用 \Code{vector}、\Code{map}、\Code{set} 还是堆。
  \item 每个命令对应哪个函数。
  \item 比较器的排序规则。
  \item 不存在、重复、越界怎么处理。
\end{itemize}

\subsection{最常用代码形状}

\begin{lstlisting}
struct Obj {
    int id;
    int status;
};

map<int, Obj> obj;

void handleAdd() {}
void handleDelete() {}
void handleQuery() {}

void solve() {
    int q;
    cin >> q;
    while (q--) {
        string op;
        cin >> op;
        if (op == "add") handleAdd();
        else if (op == "delete") handleDelete();
        else if (op == "query") handleQuery();
    }
}
\end{lstlisting}

\subsection{最后 10 分钟检查}

\begin{itemize}
  \item 删除调试输出。
  \item 检查所有命令是否都处理。
  \item 检查输出大小写、空格、换行。
  \item 检查比较器并列规则。
  \item 检查空集合、找不到对象。
  \item 检查最后一个事件。
  \item 提交一个当前最稳版本。
\end{itemize}

\section{C 题提交前检查}\label{sec:07-42}

\begin{itemize}
  \item 是否删除所有调试输出。
  \item 每种命令是否都处理了。
  \item 每种状态是否都有初始值。
  \item 不存在对象时是否按题意处理。
  \item 重复对象是否按题意处理。
  \item 排序和优先级比较器是否正确。
  \item 字符串读取是否会漏掉空格。
  \item 多组数据是否清空所有容器。
  \item 最后一个事件或最后一段数据是否处理。
  \item 输出顺序是否符合题目。
\end{itemize}

\section{C 题模板练习清单}\label{sec:07-43}

\begin{longtable}{p{4cm}p{4cm}p{6cm}}
\toprule
模板 & 熟练要求 & 用途 \\
\midrule
\endfirsthead
\toprule
模板 & 熟练要求 & 用途 \\
\midrule
\endhead
结构体 + map 保存对象 & 建议掌握 & 对象状态模拟。 \\
命令分发 & 建议掌握 & 多命令题。 \\
getline + stringstream & 建议掌握 & 整行命令。 \\
split & 照着能写 & 路径、配置解析。 \\
状态机 & 照着能写 & 多状态对象。 \\
事件排序 & 照着能写 & 时间顺序模拟。 \\
priority\_queue 自定义比较 & 练过再用 & 优先级模拟。 \\
括号匹配 & 照着能写 & 表达式/嵌套结构。 \\
带状态 BFS & 理解后再用 & 网格状态题。 \\
\bottomrule
\end{longtable}

\chapter{D/E 题部分分策略与高级算法入口}\label{ch:08}

\section{D/E 题定位}\label{sec:08-1}

D/E 题可以先找小数据、特殊结构，或自己练过的做法。能独立完成一个子任务，就先把它写好；暂时不会的优化留到后面。

\begin{itemize}
  \item D 题先看子任务，按会写的做法估算能覆盖哪些数据。
  \item E 题也可以找小数据或特殊性质，但要给前面会做的题留够时间。
  \item 先完成选定子任务，再考虑扩大适用范围。
  \item 思考很久仍没进展时，可以先保存思路，转做其他题。
\end{itemize}

如果已经花了十几二十分钟仍没有完整思路，可以重新看子任务表，或者暂时换题。

\section{D/E 题读题流程}\label{sec:08-2}

\begin{enumerate}
  \item 先看有没有子任务或部分分说明。
  \item 抄下所有数据范围，尤其是小数据档。
  \item 找特殊性质：链、树、DAG、无修改、无查询、边权相同、参数很小。
  \item 写出最暴力做法。
  \item 估算暴力能拿哪些数据点。
  \item 如果会优化，再逐步优化。
\end{enumerate}

\section{部分分决策树}\label{sec:08-3}

\begin{longtable}{p{4.5cm}p{5cm}p{4.5cm}}
\toprule
题面或数据特征 & 可写做法 & 目标 \\
\midrule
\endfirsthead
\toprule
题面或数据特征 & 可写做法 & 目标 \\
\midrule
\endhead
$n \le 20$ & DFS、枚举子集、状态压缩 & 小数据分 \\
$n \le 100$ & $O(n^3)$、Floyd、区间 DP & 小中数据 \\
$n \le 1000$ & $O(n^2)$ & 中小数据 \\
图是树 & DFS、树形 DP、LCA 入门 & 特殊性质分 \\
图是链 & 转成数组问题 & 特殊性质分 \\
边权全为 1 & BFS & 替代 Dijkstra \\
没有修改 & 前缀和、静态预处理 & 静态数据分 \\
查询很少 & 每次暴力 & 小 $q$ 分 \\
修改很少 & 修改后重算 & 小修改分 \\
参数 $k$ 很小 & 状压、枚举 $k$ & 参数小分 \\
求阈值且可行性单调 & 二分答案（check 须正确） & 进一步优化 \\
\bottomrule
\end{longtable}

\section{暴力模板}\label{sec:08-4}

\subsection{枚举所有点对}

\begin{lstlisting}
long long ans = 0;
for (int i = 0; i < n; i++) {
    for (int j = i + 1; j < n; j++) {
        // check pair (i, j)
    }
}
\end{lstlisting}

适用：$n \le 5000$ 需谨慎，$n \le 1000$ 通常可试。

\subsection{枚举所有区间}

\begin{lstlisting}
for (int l = 0; l < n; l++) {
    long long sum = 0;
    for (int r = l; r < n; r++) {
        sum += a[r];
        // process interval [l, r]
    }
}
\end{lstlisting}

适用：$n \le 5000$ 以内看时限；如果 $n=10^5$，只能拿小数据分。

\subsection{DFS 枚举选择}

\begin{lstlisting}
int n;
vector<int> choose;

void dfs(int idx) {
    if (idx == n) {
        // evaluate choose
        return;
    }

    // not choose idx
    dfs(idx + 1);

    // choose idx
    choose.push_back(idx);
    dfs(idx + 1);
    choose.pop_back();
}
\end{lstlisting}

适用：$n \le 20$ 左右。

\section{小数据常用算法}\label{sec:08-5}

\subsection{Floyd}

点数小且问任意两点最短路：

\begin{lstlisting}
for (int k = 1; k <= n; k++) {
    for (int i = 1; i <= n; i++) {
        for (int j = 1; j <= n; j++) {
            if (dist[i][k] == INF || dist[k][j] == INF) continue;
            dist[i][j] = min(dist[i][j], dist[i][k] + dist[k][j]);
        }
    }
}
\end{lstlisting}

\subsection{全排列}

\begin{lstlisting}
vector<int> p(n);
iota(p.begin(), p.end(), 0);
do {
    // evaluate permutation
} while (next_permutation(p.begin(), p.end()));
\end{lstlisting}

适用：$n \le 9$ 或 $10$ 左右。

\subsection{状压枚举子集}

\begin{lstlisting}
for (int mask = 0; mask < (1 << n); mask++) {
    for (int i = 0; i < n; i++) {
        if (mask & (1 << i)) {
            // i is selected
        }
    }
}
\end{lstlisting}

适用：$n \le 20$ 左右。若 $n \ge 31$，注意 \Code{1LL << n} 和复杂度。

\section{特殊性质策略}\label{sec:08-6}

\subsection{图是链}

如果图是一条链，很多图问题可以转成数组问题：

\begin{itemize}
  \item 距离变成下标差或前缀和。
  \item 子树变成连续区间。
  \item 路径查询变成区间查询。
\end{itemize}

\subsection{图是树}

树比一般图简单：

\begin{itemize}
  \item 没有环。
  \item 两点之间路径唯一。
  \item 可用 DFS 统计子树。
  \item 可用树形 DP。
\end{itemize}

\subsection{边权相同}

如果所有边权都是 1：

\begin{itemize}
  \item 最短路用 BFS。
  \item 不需要 Dijkstra。
\end{itemize}

\subsection{无修改}

如果没有修改，只有查询：

\begin{itemize}
  \item 优先预处理。
  \item 区间问题考虑前缀和。
  \item 图问题考虑一次 BFS/Dijkstra 后回答多次查询。
\end{itemize}

\section{按题型拿部分分}\label{sec:08-7}

\subsection{图论题}

\begin{longtable}{p{4cm}p{5cm}p{5cm}}
\toprule
不会满分时 & 可写做法 & 可能覆盖 \\
\midrule
\endfirsthead
\toprule
不会满分时 & 可写做法 & 可能覆盖 \\
\midrule
\endhead
不会复杂最短路 & 小图 Floyd & $n$ 小的数据 \\
不会带权最短路 & 无权 BFS & 边权全 1 的数据 \\
不会动态图 & 每次修改后重跑 BFS/DFS & 操作数小的数据 \\
不会强连通 & 小图逐点搜索，判断双向可达 & $O(n(n+m))$，单次 DFS 不能代替强连通 \\
不会树上高级算法 & 从每个查询端点暴力爬/DFS & $n,q$ 小的数据 \\
\bottomrule
\end{longtable}

\subsection{区间数据结构题}

\begin{longtable}{p{4cm}p{5cm}p{5cm}}
\toprule
不会满分时 & 可写做法 & 可能覆盖 \\
\midrule
\endfirsthead
\toprule
不会满分时 & 可写做法 & 可能覆盖 \\
\midrule
\endhead
不会线段树 & 每次暴力扫描区间 & 小 $nq$ 数据 \\
只有查询无修改 & 前缀和 & 静态数据 \\
只有区间修改最后输出 & 差分 & 离线数据 \\
单点修改区间查询 & 树状数组入口 & 中等数据 \\
区间修改单点查询 & 差分思想 / 树状数组 & 中等数据 \\
\bottomrule
\end{longtable}

\subsection{DP 题}

\begin{longtable}{p{4cm}p{5cm}p{5cm}}
\toprule
不会满分时 & 可写做法 & 可能覆盖 \\
\midrule
\endfirsthead
\toprule
不会满分时 & 可写做法 & 可能覆盖 \\
\midrule
\endhead
状态想不清 & DFS 暴力枚举 & $n$ 小的数据 \\
背包优化不会 & 二维 DP & 小容量数据 \\
状态压缩不会 & 回溯搜索 & $n$ 小的数据 \\
转移太复杂 & 只处理部分状态或特殊参数 & 特殊性质数据 \\
\bottomrule
\end{longtable}

\subsection{字符串题}

\begin{longtable}{p{4cm}p{5cm}p{5cm}}
\toprule
不会满分时 & 可写做法 & 可能覆盖 \\
\midrule
\endfirsthead
\toprule
不会满分时 & 可写做法 & 可能覆盖 \\
\midrule
\endhead
不会 KMP & \Code{find} 或暴力匹配 & 短字符串数据 \\
不会 Trie & map 存完整字符串 & 字符串数量较少 \\
不会复杂解析 & split 简单格式 & 无嵌套数据 \\
不会哈希 & 直接比较字符串 & 小数据 \\
\bottomrule
\end{longtable}

\section{树状数组入口}\label{sec:08-8}

\subsection{什么题会用到}

\begin{itemize}
  \item 单点修改，前缀和或区间和查询。
  \item 需要动态维护数组和。
  \item 数据范围 $n,q \le 10^5$。
\end{itemize}

\subsection{树状数组模板}

\begin{lstlisting}
struct BIT {
    int n;
    vector<long long> tree;

    BIT(int n) : n(n), tree(n + 1, 0) {}

    void add(int idx, long long val) {
        assert(1 <= idx && idx <= n); // idx=0 would never advance
        for (; idx <= n; idx += idx & -idx) {
            tree[idx] += val;
        }
    }

    long long sumPrefix(int idx) {
        assert(0 <= idx && idx <= n);
        long long res = 0;
        for (; idx > 0; idx -= idx & -idx) {
            res += tree[idx];
        }
        return res;
    }

    long long rangeSum(int l, int r) {
        return sumPrefix(r) - sumPrefix(l - 1);
    }
};
\end{lstlisting}

\Warn{树状数组一般用 1 下标。}

\section{线段树入口}\label{sec:08-9}

\subsection{什么题会用到}

\begin{itemize}
  \item 区间查询和区间修改交替出现。
  \item 需要维护区间和、区间最大值、区间最小值。
  \item $n,q \le 10^5$，暴力扫描会超时。
\end{itemize}

\subsection{区间和线段树框架}

\begin{lstlisting}
struct SegTree {
    int n;
    vector<long long> tree, lazy;

    SegTree(int n) : n(n), tree(4 * n + 4), lazy(4 * n + 4) {}

    void apply(int node, int l, int r, long long val) {
        tree[node] += val * (r - l + 1);
        lazy[node] += val;
    }

    void push(int node, int l, int r) {
        if (lazy[node] == 0 || l == r) return;
        int mid = (l + r) / 2;
        apply(node * 2, l, mid, lazy[node]);
        apply(node * 2 + 1, mid + 1, r, lazy[node]);
        lazy[node] = 0;
    }

    void rangeAdd(int node, int l, int r, int ql, int qr, long long val) {
        if (ql <= l && r <= qr) {
            apply(node, l, r, val);
            return;
        }
        push(node, l, r);
        int mid = (l + r) / 2;
        if (ql <= mid) rangeAdd(node * 2, l, mid, ql, qr, val);
        if (qr > mid) rangeAdd(node * 2 + 1, mid + 1, r, ql, qr, val);
        tree[node] = tree[node * 2] + tree[node * 2 + 1];
    }

    long long query(int node, int l, int r, int ql, int qr) {
        if (ql <= l && r <= qr) return tree[node];
        push(node, l, r);
        int mid = (l + r) / 2;
        long long res = 0;
        if (ql <= mid) res += query(node * 2, l, mid, ql, qr);
        if (qr > mid) res += query(node * 2 + 1, mid + 1, r, ql, qr);
        return res;
    }
};
\end{lstlisting}

调用：

\begin{lstlisting}
// Require n >= 1; initial array is all zero.
SegTree seg(n);
// For each initial a[i], call rangeAdd(1,1,n,i,i,a[i]).
seg.rangeAdd(1, 1, n, l, r, v);
long long ans = seg.query(1, 1, n, l, r);
\end{lstlisting}

\subsection{不会线段树时怎么办}

\begin{itemize}
  \item 如果没有修改，用前缀和。
  \item 如果只有区间修改、最后查询，用差分。
  \item 如果是单点修改、区间查询，用树状数组。
  \item 如果数据小，直接暴力扫描。
\end{itemize}

\section{Kruskal 最小生成树入口}\label{sec:08-10}

\subsection{什么题会用到}

\begin{itemize}
  \item 连接所有点。
  \item 总代价最小。
  \item 无向图。
  \item 选择若干边让图连通。
\end{itemize}

\subsection{Kruskal 模板}

\begin{lstlisting}
struct Edge {
    int u, v;
    long long w;
};

sort(edges.begin(), edges.end(), [](const Edge& a, const Edge& b) {
    return a.w < b.w;
});

DSU dsu(n);
long long ans = 0;
int cnt = 0;
for (auto e : edges) {
    if (dsu.unite(e.u, e.v)) {
        ans += e.w;
        cnt++;
    }
}

if (cnt == n - 1) cout << ans << '\n';
else cout << "Impossible\n";
\end{lstlisting}

\section{树上问题入口}\label{sec:08-11}

\subsection{什么题会用到}

\begin{itemize}
  \item 给定一棵树。
  \item 查询两点路径。
  \item 查询祖先、最近公共祖先。
  \item 对树上路径加值或统计经过次数。
\end{itemize}

\subsection{树上暴力拿分}

如果不会 LCA，且 $n,q$ 小，可以每次 DFS 找路径。

\begin{lstlisting}
bool dfsPath(int u, int parent, int target, vector<int>& path) {
    if (u == target) {
        path.push_back(u);
        return true;
    }
    for (int v : g[u]) {
        if (v == parent) continue;
        if (dfsPath(v, u, target, path)) {
            path.push_back(u);
            return true;
        }
    }
    return false;
}
\end{lstlisting}

\subsection{LCA 倍增入口}

\begin{lstlisting}
int LOG; // compute from n before allocating up
vector<vector<int>> up;
vector<int> depth;
vector<vector<int>> g;

void dfsLca(int u, int p) {
    up[0][u] = p;
    for (int k = 1; k < LOG; k++) {
        up[k][u] = up[k - 1][up[k - 1][u]];
    }
    for (int v : g[u]) {
        if (v == p) continue;
        depth[v] = depth[u] + 1;
        dfsLca(v, u);
    }
}

int lca(int a, int b) {
    if (depth[a] < depth[b]) swap(a, b);
    int diff = depth[a] - depth[b];
    for (int k = 0; k < LOG; k++) {
        if (diff & (1 << k)) a = up[k][a];
    }
    if (a == b) return a;
    for (int k = LOG - 1; k >= 0; k--) {
        if (up[k][a] != up[k][b]) {
            a = up[k][a];
            b = up[k][b];
        }
    }
    return up[0][a];
}
\end{lstlisting}

先读入 n 并将无向树边存入 g。下面根为 1，根父亲为自身；深链应改用 \Tref{209} 的无递归遍历。初始化：

\begin{lstlisting}
LOG = 1;
while ((1LL << LOG) <= n) LOG++;
up.assign(LOG, vector<int>(n + 1));
depth.assign(n + 1, 0);
dfsLca(1, 1);
\end{lstlisting}

\subsection{树上差分入口}

多次给路径 $(u,v)$ 加 1，最后统计每个点或边被经过次数。

点差分采用根的父亲为自身的约定，经过根的路径不能再减一次根；以下函数与 LCA 共用 g、up。处理全部路径后调用 collect(1,1)。

\begin{lstlisting}
vector<long long> diff(n + 1, 0);

void addPath(int u, int v) {
    int w = lca(u, v);
    diff[u]++;
    diff[v]++;
    diff[w]--;
    if (up[0][w] != w) diff[up[0][w]]--; // root is its own parent
}

void collect(int u, int p) {
    for (int v : g[u]) {
        if (v == p) continue;
        collect(v, u);
        diff[u] += diff[v];
    }
}
\end{lstlisting}

点差分和边差分的公式不同。拿不准时，可先逐条路径更新，用小数据验证；是否能得部分分要看子任务范围。

\section{状态压缩 DP 入口}\label{sec:08-12}

\subsection{什么题会用到}

\begin{itemize}
  \item $n$ 很小，例如 $n \le 20$。
  \item 每个元素只有选或不选。
  \item 状态可以用二进制集合表示。
  \item 旅行、分组、覆盖、集合最优值。
\end{itemize}

\subsection{集合 DP 基本形状}

\begin{lstlisting}
int N = 1 << n;
vector<int> dp(N, INF);
dp[0] = 0;

for (int mask = 0; mask < N; mask++) {
    for (int i = 0; i < n; i++) {
        if ((mask & (1 << i)) == 0) {
            int nmask = mask | (1 << i);
            dp[nmask] = min(dp[nmask], dp[mask] + cost[i]);
        }
    }
}
\end{lstlisting}

\subsection{状压易错点}

\begin{itemize}
  \item $1 << n$ 当 $n \ge 31$ 会溢出，要用 \Code{1LL << n}。
  \item 状态数是 $2^n$，$n=25$ 已经很大。
  \item 数组大小和内存要提前估算。
\end{itemize}

\section{KMP 字符串匹配入口}\label{sec:08-13}

\subsection{什么题会用到}

\begin{itemize}
  \item 长文本中多次查找模式串。
  \item 字符串长度很大，暴力匹配会超时。
  \item 需要统计模式串出现次数。
\end{itemize}

\subsection{KMP 模板}

\begin{lstlisting}
vector<int> prefixFunction(const string& p) {
    int n = (int)p.size();
    vector<int> pi(n, 0);
    for (int i = 1; i < n; i++) {
        int j = pi[i - 1];
        while (j > 0 && p[i] != p[j]) j = pi[j - 1];
        if (p[i] == p[j]) j++;
        pi[i] = j;
    }
    return pi;
}

int countMatch(const string& text, const string& p) {
    if (p.empty()) return (int)text.size() + 1; // empty pattern: all gaps
    vector<int> pi = prefixFunction(p);
    int j = 0;
    int cnt = 0;
    for (char c : text) {
        while (j > 0 && c != p[j]) j = pi[j - 1];
        if (c == p[j]) j++;
        if (j == (int)p.size()) {
            cnt++;
            j = pi[j - 1];
        }
    }
    return cnt;
}
\end{lstlisting}

\subsection{不会 KMP 时}

\begin{itemize}
  \item 字符串短时用暴力匹配。
  \item C++ 可用 \Code{s.find(t)} 拿部分数据。
  \item 多模式串匹配可能要 Trie 或 AC 自动机，但新手优先拿小数据。
\end{itemize}

\section{Trie 字典树入口}\label{sec:08-14}

\subsection{什么题会用到}

\begin{itemize}
  \item 很多字符串插入和查询。
  \item 前缀查询。
  \item 字典、单词集合。
\end{itemize}

\subsection{小写字母 Trie}

\begin{lstlisting}
struct Trie {
    struct Node {
        int child[26];
        bool end = false;
        Node() {
            memset(child, -1, sizeof child);
        }
    };

    vector<Node> tr;

    Trie() {
        tr.push_back(Node());
    }

    void insert(const string& s) {
        int u = 0;
        for (char c : s) {
            int x = c - 'a';
            if (tr[u].child[x] == -1) {
                tr[u].child[x] = (int)tr.size();
                tr.push_back(Node());
            }
            u = tr[u].child[x];
        }
        tr[u].end = true;
    }

    bool find(const string& s) {
        int u = 0;
        for (char c : s) {
            int x = c - 'a';
            if (tr[u].child[x] == -1) return false;
            u = tr[u].child[x];
        }
        return tr[u].end;
    }
};
\end{lstlisting}

\subsection{不会 Trie 时}

\begin{itemize}
  \item 只查完整字符串，用 \Code{set<string>} 或 \Code{unordered\_set<string>}。
  \item 数据小，用逐个字符串比较。
  \item 前缀查询可以用排序后 \Code{lower\_bound} 尝试部分分。
\end{itemize}

\section{D/E 考场时间策略}\label{sec:08-15}

\begin{longtable}{p{3cm}p{8cm}p{3cm}}
\toprule
剩余时间 & 建议动作 & 目标 \\
\midrule
\endfirsthead
\toprule
剩余时间 & 建议动作 & 目标 \\
\midrule
\endhead
60 分钟以上 & 先思考满分算法，再写部分分保底 & 进一步优化 \\
30--60 分钟 & 先写暴力或特殊性质，再看能否优化 & 稳定拿分 \\
15--30 分钟 & 只写最确定的小数据做法 & 不空题 \\
15 分钟以内 & 完成一个明确子任务的正确实现，清理输出 & 完成可做的部分 \\
\bottomrule
\end{longtable}

\section{常见高级算法入口}\label{sec:08-16}

\begin{longtable}{p{4cm}p{5cm}p{5cm}}
\toprule
题面信号 & 可能算法 & 先翻哪里 \\
\midrule
\endfirsthead
\toprule
题面信号 & 可能算法 & 先翻哪里 \\
\midrule
\endhead
最大化最小值、最小化最大值 & 二分答案 & B 题二分答案 \\
区间修改查询交替 & 线段树 / 树状数组 & \Tref{401} / \Tref{402} \\
树上路径、祖先 & LCA / 树上差分 & \Tref{208}、本章树上差分 \\
强连通、缩点 & Tarjan & 本书未收录 Tarjan；小图可枚举可达性 \\
最小连接代价 & 最小生成树 & Kruskal / Prim \\
网络最短距离 & Dijkstra / Floyd & B 题图论入口 \\
容量和价值 & 背包 DP & B 题 DP 入口 \\
集合状态很小 & 状态压缩 DP & 状压入口 \\
字符串匹配很多次 & KMP / Trie & \Tref{502} / \Tref{503} \\
\bottomrule
\end{longtable}

\section{高级算法识别速查}\label{sec:08-17}

\begin{longtable}{p{4cm}p{5cm}p{5cm}}
\toprule
关键词 & 优先想到 & 不会时的替代 \\
\midrule
\endfirsthead
\toprule
关键词 & 优先想到 & 不会时的替代 \\
\midrule
\endhead
区间加、区间查询 & 线段树 & 差分、树状数组、暴力 \\
单点修改、区间和 & 树状数组 & 前缀和静态版、暴力 \\
树上路径查询 & LCA & 每次 DFS 找路径 \\
多次树上路径加 & 树上差分 & 暴力路径加 \\
连接所有点最小代价 & Kruskal & 枚举小图、Prim 小数据 \\
集合状态很小 & 状压 DP & DFS 枚举 \\
长串匹配短串 & KMP & \Code{find} / 暴力 \\
前缀查询很多 & Trie & set / 排序后二分 \\
最大化最小值 & 二分答案 & 枚举答案小范围 \\
任务依赖 & 拓扑排序 & 小图 DFS 检查 \\
\bottomrule
\end{longtable}

\section{D/E 模板练习建议}\label{sec:08-18}

\begin{longtable}{p{4cm}p{4cm}p{6cm}}
\toprule
模板 & 优先级 & 原因 \\
\midrule
\endfirsthead
\toprule
模板 & 优先级 & 原因 \\
\midrule
\endhead
二分答案 & 高 & B/D 都常见，代码短，提分明显。 \\
树状数组 & 高 & 区间题常见，比线段树简单。 \\
Dijkstra & 高 & 图论常见，B/D 都可能用。 \\
Kruskal & 中 & 最小生成树题专用，模板短。 \\
线段树 & 中高 & 强但代码长，建议练熟基础版。 \\
LCA & 中 & 树上题常见，但新手可先会暴力。 \\
状态压缩 DP & 中 & 看到 $n \le 20$ 很有用。 \\
KMP & 中低 & 字符串题有用，但 CSP 中不一定高频。 \\
Trie & 中低 & 前缀题有用，完整字符串可用 set 替代。 \\
树上差分 & 低到中 & 需要 LCA 基础，适合进一步优化。 \\
\bottomrule
\end{longtable}

\section{完整做法暂时不会时}\label{sec:08-19}

出现以下情况，应转为部分分：

\begin{itemize}
  \item 已经想了较久，仍没有可实现的做法。
  \item 会算法名字，但模板完全不会写。
  \item 写到一半发现状态定义不清。
  \item 样例都无法解释。
  \item A/B/C 还没稳。
\end{itemize}

\section{D/E 题提交前检查}\label{sec:08-20}

\begin{itemize}
  \item 是否至少能过样例。
  \item 是否能过自己造的小数据。
  \item 是否处理了最小输入。
  \item 如果是暴力，是否明确知道能拿哪些数据。
  \item 是否删除调试输出。
  \item 是否没有因为追求优化改坏暴力基础逻辑。
\end{itemize}


\part{模板与检查}
\chapter{高频算法模板库}

\section{模板库使用方法}

本章把前面章节中反复出现的算法模板统一编号。考场上建议先通过决策树确定题型，再来本章找对应模板。

\begin{itemize}
  \item \Code{T001--T099}：基础实现和 STL。
  \item \Code{T100--T199}：前缀和、差分、排序、哈希、二分。
  \item \Code{T200--T299}：搜索和图论。
  \item \Code{T300--T399}：动态规划。
  \item \Code{T400--T499}：数据结构。
  \item \Code{T500--T599}：字符串、数学、日期。
\end{itemize}

\section{模板速查表}

\begin{longtable}{p{2cm}p{4cm}p{5cm}p{3cm}}
\toprule
编号 & 名称 & 识别信号 & 复杂度 \\
\midrule
T101 & 一维前缀和 & 多次静态区间和查询 & 预处理 $O(n)$，查询 $O(1)$ \\
T102 & 二维前缀和 & 多次子矩阵和查询 & 预处理 $O(nm)$，查询 $O(1)$ \\
T103 & 一维差分 & 多次区间加，最后查询 & $O(n+q)$ \\
T104 & 二维差分 & 多次矩形加，最后输出 & $O(nm+q)$ \\
T105 & 排序自定义比较器 & 多关键字排序、排名 & $O(n\log n)$ \\
T106 & 哈希统计 & 出现次数、查重、频率 & 平均 $O(n)$ \\
T107 & 二分答案 & 最大化最小值、最小化最大值 & $O(\log V \cdot check)$ \\
T108 & 双指针 & 排序后找对、滑动窗口 & $O(n)$ 或 $O(n\log n)$ \\
T109 & 坐标离散化 & 坐标大但点数少 & $O(n\log n)$ \\
T110 & 事件排序 & 按时间发生、扫描线 & $O(n\log n)$ \\
T201 & 网格 BFS & 最短步数、迷宫 & $O(nm)$ \\
T202 & DFS 连通块 & 区域大小、连通块数量 & $O(n+m)$ 或 $O(nm)$ \\
T203 & 并查集 & 合并集合、判断同组 & 近似 $O(1)$ \\
T204 & Dijkstra & 正权最短路 & $O((n+m)\log n)$ \\
T205 & 拓扑排序 & 依赖关系、先后顺序 & $O(n+m)$ \\
T206 & Floyd & 小图任意两点最短路 & $O(n^3)$ \\
T207 & Kruskal & 最小生成树 & $O(m\log m)$ \\
T208 & LCA 倍增 & 树上最近公共祖先 & $O(n\log n)$ 预处理，$O(\log n)$ 查询 \\
T301 & 01 背包 & 每个物品选或不选 & $O(nV)$ \\
T302 & 完全背包 & 每个物品可选多次 & $O(nV)$ \\
T401 & 树状数组 & 单点修改、区间和 & $O(\log n)$ \\
T402 & 线段树 & 区间修改、区间查询 & $O(\log n)$ \\
T403 & 单调栈 & 最近更大/更小元素 & $O(n)$ \\
T404 & 单调队列 & 滑动窗口最大/最小 & $O(n)$ \\
T501 & split 字符串分割 & 路径、配置、命令解析 & $O(|s|)$ \\
T502 & KMP & 长文本匹配模式串 & $O(n+m)$ \\
T503 & Trie & 字符串集合、前缀查询 & $O(\sum |s|)$ \\
T510 & 整除取整和标准取模 & 分组数、负数取模、数学 floor/ceil & $O(1)$ \\
T511 & 快速幂 & $a^b \bmod m$ & $O(\log b)$ \\
T512 & 素数和约数 & 判断素数、枚举约数 & $O(\sqrt n)$ \\
T513 & 进制转换 & 任意进制与十进制 & $O(\text{位数})$ \\
T514 & 日期处理 & 闰年、每月天数、日期推进 & 按实现而定 \\
T601 & 大模拟骨架 & 多对象、多命令、长题面 & 按题目而定 \\
T602 & 命令分发 & 每行一个操作名 & 按命令复杂度 \\
T603 & 状态机 & 对象多状态切换 & 按操作复杂度 \\
T604 & 懒删除优先队列 & 动态删除/更新后仍取最优 & $O(\log n)$ 均摊 \\
T605 & LRU 缓存 & 最近最少使用淘汰 & $O(1)$ 均摊 \\
T606 & 路径规范化 & \Code{/a/../b}、目录路径 & $O(|s|)$ \\
T607 & URL 参数解析 & \Code{path?a=1\&b=2} & $O(|s|)$ \\
T608 & 括号匹配 & 括号合法性、嵌套结构 & $O(n)$ \\
T609 & 简单表达式求值 & 只含加减的表达式 & $O(n)$ \\
\bottomrule
\end{longtable}

\section{T101 一维前缀和}

\subsection{适用场景}

\begin{itemize}
  \item 数组不修改。
  \item 多次查询 $[l,r]$ 的和。
  \item 需要快速得到一段累计值。
\end{itemize}

\subsection{模板}

\begin{lstlisting}
int n, q;
cin >> n >> q;
vector<long long> pre(n + 1, 0);

for (int i = 1; i <= n; i++) {
    long long x;
    cin >> x;
    pre[i] = pre[i - 1] + x;
}

while (q--) {
    int l, r;
    cin >> l >> r;
    cout << pre[r] - pre[l - 1] << '\n';
}
\end{lstlisting}

\subsection{易错点}

\begin{itemize}
  \item 模板默认下标从 1 开始。
  \item \Code{pre[0]} 是 0。
  \item 求和用 \Code{long long}。
  \item 如果有修改，普通前缀和不适用。
\end{itemize}

\section{T102 二维前缀和}

\subsection{适用场景}

多次查询子矩阵和。

\begin{lstlisting}
int n, m, q;
cin >> n >> m >> q;
vector<vector<long long>> pre(n + 1, vector<long long>(m + 1, 0));

for (int i = 1; i <= n; i++) {
    for (int j = 1; j <= m; j++) {
        long long x;
        cin >> x;
        pre[i][j] = pre[i - 1][j] + pre[i][j - 1]
                  - pre[i - 1][j - 1] + x;
    }
}

while (q--) {
    int x1, y1, x2, y2;
    cin >> x1 >> y1 >> x2 >> y2;
    long long ans = pre[x2][y2] - pre[x1 - 1][y2]
                  - pre[x2][y1 - 1] + pre[x1 - 1][y1 - 1];
    cout << ans << '\n';
}
\end{lstlisting}

\subsection{易错点}

\begin{itemize}
  \item 四项公式不要写反。
  \item 默认坐标从 1 开始。
  \item 子矩阵端点是否保证 $x1 \le x2, y1 \le y2$。
\end{itemize}

\section{T103 一维差分}

\subsection{适用场景}

多次区间加，最后统一输出每个位置的值。

\begin{lstlisting}
int n, q;
cin >> n >> q;
vector<long long> diff(n + 2, 0);

while (q--) {
    int l, r;
    long long v;
    cin >> l >> r >> v;
    diff[l] += v;
    diff[r + 1] -= v;
}

long long cur = 0;
for (int i = 1; i <= n; i++) {
    cur += diff[i];
    cout << cur << '\n';
}
\end{lstlisting}

\subsection{易错点}

\begin{itemize}
  \item 数组开到 \Code{n+2}。
  \item 如果修改后立刻查询区间和，普通差分不够。
  \item 修改值用 \Code{long long}。
\end{itemize}

\section{T104 二维差分}

\subsection{适用场景}

多次对子矩形加值，最后输出整个矩阵。

\begin{lstlisting}
int n, m, q;
cin >> n >> m >> q;
vector<vector<long long>> diff(n + 2, vector<long long>(m + 2, 0));

while (q--) {
    int x1, y1, x2, y2;
    long long v;
    cin >> x1 >> y1 >> x2 >> y2 >> v;
    diff[x1][y1] += v;
    diff[x2 + 1][y1] -= v;
    diff[x1][y2 + 1] -= v;
    diff[x2 + 1][y2 + 1] += v;
}

for (int i = 1; i <= n; i++) {
    for (int j = 1; j <= m; j++) {
        diff[i][j] += diff[i - 1][j] + diff[i][j - 1]
                    - diff[i - 1][j - 1];
    }
}
\end{lstlisting}

\section{T105 排序自定义比较器}

\subsection{适用场景}

多关键字排序、排名、优先级排序。

\begin{lstlisting}
struct Node {
    int id;
    int score;
    int time;
};

sort(a.begin(), a.end(), [](const Node& x, const Node& y) {
    if (x.score != y.score) return x.score > y.score;
    if (x.time != y.time) return x.time < y.time;
    return x.id < y.id;
});
\end{lstlisting}

\subsection{易错点}

\begin{itemize}
  \item 相等时不能返回 \Code{true}。
  \item 每个关键字升序还是降序要写清楚。
  \item 需要保持原顺序时用 \Code{stable\_sort}。
\end{itemize}

\section{T106 哈希统计}

\subsection{值域小}

\begin{lstlisting}
vector<int> cnt(100001, 0);
for (int x : a) {
    cnt[x]++;
}
\end{lstlisting}

\subsection{值域大或字符串}

\begin{lstlisting}
unordered_map<string, int> cnt;
for (string s : words) {
    cnt[s]++;
}
\end{lstlisting}

\subsection{易错点}

\begin{itemize}
  \item 需要有序输出时用 \Code{map}。
  \item key 很大不能开计数数组。
  \item \Code{cnt[x]} 会自动创建 key。
\end{itemize}

\section{T107 二分答案}

\subsection{适用场景}

\begin{itemize}
  \item 求最小可行值。
  \item 求最大可行值。
  \item 给定答案 $mid$ 能判断是否可行。
  \item 可行性具有单调性。
\end{itemize}

\subsection{最小可行值}

\begin{lstlisting}
bool check(long long mid) {
    // return true if mid is feasible
}

long long l = 0, r = 1000000000000000000LL;
while (l < r) {
    long long mid = l + (r - l) / 2;
    if (check(mid)) r = mid;
    else l = mid + 1;
}
cout << l << '\n';
\end{lstlisting}

\subsection{最大可行值}

\begin{lstlisting}
long long l = 0, r = 1000000000000000000LL;
while (l < r) {
    long long mid = l + (r - l + 1) / 2;
    if (check(mid)) l = mid;
    else r = mid - 1;
}
cout << l << '\n';
\end{lstlisting}

\Warn{二分答案必须有单调性。}

\section{T108 双指针}

\subsection{排序后找两数和}

\begin{lstlisting}
sort(a.begin(), a.end());
int l = 0, r = n - 1;
bool ok = false;

while (l < r) {
    long long sum = 1LL * a[l] + a[r];
    if (sum == target) {
        ok = true;
        break;
    } else if (sum < target) {
        l++;
    } else {
        r--;
    }
}
\end{lstlisting}

\subsection{非负数组滑动窗口}

\begin{lstlisting}
int l = 0;
long long sum = 0;
int best = 0;

for (int r = 0; r < n; r++) {
    sum += a[r];
    while (sum > limit && l <= r) {
        sum -= a[l];
        l++;
    }
    best = max(best, r - l + 1);
}
\end{lstlisting}

\Warn{滑动窗口通常要求元素非负或条件具有单调性。}

\section{T109 坐标离散化}

\subsection{适用场景}

坐标值很大，但出现的坐标数量不多。

\begin{lstlisting}
vector<int> xs;
for (int x : original) {
    xs.push_back(x);
}

sort(xs.begin(), xs.end());
xs.erase(unique(xs.begin(), xs.end()), xs.end());

auto getId = [&](int x) {
    return (int)(lower_bound(xs.begin(), xs.end(), x) - xs.begin()) + 1;
};
\end{lstlisting}

\Warn{如果题目涉及区间长度，不能只看离散编号，还要保留真实坐标差。}

\section{T110 事件排序}

\subsection{适用场景}

按时间顺序处理事件，或扫描区间开始/结束。

\begin{lstlisting}
struct Event {
    int time;
    int type;
    int id;
};

sort(events.begin(), events.end(), [](const Event& a, const Event& b) {
    if (a.time != b.time) return a.time < b.time;
    return a.type < b.type;
});

for (const Event& e : events) {
    if (e.type == 0) {
        // start
    } else {
        // end
    }
}
\end{lstlisting}

\subsection{区间同时进行数量}

\begin{lstlisting}
vector<pair<int, int>> events;
for (auto [l, r] : intervals) {
    events.push_back({l, 1});
    events.push_back({r, -1});
}

sort(events.begin(), events.end(), [](auto a, auto b) {
    if (a.first != b.first) return a.first < b.first;
    return a.second < b.second; // end before start for [l,r)
});

int cur = 0, best = 0;
for (auto [t, delta] : events) {
    cur += delta;
    best = max(best, cur);
}
\end{lstlisting}

\Warn{同一时间开始和结束谁先处理，要按题意决定。}

\section{T201 网格 BFS}

\subsection{适用场景}

迷宫、最短步数、每一步代价相同。

\begin{lstlisting}
int n, m;
cin >> n >> m;
vector<string> g(n);
for (int i = 0; i < n; i++) cin >> g[i];

int sx, sy, tx, ty;
cin >> sx >> sy >> tx >> ty;

vector<vector<int>> dist(n, vector<int>(m, -1));
queue<pair<int, int>> q;

dist[sx][sy] = 0;
q.push({sx, sy});

int dx[4] = {-1, 0, 1, 0};
int dy[4] = {0, 1, 0, -1};

while (!q.empty()) {
    auto [x, y] = q.front();
    q.pop();

    for (int d = 0; d < 4; d++) {
        int nx = x + dx[d];
        int ny = y + dy[d];
        if (nx < 0 || nx >= n || ny < 0 || ny >= m) continue;
        if (g[nx][ny] == '#') continue;
        if (dist[nx][ny] != -1) continue;
        dist[nx][ny] = dist[x][y] + 1;
        q.push({nx, ny});
    }
}

cout << dist[tx][ty] << '\n';
\end{lstlisting}

\subsection{易错点}

\begin{itemize}
  \item 坐标是 0 下标还是 1 下标。
  \item 入队时标记访问。
  \item 边权不全相等时不要用普通 BFS。
\end{itemize}

\section{T202 DFS 连通块}

\subsection{网格 DFS}

\begin{lstlisting}
int n, m;
vector<string> g;
vector<vector<int>> vis;
int dx[4] = {-1, 0, 1, 0};
int dy[4] = {0, 1, 0, -1};

int dfs(int x, int y) {
    vis[x][y] = 1;
    int area = 1;
    for (int d = 0; d < 4; d++) {
        int nx = x + dx[d];
        int ny = y + dy[d];
        if (nx < 0 || nx >= n || ny < 0 || ny >= m) continue;
        if (vis[nx][ny]) continue;
        if (g[nx][ny] == '#') continue;
        area += dfs(nx, ny);
    }
    return area;
}
\end{lstlisting}

\subsection{易错点}

\begin{itemize}
  \item 大网格递归可能爆栈，可改 BFS。
  \item 障碍字符要按题目改。
\end{itemize}

\section{T203 并查集}

\begin{lstlisting}
struct DSU {
    vector<int> parent, sz;

    DSU(int n) {
        parent.resize(n + 1);
        sz.assign(n + 1, 1);
        for (int i = 1; i <= n; i++) parent[i] = i;
    }

    int find(int x) {
        if (parent[x] == x) return x;
        return parent[x] = find(parent[x]);
    }

    bool unite(int a, int b) {
        int ra = find(a);
        int rb = find(b);
        if (ra == rb) return false;
        if (sz[ra] < sz[rb]) swap(ra, rb);
        parent[rb] = ra;
        sz[ra] += sz[rb];
        return true;
    }

    bool same(int a, int b) {
        return find(a) == find(b);
    }
};
\end{lstlisting}

\subsection{适用场景}

合并集合、判断是否同组、连通块。

\section{T204 Dijkstra}

\begin{lstlisting}
struct Edge {
    int to;
    long long w;
};

const long long INF = (long long)4e18;

int n, m, s;
cin >> n >> m >> s;
vector<vector<Edge>> g(n + 1);

for (int i = 0; i < m; i++) {
    int u, v;
    long long w;
    cin >> u >> v >> w;
    g[u].push_back({v, w});
    g[v].push_back({u, w}); // remove if directed
}

vector<long long> dist(n + 1, INF);
priority_queue<pair<long long, int>,
               vector<pair<long long, int>>,
               greater<pair<long long, int>>> pq;

dist[s] = 0;
pq.push({0, s});

while (!pq.empty()) {
    auto [d, u] = pq.top();
    pq.pop();
    if (d != dist[u]) continue;

    for (auto e : g[u]) {
        if (dist[e.to] > dist[u] + e.w) {
            dist[e.to] = dist[u] + e.w;
            pq.push({dist[e.to], e.to});
        }
    }
}
\end{lstlisting}

\subsection{易错点}

\begin{itemize}
  \item 不能处理负边权。
  \item 无向图双向加边，有向图单向加边。
  \item 距离用 \Code{long long}。
\end{itemize}

\section{T205 拓扑排序}

\subsection{适用场景}

任务依赖、课程先修、判断有向图是否有环。

\begin{lstlisting}
int n, m;
cin >> n >> m;
vector<vector<int>> g(n + 1);
vector<int> indeg(n + 1, 0);

for (int i = 0; i < m; i++) {
    int u, v;
    cin >> u >> v; // u before v
    g[u].push_back(v);
    indeg[v]++;
}

queue<int> q;
for (int i = 1; i <= n; i++) {
    if (indeg[i] == 0) q.push(i);
}

vector<int> order;
while (!q.empty()) {
    int u = q.front();
    q.pop();
    order.push_back(u);
    for (int v : g[u]) {
        indeg[v]--;
        if (indeg[v] == 0) q.push(v);
    }
}

if ((int)order.size() == n) {
    // has topological order
} else {
    // has cycle
}
\end{lstlisting}

\section{T206 Floyd}

\subsection{适用场景}

点数小，任意两点最短路。

\begin{lstlisting}
const long long INF = (long long)4e18;

int n, m;
cin >> n >> m;
vector<vector<long long>> dist(n + 1, vector<long long>(n + 1, INF));

for (int i = 1; i <= n; i++) dist[i][i] = 0;

for (int i = 0; i < m; i++) {
    int u, v;
    long long w;
    cin >> u >> v >> w;
    dist[u][v] = min(dist[u][v], w);
    dist[v][u] = min(dist[v][u], w); // remove if directed
}

for (int k = 1; k <= n; k++) {
    for (int i = 1; i <= n; i++) {
        for (int j = 1; j <= n; j++) {
            if (dist[i][k] == INF || dist[k][j] == INF) continue;
            dist[i][j] = min(dist[i][j], dist[i][k] + dist[k][j]);
        }
    }
}
\end{lstlisting}

\Warn{三层循环顺序必须是 \Code{k,i,j}。}

\section{T207 Kruskal}

\subsection{适用场景}

无向图，连接所有点的最小总代价。

\begin{lstlisting}
struct Edge {
    int u, v;
    long long w;
};

vector<Edge> edges;
sort(edges.begin(), edges.end(), [](const Edge& a, const Edge& b) {
    return a.w < b.w;
});

DSU dsu(n);
long long ans = 0;
int cnt = 0;

for (auto e : edges) {
    if (dsu.unite(e.u, e.v)) {
        ans += e.w;
        cnt++;
    }
}

if (cnt == n - 1) cout << ans << '\n';
else cout << "Impossible\n";
\end{lstlisting}

\section{T208 LCA 倍增}

\subsection{适用场景}

树上最近公共祖先、树上路径查询。

\begin{lstlisting}
const int LOG = 20;
vector<vector<int>> up;
vector<int> depth;
vector<vector<int>> g;

void dfsLca(int u, int p) {
    up[0][u] = p;
    for (int k = 1; k < LOG; k++) {
        up[k][u] = up[k - 1][up[k - 1][u]];
    }
    for (int v : g[u]) {
        if (v == p) continue;
        depth[v] = depth[u] + 1;
        dfsLca(v, u);
    }
}

int lca(int a, int b) {
    if (depth[a] < depth[b]) swap(a, b);
    int diff = depth[a] - depth[b];
    for (int k = 0; k < LOG; k++) {
        if (diff & (1 << k)) a = up[k][a];
    }
    if (a == b) return a;
    for (int k = LOG - 1; k >= 0; k--) {
        if (up[k][a] != up[k][b]) {
            a = up[k][a];
            b = up[k][b];
        }
    }
    return up[0][a];
}
\end{lstlisting}

初始化：

\begin{lstlisting}
up.assign(LOG, vector<int>(n + 1));
depth.assign(n + 1, 0);
dfsLca(1, 1);
\end{lstlisting}

\section{T301 01 背包}

\begin{lstlisting}
int n, V;
cin >> n >> V;
vector<int> w(n + 1), val(n + 1);
for (int i = 1; i <= n; i++) {
    cin >> w[i] >> val[i];
}

vector<int> dp(V + 1, 0);
for (int i = 1; i <= n; i++) {
    for (int j = V; j >= w[i]; j--) {
        dp[j] = max(dp[j], dp[j - w[i]] + val[i]);
    }
}
cout << dp[V] << '\n';
\end{lstlisting}

\Warn{01 背包容量倒序。}

\section{T302 完全背包}

\begin{lstlisting}
int n, V;
cin >> n >> V;
vector<int> w(n + 1), val(n + 1);
for (int i = 1; i <= n; i++) {
    cin >> w[i] >> val[i];
}

vector<int> dp(V + 1, 0);
for (int i = 1; i <= n; i++) {
    for (int j = w[i]; j <= V; j++) {
        dp[j] = max(dp[j], dp[j - w[i]] + val[i]);
    }
}
cout << dp[V] << '\n';
\end{lstlisting}

\Warn{完全背包容量正序。}

\section{T401 树状数组}

\begin{lstlisting}
struct BIT {
    int n;
    vector<long long> tree;

    BIT(int n) : n(n), tree(n + 1, 0) {}

    void add(int idx, long long val) {
        for (; idx <= n; idx += idx & -idx) {
            tree[idx] += val;
        }
    }

    long long sumPrefix(int idx) {
        long long res = 0;
        for (; idx > 0; idx -= idx & -idx) {
            res += tree[idx];
        }
        return res;
    }

    long long rangeSum(int l, int r) {
        return sumPrefix(r) - sumPrefix(l - 1);
    }
};
\end{lstlisting}

\subsection{适用场景}

单点修改、区间和查询。

\section{T402 线段树}

\subsection{适用场景}

区间加、区间和查询。模板较长，只有练过才建议考场使用。

\begin{lstlisting}
struct SegTree {
    int n;
    vector<long long> tree, lazy;

    SegTree(int n) : n(n), tree(4 * n + 4), lazy(4 * n + 4) {}

    void apply(int node, int l, int r, long long val) {
        tree[node] += val * (r - l + 1);
        lazy[node] += val;
    }

    void push(int node, int l, int r) {
        if (lazy[node] == 0 || l == r) return;
        int mid = (l + r) / 2;
        apply(node * 2, l, mid, lazy[node]);
        apply(node * 2 + 1, mid + 1, r, lazy[node]);
        lazy[node] = 0;
    }

    void rangeAdd(int node, int l, int r, int ql, int qr, long long val) {
        if (ql <= l && r <= qr) {
            apply(node, l, r, val);
            return;
        }
        push(node, l, r);
        int mid = (l + r) / 2;
        if (ql <= mid) rangeAdd(node * 2, l, mid, ql, qr, val);
        if (qr > mid) rangeAdd(node * 2 + 1, mid + 1, r, ql, qr, val);
        tree[node] = tree[node * 2] + tree[node * 2 + 1];
    }

    long long query(int node, int l, int r, int ql, int qr) {
        if (ql <= l && r <= qr) return tree[node];
        push(node, l, r);
        int mid = (l + r) / 2;
        long long res = 0;
        if (ql <= mid) res += query(node * 2, l, mid, ql, qr);
        if (qr > mid) res += query(node * 2 + 1, mid + 1, r, ql, qr);
        return res;
    }
};
\end{lstlisting}

调用：

\begin{lstlisting}
SegTree seg(n);
seg.rangeAdd(1, 1, n, l, r, v);
long long ans = seg.query(1, 1, n, l, r);
\end{lstlisting}

\subsection{易错点}

\begin{itemize}
  \item 默认区间是 1 到 $n$。
  \item 懒标记下传不要漏。
  \item 求和用 \Code{long long}。
\end{itemize}

\section{T403 单调栈}

\subsection{适用场景}

找每个元素右边第一个更大元素、左边第一个更小元素等。

\begin{lstlisting}
vector<int> ans(n, -1);
stack<int> st; // indices

for (int i = 0; i < n; i++) {
    while (!st.empty() && a[i] > a[st.top()]) {
        ans[st.top()] = i;
        st.pop();
    }
    st.push(i);
}
\end{lstlisting}

\subsection{易错点}

\begin{itemize}
  \item 比较符号决定严格更大还是大于等于。
  \item 栈里通常存下标，不是值。
  \item 全相等数据要测试。
\end{itemize}

\section{T404 单调队列}

\subsection{适用场景}

固定长度滑动窗口最大值或最小值。

\begin{lstlisting}
deque<int> dq; // indices, values decreasing
for (int i = 0; i < n; i++) {
    while (!dq.empty() && dq.front() <= i - k) dq.pop_front();
    while (!dq.empty() && a[dq.back()] <= a[i]) dq.pop_back();
    dq.push_back(i);

    if (i >= k - 1) {
        cout << a[dq.front()] << '\n';
    }
}
\end{lstlisting}

\subsection{易错点}

\begin{itemize}
  \item 先删除过期下标。
  \item 最大值队列保持递减，最小值队列保持递增。
  \item 输出窗口从 \Code{i >= k - 1} 开始。
\end{itemize}

\section{T501 split 字符串分割}

\subsection{按指定字符分割}

\begin{lstlisting}
vector<string> split(const string& s, char sep) {
    vector<string> res;
    string cur;
    for (char c : s) {
        if (c == sep) {
            res.push_back(cur);
            cur.clear();
        } else {
            cur.push_back(c);
        }
    }
    res.push_back(cur);
    return res;
}
\end{lstlisting}

\subsection{按空格分割}

\begin{lstlisting}
stringstream ss(line);
string word;
while (ss >> word) {
    // use word
}
\end{lstlisting}

\section{T502 KMP}

\begin{lstlisting}
vector<int> prefixFunction(const string& p) {
    int n = (int)p.size();
    vector<int> pi(n, 0);
    for (int i = 1; i < n; i++) {
        int j = pi[i - 1];
        while (j > 0 && p[i] != p[j]) j = pi[j - 1];
        if (p[i] == p[j]) j++;
        pi[i] = j;
    }
    return pi;
}

int countMatch(const string& text, const string& p) {
    vector<int> pi = prefixFunction(p);
    int j = 0, cnt = 0;
    for (char c : text) {
        while (j > 0 && c != p[j]) j = pi[j - 1];
        if (c == p[j]) j++;
        if (j == (int)p.size()) {
            cnt++;
            j = pi[j - 1];
        }
    }
    return cnt;
}
\end{lstlisting}

\section{T503 Trie}

\begin{lstlisting}
struct Trie {
    struct Node {
        int child[26];
        bool end = false;
        Node() {
            memset(child, -1, sizeof child);
        }
    };

    vector<Node> tr;

    Trie() {
        tr.push_back(Node());
    }

    void insert(const string& s) {
        int u = 0;
        for (char c : s) {
            int x = c - 'a';
            if (tr[u].child[x] == -1) {
                tr[u].child[x] = (int)tr.size();
                tr.push_back(Node());
            }
            u = tr[u].child[x];
        }
        tr[u].end = true;
    }

    bool find(const string& s) {
        int u = 0;
        for (char c : s) {
            int x = c - 'a';
            if (tr[u].child[x] == -1) return false;
            u = tr[u].child[x];
        }
        return tr[u].end;
    }
};
\end{lstlisting}

\Warn{这个 Trie 模板只适合小写英文字母。}

\section{T510 整除取整和标准取模}

\subsection{适用场景}

\begin{itemize}
  \item 分组、分页、分块，要求正整数向上取整。
  \item 坐标、时间差、偏移量可能为负，要求数学意义的向上或向下取整。
  \item 同余、周期、哈希下标中需要把负数余数转成 $0$ 到 $m-1$。
\end{itemize}

\subsection{正整数向上取整}

若 $a \ge 0, b > 0$：

\begin{lstlisting}
long long ceil_div_positive(long long a, long long b) {
    return (a + b - 1) / b;
}
\end{lstlisting}

示例：

\begin{lstlisting}
long long pages = ceil_div_positive(n, pageSize);
long long blocks = ceil_div_positive(n, blockSize);
\end{lstlisting}

\subsection{有符号通用版本}

\begin{lstlisting}
long long floor_div(long long a, long long b) {
    long long q = a / b;
    long long r = a % b;
    if (r != 0 && ((r > 0) != (b > 0))) q--;
    return q;
}

long long ceil_div(long long a, long long b) {
    long long q = a / b;
    long long r = a % b;
    if (r != 0 && ((r > 0) == (b > 0))) q++;
    return q;
}
\end{lstlisting}

\subsection{标准非负余数}

\begin{lstlisting}
long long norm_mod(long long x, long long m) {
    long long r = x % m;
    if (r < 0) r += m;
    return r;
}
\end{lstlisting}

\subsection{易错点}

\begin{itemize}
  \item \Code{(a+b-1)/b} 只适合 $a \ge 0, b > 0$，不是通用向上取整。
  \item C++ 整数除法向 0 截断，\Code{-3 / 2} 是 \Code{-1}。
  \item C++ 负数取模可能得到负数，\Code{-1 \% 5} 是 \Code{-1}。
  \item 除数 \Code{b} 或模数 \Code{m} 不能为 0。
  \item 若 \Code{a + b - 1} 可能溢出，需要改写或先确认范围。
\end{itemize}

\section{T511 快速幂}

\subsection{适用场景}

计算 $a^b \bmod mod$。

\begin{lstlisting}
long long modPow(long long a, long long b, long long mod) {
    long long res = 1 % mod;
    a %= mod;
    while (b > 0) {
        if (b & 1) res = res * a % mod;
        a = a * a % mod;
        b >>= 1;
    }
    return res;
}
\end{lstlisting}

\section{T512 素数和约数}

\subsection{gcd 和 lcm}

\begin{lstlisting}
long long g = gcd(a, b);
long long l = a / g * b;
\end{lstlisting}

\subsection{判断素数}

\begin{lstlisting}
bool isPrime(long long x) {
    if (x < 2) return false;
    for (long long d = 2; d * d <= x; d++) {
        if (x % d == 0) return false;
    }
    return true;
}
\end{lstlisting}

\subsection{枚举约数}

\begin{lstlisting}
vector<long long> divisors;
for (long long d = 1; d * d <= x; d++) {
    if (x % d == 0) {
        divisors.push_back(d);
        if (d != x / d) divisors.push_back(x / d);
    }
}
sort(divisors.begin(), divisors.end());
\end{lstlisting}

\section{T513 进制转换}

\subsection{任意进制转十进制}

\begin{lstlisting}
long long toDecimal(const string& s, int base) {
    long long ans = 0;
    for (char c : s) {
        int v;
        if ('0' <= c && c <= '9') v = c - '0';
        else if ('A' <= c && c <= 'F') v = c - 'A' + 10;
        else v = c - 'a' + 10;
        ans = ans * base + v;
    }
    return ans;
}
\end{lstlisting}

\subsection{十进制转任意进制}

\begin{lstlisting}
string fromDecimal(long long x, int base) {
    if (x == 0) return "0";
    string digits = "0123456789ABCDEF";
    string res;
    while (x > 0) {
        res.push_back(digits[x % base]);
        x /= base;
    }
    reverse(res.begin(), res.end());
    return res;
}
\end{lstlisting}

\Warn{数字很长时不能转成 \Code{long long}，要用字符串模拟。}

\section{T514 日期处理}

\subsection{闰年}

\begin{lstlisting}
bool leap(int y) {
    return (y % 400 == 0) || (y % 4 == 0 && y % 100 != 0);
}
\end{lstlisting}

\subsection{每月天数}

\begin{lstlisting}
int mdays(int y, int m) {
    int d[] = {0,31,28,31,30,31,30,31,31,30,31,30,31};
    if (m == 2 && leap(y)) return 29;
    return d[m];
}
\end{lstlisting}

\subsection{下一天}

\begin{lstlisting}
void nextDay(int& y, int& m, int& d) {
    d++;
    if (d > mdays(y, m)) {
        d = 1;
        m++;
        if (m > 12) {
            m = 1;
            y++;
        }
    }
}
\end{lstlisting}

\subsection{时间转秒}

\begin{lstlisting}
int toSeconds(int h, int m, int s) {
    return h * 3600 + m * 60 + s;
}
\end{lstlisting}

\section{T601 大模拟骨架}

\subsection{适用场景}

C 题长题面、多对象、多命令、多状态。

\begin{lstlisting}
struct Item {
    int id;
    string name;
    int state;
};

int n, q;
vector<Item> items;

void readInput() {
    cin >> n >> q;
    items.resize(n);
    for (int i = 0; i < n; i++) {
        cin >> items[i].id >> items[i].name;
        items[i].state = 0;
    }
}

void solve() {
    while (q--) {
        string op;
        cin >> op;
        // dispatch command
    }
}

int main() {
    ios::sync_with_stdio(false);
    cin.tie(nullptr);

    readInput();
    solve();
    return 0;
}
\end{lstlisting}

\section{T602 命令分发}

\begin{lstlisting}
void handleAdd() {
    int id, x;
    cin >> id >> x;
    // add logic
}

void handleDelete() {
    int id;
    cin >> id;
    // delete logic
}

void handleQuery() {
    int id;
    cin >> id;
    // query logic
}

void handleCommand() {
    string op;
    cin >> op;
    if (op == "add") handleAdd();
    else if (op == "delete") handleDelete();
    else if (op == "query") handleQuery();
}
\end{lstlisting}

\subsection{易错点}

\begin{itemize}
  \item 每种命令的参数个数不同。
  \item 命令参数可能包含空格时要用 \Code{getline}。
  \item 查询命令通常直接输出，修改命令通常不输出。
\end{itemize}

\section{T603 状态机}

\begin{lstlisting}
const int WAITING = 0;
const int RUNNING = 1;
const int FINISHED = 2;

struct Job {
    int id;
    int status;
    int score;
};

void start(Job& job) {
    if (job.status != WAITING) return;
    job.status = RUNNING;
}

void finish(Job& job) {
    if (job.status != RUNNING) return;
    job.status = FINISHED;
    job.score += 10;
}
\end{lstlisting}

\subsection{检查点}

\begin{itemize}
  \item 初始状态是什么。
  \item 每个操作允许在哪些状态发生。
  \item 非法状态转换是忽略还是输出错误。
  \item 状态变化时是否要同步修改其他字段。
\end{itemize}

\section{T604 懒删除优先队列}

\subsection{适用场景}

动态更新或删除元素，同时需要不断取当前最优元素。

\begin{lstlisting}
struct Task {
    int priority;
    int id;
};

struct Cmp {
    bool operator()(const Task& a, const Task& b) const {
        if (a.priority != b.priority) return a.priority < b.priority;
        return a.id > b.id;
    }
};

priority_queue<Task, vector<Task>, Cmp> pq;
map<int, int> currentPriority;

void addOrUpdate(int id, int priority) {
    currentPriority[id] = priority;
    pq.push({priority, id});
}

void removeTask(int id) {
    currentPriority.erase(id);
}

Task getTop() {
    while (!pq.empty()) {
        Task t = pq.top();
        if (currentPriority.count(t.id) &&
            currentPriority[t.id] == t.priority) {
            return t;
        }
        pq.pop();
    }
    return {-1, -1};
}
\end{lstlisting}

\subsection{易错点}

\begin{itemize}
  \item 每次取堆顶前都要清理过期元素。
  \item 更新元素时直接压入新值，不删除旧值。
  \item 判断有效性必须和当前状态表一致。
\end{itemize}

\section{T605 LRU 缓存}

\subsection{适用场景}

最近最少使用淘汰，访问后变成最新。

\begin{lstlisting}
int cap;
list<int> order; // front is most recent
unordered_map<int, list<int>::iterator> pos;

void access(int x) {
    if (cap == 0) return;

    if (pos.count(x)) {
        order.erase(pos[x]);
    } else if ((int)order.size() == cap) {
        int old = order.back();
        order.pop_back();
        pos.erase(old);
    }

    order.push_front(x);
    pos[x] = order.begin();
}
\end{lstlisting}

\subsection{易错点}

\begin{itemize}
  \item 容量为 0 要特判。
  \item 删除链表节点后，map 里的迭代器也要更新。
  \item \Code{front} 是最新还是最旧要统一。
\end{itemize}

\section{T606 路径规范化}

\subsection{适用场景}

文件路径、目录路径，处理 \Code{.}、\Code{..}、多余斜杠。

\begin{lstlisting}
vector<string> normalizePath(const string& path) {
    vector<string> st;
    vector<string> parts = split(path, '/');
    for (string part : parts) {
        if (part.empty() || part == ".") {
            continue;
        } else if (part == "..") {
            if (!st.empty()) st.pop_back();
        } else {
            st.push_back(part);
        }
    }
    return st;
}

string joinPath(const vector<string>& parts) {
    if (parts.empty()) return "/";
    string res;
    for (const string& p : parts) {
        res += "/";
        res += p;
    }
    return res;
}
\end{lstlisting}

\Warn{如果题目允许相对路径，是否能越过根目录要按题意处理。}

\section{T607 URL 参数解析}

\subsection{适用场景}

形如 \Code{/path/to?a=1\&b=2} 的字符串。

\begin{lstlisting}
struct Url {
    string path;
    map<string, string> query;
};

Url parseUrl(const string& s) {
    Url res;
    int qpos = (int)s.find('?');
    if (qpos == (int)string::npos) {
        res.path = s;
        return res;
    }

    res.path = s.substr(0, qpos);
    string qs = s.substr(qpos + 1);
    vector<string> pairs = split(qs, '&');
    for (string p : pairs) {
        int eq = (int)p.find('=');
        if (eq == (int)string::npos) {
            res.query[p] = "";
        } else {
            string key = p.substr(0, eq);
            string val = p.substr(eq + 1);
            res.query[key] = val;
        }
    }
    return res;
}
\end{lstlisting}

\subsection{易错点}

\begin{itemize}
  \item URL 可能没有 \Code{?}。
  \item 参数值可能为空。
  \item 重复参数时保留第一个还是最后一个要看题目。
\end{itemize}

\section{T608 括号匹配}

\subsection{适用场景}

判断括号是否合法，或处理简单嵌套结构。

\begin{lstlisting}
bool checkBrackets(const string& s) {
    stack<char> st;
    for (char c : s) {
        if (c == '(' || c == '[' || c == '{') {
            st.push(c);
        } else if (c == ')' || c == ']' || c == '}') {
            if (st.empty()) return false;
            char t = st.top();
            st.pop();
            if (c == ')' && t != '(') return false;
            if (c == ']' && t != '[') return false;
            if (c == '}' && t != '{') return false;
        }
    }
    return st.empty();
}
\end{lstlisting}

\section{T609 简单表达式求值}

\subsection{只含加减和非负整数}

\begin{lstlisting}
long long evalPlusMinus(const string& s) {
    long long ans = 0;
    long long num = 0;
    int sign = 1;

    for (int i = 0; i <= (int)s.size(); i++) {
        if (i < (int)s.size() && isdigit(s[i])) {
            num = num * 10 + (s[i] - '0');
        } else {
            ans += sign * num;
            num = 0;
            if (i < (int)s.size()) {
                if (s[i] == '+') sign = 1;
                else if (s[i] == '-') sign = -1;
            }
        }
    }
    return ans;
}
\end{lstlisting}

\subsection{易错点}

\begin{itemize}
  \item 上面模板不处理乘除和括号。
  \item 如果有空格，先跳过空格。
  \item 如果有负数开头，检查是否符合模板逻辑。
\end{itemize}

\chapter{易错点总表}

\section{本章使用方法}

本章用于考前最后复习和提交前检查。很多失分不是因为算法不会，而是因为输入输出、下标、溢出、初始化、比较器等细节错误。

建议使用方式：

\begin{itemize}
  \item 考前每天快速看一遍。
  \item 每次模拟赛后，把自己犯过的错补进本章。
  \item 考场提交前，用最后的总检查表快速扫一遍。
\end{itemize}

\section{输入输出易错点}

\begin{longtable}{p{4cm}p{5cm}p{5cm}}
\toprule
问题 & 表现 & 检查方法 \\
\midrule
多输出提示语 & 答案格式错误 & 不要输出 \Code{answer=}、提示文字。 \\
少换行或多内容 & 样例看似接近但不完全一致 & 严格对照题目输出格式。 \\
大小写错误 & \Code{YES} 写成 \Code{Yes} & 看题目要求。 \\
小数位错误 & 精度或格式错误 & 用 \Code{fixed << setprecision(k)}。 \\
\Code{cin} 后直接 \Code{getline} & 读到空行 & 加 \Code{cin.ignore()}。 \\
多组数据未清空 & 第一组对，后面错 & 每组重置数组、图、map、队列。 \\
\bottomrule
\end{longtable}

\section{数据类型和溢出}

\begin{longtable}{p{4cm}p{5cm}p{5cm}}
\toprule
场景 & 风险 & 正确做法 \\
\midrule
求和 & \Code{int} 溢出 & 答案用 \Code{long long}。 \\
两个大整数相乘 & 先按 \Code{int} 溢出 & 写 \Code{1LL * a * b}。 \\
最短路距离 & 路径和很大 & \Code{dist} 用 \Code{long long}。 \\
计数方案数 & 方案数巨大 & 看是否取模，用 \Code{long long}。 \\
位运算左移 & \Code{1 << k} 溢出 & 大位数用 \Code{1LL << k}。 \\
INF 太小 & 最短路被错误更新 & \Code{long long INF = 4e18}。 \\
\bottomrule
\end{longtable}

\section{整除、取整和取模}

\begin{longtable}{p{4cm}p{5cm}p{5cm}}
\toprule
场景 & 常见错误 & 正确处理 \\
\midrule
正整数分组数 & 少算最后不足一组 & $a \ge 0,b>0$ 时用 \Code{(a+b-1)/b}。 \\
负数参与除法 & 误以为 \Code{/} 是向下取整 & C++ 整数除法向 0 截断，必要时用 \Code{floor\_div/ceil\_div}。 \\
负数参与取模 & 余数下标为负导致越界 & 用 \Code{r=x\%m; if (r<0) r+=m;}。 \\
大整数取整 & 用 \Code{double} 后丢精度 & 尽量用整数公式；接近 $10^{18}$ 时避免浮点取整。 \\
除数或模数 & 没判断 0 & 除法和取模前确认分母不为 0。 \\
\bottomrule
\end{longtable}

\section{下标和边界}

\begin{itemize}
  \item 题目编号从 1 开始，代码数组是否开了 \Code{n+1}。
  \item 0 下标数组的第 $k$ 个元素是 \Code{a[k-1]}。
  \item 区间 $[l,r]$ 是否包含两端。
  \item 前缀和 1 下标公式是 \Code{pre[r] - pre[l-1]}。
  \item 访问 \Code{r+1} 时数组是否开到 \Code{n+2}。
  \item 空数组不能访问 \Code{a[0]} 或 \Code{a.back()}。
  \item 循环是 \Code{i < n} 还是 \Code{i <= n}。
\end{itemize}

\section{初始化和清空}

\begin{longtable}{p{4cm}p{5cm}p{5cm}}
\toprule
对象 & 常见错误 & 正确处理 \\
\midrule
全局数组 & 以为每组自动清空 & 多组数据手动清空。 \\
局部数组 & 未初始化 & 定义时初始化或使用 \Code{vector}。 \\
\Code{vector} & 上组残留 & \Code{assign} 或重新定义。 \\
\Code{map/set} & 上组残留 & \Code{clear()}。 \\
\Code{queue} & 没有 \Code{clear} & 每组内部重新定义。 \\
\Code{priority\_queue} & 没有 \Code{clear} & 每组内部重新定义。 \\
答案变量 & 没重置 & 每组开始设为 0 或初始值。 \\
\bottomrule
\end{longtable}

\section{排序和比较器}

\begin{itemize}
  \item 比较器表示“谁排前面”。
  \item 相等时不能返回 \Code{true}。
  \item 多关键字排序要逐项判断。
  \item 升序和降序不要写反。
  \item 需要保持原输入顺序时用 \Code{stable\_sort}。
  \item \Code{priority\_queue} 的比较器和 \Code{sort} 直觉相反，必须手测。
\end{itemize}

错误示例：

\begin{lstlisting}
return a.score >= b.score; // wrong
\end{lstlisting}

正确示例：

\begin{lstlisting}
if (a.score != b.score) return a.score > b.score;
return a.id < b.id;
\end{lstlisting}

\section{STL 容器易错点}

\begin{longtable}{p{4cm}p{5cm}p{5cm}}
\toprule
容器/函数 & 易错点 & 提醒 \\
\midrule
\Code{vector} & 越界访问 & 确认大小和下标。 \\
\Code{map[key]} & 不存在时会创建 key & 只判断存在用 \Code{count/find}。 \\
\Code{unordered\_map} & 遍历无序 & 要有序输出用 \Code{map}。 \\
\Code{set.lower\_bound} & 返回 \Code{end} 后解引用 & 先判断 \Code{it != end}。 \\
\Code{queue.front} & 空队列访问 & 先判断 \Code{empty()}。 \\
\Code{stack.top} & 空栈访问 & 先判断 \Code{empty()}。 \\
\Code{priority\_queue} & 不能删除中间元素 & 用懒删除或 set。 \\
\Code{erase} & 迭代器失效 & 删除时小心循环写法。 \\
\bottomrule
\end{longtable}

\section{字符串易错点}

\begin{itemize}
  \item \Code{cin >> s} 不能读取空格。
  \item \Code{getline} 可能读到上一行剩余换行。
  \item \Code{s.find(t)} 找不到返回 \Code{string::npos}。
  \item 不要写 \Code{if (s.find(t))}。
  \item \Code{substr(pos, len)} 第二个参数是长度。
  \item 字符串数字可能超过 \Code{long long}。
  \item 连续分隔符、开头分隔符、结尾分隔符要测试。
\end{itemize}

\section{图论易错点}

\begin{itemize}
  \item 无向图要双向加边，有向图不要。
  \item 点编号从 0 还是 1 开始。
  \item BFS 应在入队时标记访问。
  \item Dijkstra 不能处理负权边。
  \item Dijkstra 要跳过过期状态。
  \item Floyd 三层循环顺序是 \Code{k,i,j}。
  \item 并查集初始化每个点的父亲。
  \item 拓扑排序边方向不要反。
\end{itemize}

\section{动态规划易错点}

\begin{itemize}
  \item 先定义 \Code{dp[i]} 的含义，再写转移。
  \item 初始状态不要漏。
  \item 01 背包容量倒序。
  \item 完全背包容量正序。
  \item 求最大值时初始值是否应该是负无穷。
  \item 方案数是否需要取模。
  \item 多组数据是否清空 DP 数组。
\end{itemize}

\section{数据结构易错点}

\begin{itemize}
  \item 树状数组通常用 1 下标。
  \item 线段树递归区间边界要统一。
  \item 懒标记下传不要漏。
  \item 单调栈比较符号决定是否允许相等。
  \item 单调队列先删除过期元素。
  \item LCA 的 \Code{LOG} 要足够大。
  \item 递归 DFS 在大数据上可能爆栈。
\end{itemize}

\section{提交前总检查}

\begin{enumerate}
  \item 删除所有调试输出。
  \item 确认没有输出提示语。
  \item 确认输入格式和题目一致。
  \item 确认多组数据清空状态。
  \item 确认答案和中间乘法使用 \Code{long long}。
  \item 确认数组大小足够。
  \item 确认下标从 0 还是 1 开始。
  \item 确认排序比较器正确。
  \item 确认特殊边界测试过。
  \item 确认提交的是当前最新代码。
\end{enumerate}

\section{最后十分钟清单}

如果考试快结束：

\begin{itemize}
  \item 不要重构。
  \item 不要大改已经能过样例的代码。
  \item 删除调试输出。
  \item 提交当前最稳版本。
  \item 如果还有时间，再做小修小补。
  \item 对没把握的题，至少提交能过样例或小数据的版本。
\end{itemize}


\part{真题校准}
\chapter{历年真题索引与模板映射}\label{ch:11}

\section{怎样查题目}\label{sec:11-1}
这一章可以按场次、题号查找，也可以按后面的题面关键词反查。前面列出近两年的题目，后面按 A、B、C、D/E 分类整理历年题目。

模板编号指向能借用的代码片段。例如 LCA 可以帮助找树上路径，却不会直接算出路径 mex；区间加线段树也不能原样用来做区间赋值。练习时仍要结合题目的条件。

\section{按场次查找}\label{sec:recent}
这里收录 2024 年 9 月至 2026 年 5 月的八场题目，共 40 题。题号链接到公开题面，表中列出可参考的写法和容易漏掉的条件。

各场起始页：202409（第\pageref{recent:202409}页）、202412（第\pageref{recent:202412}页）、202503（第\pageref{recent:202503}页）、202506（第\pageref{recent:202506}页）、202509（第\pageref{recent:202509}页）、202512（第\pageref{recent:202512}页）、202603（第\pageref{recent:202603}页）、202605（第\pageref{recent:202605}页）。

\Needspace{20\baselineskip}
\section{202409 场次速查}\label{recent:202409}
\begingroup
\small
\setlength{\tabcolsep}{4pt}
\begin{longtable}{p{3.1cm}p{5.5cm}p{6.2cm}}
\toprule
题号与题名 & 可参考的做法 & 要留意的条件 \\
\midrule
\endfirsthead
\toprule
题号与题名 & 可参考的做法 & 要留意的条件 \\
\midrule
\endhead
\href{https://oj.shumeng.tech/p/CSP202409A}{202409-1} 密码 & 逐字符分类与频次统计；\Tref{106}。 & 字母是一类，大小写字符分别计数；先检查三类字符，再判断频次。 \\
\href{https://oj.shumeng.tech/p/CSP202409B}{202409-2} 字符串变换 & 有限字符映射重复执行；\Tref{113} 的倍增或逐环取模。 & 含空格，用 getline；每个询问从原串出发；未给出的映射保持自身。 \\
\href{https://oj.shumeng.tech/p/CSP202409C}{202409-3} 补丁应用 & 按行解析块、校验旧内容、搜索偏移并替换；\Tref{601}、\Tref{602}。 & 保留行首空格；先检验合法性，再修改。简化子任务：合法、无偏移、小文本。 \\
\href{https://oj.shumeng.tech/p/CSP202409D}{202409-4} 通讯延迟 & 节点与覆盖它的基站建二部图，边权分别为延迟/0；\Tref{204}。 & 小数据可按基站枚举点对建边；正方形含边界；不可达输出 Nan。 \\
\href{https://oj.shumeng.tech/p/CSP202409E}{202409-5} 木板切割 & 序列分裂后统计不同色与颜色段；先做小规模 vector 模拟。 & 段编号固定，不随切割重排；空木板两项均为 0。小数据 n,m,k 不超过 2000。 \\
\bottomrule
\end{longtable}
\endgroup
\Needspace{20\baselineskip}
\section{202412 场次速查}\label{recent:202412}
\begingroup
\small
\setlength{\tabcolsep}{4pt}
\begin{longtable}{p{3.1cm}p{5.5cm}p{6.2cm}}
\toprule
题号与题名 & 可参考的做法 & 要留意的条件 \\
\midrule
\endfirsthead
\toprule
题号与题名 & 可参考的做法 & 要留意的条件 \\
\midrule
\endhead
\href{https://oj.shumeng.tech/p/CSP202412A}{202412-1} 移动 & 逐命令尝试新坐标；第\ref{ch:05}章方向与边界。 & 越界指令被忽略；题目 f/b 改 y，l/r 改 x，不要套错坐标轴。 \\
\href{https://oj.shumeng.tech/p/CSP202412B}{202412-2} 梦境巡查 & 将各段最低初始能量写成前缀消耗，维护前后缀最大值。 & 去掉第 i 处补给，只会抬高它之后的要求；小数据逐个禁用补给重算。 \\
\href{https://oj.shumeng.tech/p/CSP202412C}{202412-3} 缓存模拟 & 每个缓存组独立 LRU，附有效位、脏位；\Tref{605} 提供顺序维护。 & 淘汰脏行先写回，再载入；命中读写都更新次序。单路或无需淘汰子任务可先做。 \\
\href{https://oj.shumeng.tech/p/CSP202412D}{202412-4} 跳房子 & 状态是回退后的格子；BFS 枚举可落点；\Tref{201} 提供框架。 & 一次跳到 j 后必须回退到 j-a[j]；n 不超过 1000 时可显式枚举边。 \\
\href{https://oj.shumeng.tech/p/CSP202412E}{202412-5} 梦魔 & 两侧扩张与最低初始攻击力；小数据先逐起点重算，需专门推导。 & 只有左右最近的梦魔能打；收益为非负时可优先打当前可击杀对象。更新后不能沿用旧答案。 \\
\bottomrule
\end{longtable}
\endgroup
\Needspace{20\baselineskip}
\section{202503 场次速查}\label{recent:202503}
\begingroup
\small
\setlength{\tabcolsep}{4pt}
\begin{longtable}{p{3.1cm}p{5.5cm}p{6.2cm}}
\toprule
题号与题名 & 可参考的做法 & 要留意的条件 \\
\midrule
\endfirsthead
\toprule
题号与题名 & 可参考的做法 & 要留意的条件 \\
\midrule
\endhead
\href{https://oj.shumeng.tech/p/CSP202503A}{202503-1} 数值积分 & 枚举闭区间内偶数 x，累加题设公式；第\ref{ch:05}章公式题。 & 这是题设离散估计，不是求解析积分；起始偶数、端点与最后乘 2。 \\
\href{https://oj.shumeng.tech/p/CSP202503B}{202503-2} 机器人饲养指南 & 完全背包：一天吃 i 个视为重量 i、价值 A[i]；\Tref{302}、\Tref{303}。 & 须喂完 n 个，使用恰好状态；收益未必单调，不按单位收益贪心。 \\
\href{https://oj.shumeng.tech/p/CSP202503C}{202503-3} 模板展开 & 直接赋值存长度值，间接赋值存表达式依赖；按需递归求长度并取模。 & 间接引用每次求值要反映当前变量；自引用直接赋值先算旧值；只求长度，无需构造巨串。 \\
\href{https://oj.shumeng.tech/p/CSP202503D}{202503-4} 集体锻炼 & 区间 gcd 的贡献求和；\Tref{512}。小数据双循环滚动更新 gcd。 & 贡献还乘左右端点 l、r，及时取模；满分需压缩相同 gcd 状态，未收专用模板。 \\
\href{https://oj.shumeng.tech/p/CSP202503E}{202503-5} 收费标准评估 & 带负权的最大连通子树；小数据更新后重跑树 DP。 & 只操作 1/2 的小数据先做；dp[u]=a[u]+各孩子 max(0,dp[v])，非空连通块不能全部取 0。 \\
\bottomrule
\end{longtable}
\endgroup
\Needspace{20\baselineskip}
\section{202506 场次速查}\label{recent:202506}
\begingroup
\small
\setlength{\tabcolsep}{4pt}
\begin{longtable}{p{3.1cm}p{5.5cm}p{6.2cm}}
\toprule
题号与题名 & 可参考的做法 & 要留意的条件 \\
\midrule
\endfirsthead
\toprule
题号与题名 & 可参考的做法 & 要留意的条件 \\
\midrule
\endhead
\href{https://oj.shumeng.tech/p/CSP202506A}{202506-1} 正态分布 & 只求表格行列，用整数 t=100(n-mu)/sigma；行 t/10+1，列 t\%10+1。 & 题设 sigma 是 100 的因子，整数计算可避免浮点截断误差；无需计算概率函数。 \\
\href{https://oj.shumeng.tech/p/CSP202506B}{202506-2} 机器人复健指南 & 马步八方向 BFS，统计距离不超过 k 的格子；\Tref{201}。 & 移动是 (1,2)/(2,1) 的带符号组合，不是邻接八格；起点算入答案且只计一次。 \\
\href{https://oj.shumeng.tech/p/CSP202506C}{202506-3} 消息解码 & 定长位字段拆解、字符集转换、最近历史映射；\Tref{513} 仅提供进制基础。 & 先用旧历史解本条，再更新历史；间接推断的代号不能回填成直接收到。先做指定消息类型子任务。 \\
\href{https://oj.shumeng.tech/p/CSP202506D}{202506-4} 月票发行 & 固定短模式的自动机/状态 DP；跟踪先 ccf 后 cspark。 & 短串子任务也不宜枚举所有 26 的 n 次方字符串；普通 KMP 只匹配，不能直接计数所有构造。 \\
\href{https://oj.shumeng.tech/p/CSP202506E}{202506-5} 博物馆 & 删除一个点后的连通块统计；小图枚举施工点，BFS 各块的攻略点数。 & 每个施工点取某连通块中最多的攻略点数再求和；施工点不能参观；大图需割点/点双等未收内容。 \\
\bottomrule
\end{longtable}
\endgroup
\Needspace{20\baselineskip}
\section{202509 场次速查}\label{recent:202509}
\begingroup
\small
\setlength{\tabcolsep}{4pt}
\begin{longtable}{p{3.1cm}p{5.5cm}p{6.2cm}}
\toprule
题号与题名 & 可参考的做法 & 要留意的条件 \\
\midrule
\endfirsthead
\toprule
题号与题名 & 可参考的做法 & 要留意的条件 \\
\midrule
\endhead
\href{https://oj.shumeng.tech/p/CSP202509A}{202509-1} 蒙特卡洛 & 统计给定点满足 x*x+y*y<=a*a，输出 4.0*count/n。 & 输入已经给出样本，不要自己随机生成；含圆周；两位小数可放大 100 后用整数平方比较。 \\
\href{https://oj.shumeng.tech/p/CSP202509B}{202509-2} 水印检查 & 枚举 5x9 位置；每块得到阈值区间，差分合并；\Tref{103}。 & 黑点要求 k 大于其最大灰度，白点要求 k 不超过其最小灰度；区间为 [maxBlack+1,minWhite]。 \\
\href{https://oj.shumeng.tech/p/CSP202509C}{202509-3} HTTP 头信息 & 静态表、动态表、三类指令与 Huffman 解码；\Tref{602}、\Tref{601}。 & 区分普通 H 转义与编码串，去掉末尾填充位。只含表引用，或无 Huffman 的子任务可先做。 \\
\href{https://oj.shumeng.tech/p/CSP202509D}{202509-4} 造题计划（上） & 树上路径 mex；小数据沿路径收集值再扫描最小缺失数；\Tref{208}。 & mex 从 0 开始；题设权值为排列；LCA 只帮找路径，不直接给 mex。 \\
\href{https://oj.shumeng.tech/p/CSP202509E}{202509-5} 造题计划（下） & 有先后限制的出题/验题选择；小 n 可按天做状态 DP。 & 状态记录已出/已验数量及最小精力；同天上午出完可下午验，任何时刻已验不超过已出。 \\
\bottomrule
\end{longtable}
\endgroup
\Needspace{20\baselineskip}
\section{202512 场次速查}\label{recent:202512}
\begingroup
\small
\setlength{\tabcolsep}{4pt}
\begin{longtable}{p{3.1cm}p{5.5cm}p{6.2cm}}
\toprule
题号与题名 & 可参考的做法 & 要留意的条件 \\
\midrule
\endfirsthead
\toprule
题号与题名 & 可参考的做法 & 要留意的条件 \\
\midrule
\endhead
\href{https://oj.shumeng.tech/p/CSP202512A}{202512-1} 集合 & 比较真实集合相等与异或指纹相等是否一致；\Tref{106}、排序。 & 异或相同不能证明集合相同；输出的是该判断是否正确，不是集合是否相等。 \\
\href{https://oj.shumeng.tech/p/CSP202512B}{202512-2} 数字变换 & 输入值域仅 512，枚举全部输入算 F，再建立逆映射；\Tref{113} 的有限状态思想。 & 先按题图精确实现三段位变换；复杂度约 512*m+n，无须对每个查询重跑全部操作。 \\
\href{https://oj.shumeng.tech/p/CSP202512C}{202512-3} 图片解码 & 密钥倒序，逆转局部旋转/翻转和整体旋转；第\ref{ch:07}章矩阵模拟。 & 逆序时先撤销整体旋转，再撤销局部；Z 小或只有翻转的子任务先逐格做。去掉补齐问号。 \\
\href{https://oj.shumeng.tech/p/CSP202512D}{202512-4} C 形阵 & 由等式得到 F=B，枚举 C,E 为 B 平方的约数，再算 A,G,D；\Tref{512}。 & 小 n 枚举并验证整除、所有等式与不同元素数；完美要求恰好 6 种，不是 7 种。 \\
\href{https://oj.shumeng.tech/p/CSP202512E}{202512-5} 数据抢修 & 可重集合划分为两两异或至少 W 的组；最小数据回溯分组并剪枝。 & 重复权值也占独立元素；合并后 v 离线；不要套只含小写字母的 Trie。大数据专用解法未收录。 \\
\bottomrule
\end{longtable}
\endgroup
\Needspace{20\baselineskip}
\section{202603 场次速查}\label{recent:202603}
\begingroup
\small
\setlength{\tabcolsep}{4pt}
\begin{longtable}{p{3.1cm}p{5.5cm}p{6.2cm}}
\toprule
题号与题名 & 可参考的做法 & 要留意的条件 \\
\midrule
\endfirsthead
\toprule
题号与题名 & 可参考的做法 & 要留意的条件 \\
\midrule
\endhead
\href{https://oj.shumeng.tech/p/CSP202603A}{202603-1} 平衡数 & 逐位统计正整数二进制中的 0/1；第\ref{ch:05}章位运算。 & 不统计最高 1 之前的补零；重复输入分别计数。题名是平衡数。 \\
\href{https://oj.shumeng.tech/p/CSP202603B}{202603-2} 机器人项目管理 & 普通任务 01 背包，灵活任务按收益/咖啡排序；枚举预算分配；\Tref{301}。 & 只对可分割的灵活任务用单位收益贪心；普通任务不能拆；结果保留足够小数。 \\
\href{https://oj.shumeng.tech/p/CSP202603C}{202603-3} 进程通信 & 进程接口、队列与内存区间模拟；\Tref{601}、\Tref{602}、第\ref{ch:07}章状态设计。 & 分开维护占用和有对象两种状态；最优适应先最短再最左。小范围先逐地址/区间模拟。 \\
\href{https://oj.shumeng.tech/p/CSP202603D}{202603-4} 异或 & k 进制逐位不进位加法、f 值与区间操作；先做小范围直接模拟。 & k 是奇数，不能直接用 C++ 二进制异或符号代替；\Tref{402} 的普通加法懒标记也不能直接套。 \\
\href{https://oj.shumeng.tech/p/CSP202603E}{202603-5} 旅游计划 & 树上路径和道路修复；小图逐计划走唯一路径，模拟备胎与维修。 & 强制在线时先按规则解码输入；计划途中轮胎状态要重置；镜像分 Easy/Hard，按原题约束区分。 \\
\bottomrule
\end{longtable}
\endgroup
\Needspace{20\baselineskip}
\section{202605 场次速查}\label{recent:202605}
\begingroup
\small
\setlength{\tabcolsep}{4pt}
\begin{longtable}{p{3.1cm}p{5.5cm}p{6.2cm}}
\toprule
题号与题名 & 可参考的做法 & 要留意的条件 \\
\midrule
\endfirsthead
\toprule
题号与题名 & 可参考的做法 & 要留意的条件 \\
\midrule
\endhead
\href{https://oj.shumeng.tech/p/CSP202605A}{202605-1} 银行家舍入 & 按字符串拆整数部分 a 与一位小数 b；第\ref{ch:04}章舍入说明。 & b=5 时向偶数取整，其余同四舍五入；不要用 llround 当成银行家舍入。 \\
\href{https://oj.shumeng.tech/p/CSP202605B}{202605-2} 机器人宿管指南 & 二分机器人数量，逐天模拟先变质后吃苹果；\Tref{107}、\Tref{510}。 & 每天损耗按当日剩余数量向上取整；第 m 天吃完恰为 0 仍可行；中间乘法用 long long。 \\
\href{https://oj.shumeng.tech/p/CSP202605C}{202605-3} 死锁优化 & 进程状态机，分申请/收益/释放三阶段；\Tref{603}、\Tref{110}。 & 普通、A/B/C 类分别处理；同段申请与强占优先级按题面。先完成无特殊进程子任务并识别死锁。 \\
\href{https://oj.shumeng.tech/p/CSP202605D}{202605-4} 石子游戏 & 胜负公式已给出；小区间枚举最后一段做划分 DP。 & 奇数位是相对子段起点的奇数位；预处理奇偶前缀异或，dp[r]=max(dp[l-1]+win(l,r))。 \\
\href{https://oj.shumeng.tech/p/CSP202605E}{202605-5} 绝世好串 & 树上聚拢与字典序；极小 n 枚举合法聚拢和拼接顺序作验证。 & 第二次起必须参与已有长度至少 2 的串；随意排序所有字符不一定可实现。满分专用解法未收录。 \\
\bottomrule
\end{longtable}
\endgroup


\section{历年整场索引（201912--202406）}

下表按场次汇总题目。各题的具体入口见后面的 A、B、C 和 D/E 索引。

\begingroup
\small
\setlength{\tabcolsep}{3pt}
\begin{longtable}{p{1.55cm}p{2.8cm}p{2.8cm}p{3.35cm}p{4.6cm}}
\toprule
场次 & A/B 题 & C 题 & D/E 题信号 & 判断树重点 \\
\midrule
\endfirsthead
\toprule
场次 & A/B 题 & C 题 & D/E 题信号 & 判断树重点 \\
\midrule
\endhead
202406 & 矩阵重塑（一/二） & 文本分词 & 货物调度、哥德尔机 & 矩阵下标、文本 token、优先队列、贪心/背包。 \\
202403 & 词频统计、相似度计算 & 化学方程式配平 & 十滴水、文件夹合并 & 计数、集合交并、表达式/方程、set/堆、DFS 序/线段树。 \\
202312 & 仓库规划、因子化简 & 树上搜索 & 宝藏、彩色路径 & 暴力枚举、质因数分解、树 DFS、前缀/分块、状压 DP。 \\
202309 & 坐标变换（一/二） & 梯度求解 & 阴阳龙、阻击 & 坐标前缀合成、后缀表达式、set、树形 DP。 \\
202305 & 重复局面、矩阵运算 & 解压缩 & 电力网络、闪耀巡航 & 字符串哈希、矩阵乘法顺序、位运算解析、最短路/状压。 \\
202303 & 田地丈量、垦田计划 & LDAP & 星际网络II、施肥 & 矩形相交、二分答案、递归表达式、线段树/离散化。 \\
202212 & 现值计算、训练计划 & JPEG 解码 & 聚集方差、星际网络 & 公式、拓扑顺序、矩阵扫描、启发式合并、线段树建图。 \\
202209 & 如此编码、何以包邮 & 防疫大数据 & 通信/调度类复杂题 & 进制分解、01 背包、时间窗口、集合维护。 \\
202206 & 归一化处理、寻宝大冒险 & 角色授权 & 高级模拟/几何题 & 小数输出、坐标哈希、权限集合、规则匹配。 \\
202203 & 未初始化警告、出行计划 & 计算资源调度器 & 图/数据结构题 & set 判断、差分区间、对象筛选、优先级排序。 \\
202112 & 序列查询、新解 & 登机牌号码 & 磁盘文件操作、极差路径 & 分段贡献、数学优化、编码解析、线段树。 \\
202109 & 数组推导、非零段划分 & 脉冲神经网络 & 收集卡牌、箱根山岳险天下 & 前缀最大反推、贡献扫描、大模拟、概率/图。 \\
202104 & 灰度直方图、邻域均值 & DHCP服务器 & 校门外的树、疫苗运输 & 计数数组、二维前缀、大模拟状态机、DP。 \\
202012 & 安全指数、最佳阈值 & 带配额的文件系统 & 食材运输、星际旅行 & 公式、排序阈值扫描、目录树/配额、大模拟到树/图。 \\
202009 & 检测点查询、风险人群筛查 & 点亮数字人生 & 星际旅行、密信与计数 & 距离排序、连续区间判断、拓扑/逻辑门、图与计数。 \\
202006 & 线性分类器、稀疏向量 & Markdown渲染器 & 1246、乔乔和牛牛逛超市 & 坐标代入、双指针、文本渲染、DP/高级优化。 \\
201912 & 报数、回收站选址 & 化学方程式 & 区块链、魔数 & 简单模拟、坐标哈希、表达式解析、图/事件模拟。 \\
\bottomrule
\end{longtable}
\endgroup

练完一套题后，可以在相应场次旁记下自己漏看的条件或想错的做法。

\section{A 题索引}

\begingroup
\small
\setlength{\tabcolsep}{3pt}
\begin{longtable}{p{1.8cm}p{2.8cm}p{2.6cm}p{3.5cm}p{3.5cm}}
\toprule
编号 & 题名 & 题型 & 主要考点 & 翻到哪里 \\
\midrule
\endfirsthead
\toprule
编号 & 题名 & 题型 & 主要考点 & 翻到哪里 \\
\midrule
\endhead
202406-1 & 矩阵重塑（其一） & 下标 / 矩阵 & 一维二维下标互转 & A 题二维数组 \\
202403-1 & 词频统计 & 统计 & 循环计数、频率 & A 题频率统计 \\
202312-1 & 仓库规划 & 枚举 / 判断 & 双重循环、支配关系 & A 题简单模拟 \\
202309-1 & 坐标变换（其一） & 坐标模拟 & 坐标平移 & A 题坐标入门 \\
202305-1 & 重复局面 & 哈希统计 & 字符串 / 局面计数 & A 题频率统计、字符串 \\
202303-1 & 田地丈量 & 坐标 / 几何 & 矩形相交面积 & A 题坐标距离、公式 \\
202212-1 & 现值计算 & 公式 / 小数 & 幂、浮点计算 & A 题公式题 \\
202209-1 & 如此编码 & 进制 / 模拟 & 数位分解 & A 题进制处理 \\
202206-1 & 归一化处理 & 小数 / 统计 & 平均、方差、小数输出 & A 题小数输出 \\
202203-1 & 未初始化警告 & 模拟 / set & 变量出现判断 & A 题存在判断、set \\
202112-1 & 序列查询 & 模拟 / 查询 & 下标、区间含义 & A 题区间入门 \\
202109-1 & 数组推导 & 模拟 / 统计 & 最小最大和 & A 题简单统计 \\
202104-1 & 灰度直方图 & 统计 & 计数数组 & A 题计数数组 \\
202012-1 & 期末预测之安全指数 & 统计 / 公式 & 加权求和、取 max & A 题简单统计 \\
202009-1 & 称检测点查询 & 排序 / 距离 & 距离计算、排序 & A 题排序、坐标距离 \\
202006-1 & 线性分类器 & 模拟 / 几何判断 & 代入直线、分类统计 & A 题公式题、坐标 \\
201912-1 & 报数 & 模拟 & 跳过含 7 或 7 倍数 & A 题简单模拟 \\
201909-1 & 小明种苹果 & 统计 & 总量、最大减少量 & A 题简单统计 \\
201903-1 & 小中大 & 排序 / 输出格式 & 中位数、整数/小数输出 & A 题排序、格式 \\
201812-1 & 小明上学 & 时间模拟 & 红绿灯时间累计 & A 题时间处理 \\
201809-1 & 卖菜 & 数组模拟 & 邻域平均 & A 题数组边界 \\
201803-1 & 跳一跳 & 状态模拟 & 连续奖励计分 & A 题状态扫描 \\
201712-1 & 最小差值 & 排序 & 相邻差最小 & A 题排序 \\
201709-1 & 打酱油 & 数学 / 贪心 & 优惠规则分解 & A 题公式题 \\
201703-1 & 分蛋糕 & 贪心模拟 & 累加达到阈值分组 & A 题简单模拟 \\
201612-1 & 中间数 & 排序 / 统计 & 比它大和小数量 & A 题排序/计数 \\
201609-1 & 最大波动 & 扫描 & 相邻差绝对值最大 & A 题简单统计 \\
201604-1 & 折点计数 & 扫描 & 左右邻居比较 & A 题简单扫描 \\
201512-1 & 数位之和 & 数位 & 十进制拆位 & A 题数字处理 \\
201509-1 & 数列分段 & 扫描 & 相邻变化计数 & A 题简单扫描 \\
201503-1 & 图像旋转 & 矩阵 & 下标变换输出 & A 题二维表格 \\
201412-1 & 门禁系统 & 统计 & 出现次数累计 & A 题频率统计 \\
201409-1 & 相邻数对 & 排序 / 计数 & 排序后相邻比较 & A 题排序 \\
201403-1 & 相反数 & 统计 / 查找 & 取相反数、集合判断 & A 题存在判断 \\
201312-1 & 出现次数最多的数 & 统计 & 计数数组 / map & A 题频率统计 \\
\bottomrule
\end{longtable}
\endgroup

\section{B 题索引}

\begingroup
\small
\setlength{\tabcolsep}{3pt}
\begin{longtable}{p{1.8cm}p{2.8cm}p{2.6cm}p{3.5cm}p{3.5cm}}
\toprule
编号 & 题名 & 题型 & 主要考点 & 翻到哪里 \\
\midrule
\endfirsthead
\toprule
编号 & 题名 & 题型 & 主要考点 & 翻到哪里 \\
\midrule
\endhead
202406-2 & 矩阵重塑（其二） & 矩阵 / 模拟 & 下标映射、批量输出 & B 题矩阵、复杂度 \\
202403-2 & 相似度计算 & 集合 & 交集、并集、去重 & B 题哈希统计 \\
202312-2 & 因子化简 & 数学 / 因数分解 & 质因数分解、幂 & B 题基础数学 \\
202309-2 & 坐标变换（其二） & 前缀 / 坐标变换 & 操作前缀合成 & B 题前缀预处理 \\
202305-2 & 矩阵运算 & 矩阵 / 模拟优化 & 计算顺序、复杂度 & B 题表格矩阵、复杂度 \\
202303-2 & 垦田计划 & 二分答案 / 贪心 & 最小时间、check & B 题二分答案 \\
202212-2 & 训练计划 & 拓扑 / 模拟 & 最早最晚时间 & B 题拓扑排序入口 \\
202209-2 & 何以包邮？ & DP / 背包 & 01 背包 & B 题简单 DP \\
202206-2 & 寻宝！大冒险！ & 坐标 / 哈希 & 点集匹配 & B 题哈希统计、坐标 \\
202203-2 & 出行计划 & 差分 / 区间统计 & 时间区间覆盖 & B 题差分 \\
202112-2 & 序列查询新解 & 前缀 / 数学 & 分段贡献 & B 题前缀和、排序扫描 \\
202109-2 & 非零段划分 & 差分 / 扫描 & 贡献、扫描统计 & B 题差分、排序扫描 \\
202104-2 & 邻域均值 & 二维前缀和 & 子矩阵平均 & B 题二维前缀和 \\
202012-2 & 期末预测之最佳阈值 & 排序 / 前缀统计 & 阈值扫描、计数 & B 题排序加扫描 \\
202009-2 & 风险人群筛查 & 模拟 / 坐标 & 连续停留判断 & B 题排序扫描、模拟 \\
202006-2 & 稀疏向量 & 双指针 / map & 有序向量点积 & B 题双指针、哈希统计 \\
201912-2 & 回收站选址 & 坐标 / 哈希 & 邻居存在、评分统计 & B 题哈希坐标 \\
201909-2 & 小明种苹果（续） & 统计 / 标记 & 掉落判断、连续组 & B 题统计扫描 \\
201903-2 & 二十四点 & 表达式 / 栈 & 运算优先级、求值 & C 题表达式入门 \\
201812-2 & 小明放学 & 时间 / 模拟 & 周期取模、信号灯 & B 题时间处理 \\
201809-2 & 买菜 & 区间重叠 & 时间段交集累计 & B 题区间扫描 \\
201803-2 & 碰撞的小球 & 模拟 & 位置更新、碰撞反向 & B 题模拟 \\
201712-2 & 游戏 & 队列模拟 & 报数淘汰 & B 题队列 \\
201709-2 & 公共钥匙盒 & 事件排序 & 还钥匙优先、时间排序 & C 题事件模拟 \\
201703-2 & 学生排队 & 序列模拟 & 移动元素 & B 题 vector 操作 \\
201612-2 & 工资计算 & 反推 / 枚举 & 税率分段、枚举答案 & B 题分段函数 \\
201609-2 & 火车购票 & 模拟 / 贪心 & 连续座位优先 & B 题模拟 \\
201604-2 & 俄罗斯方块 & 矩阵模拟 & 碰撞检测、落点 & B 题矩阵模拟 \\
201512-2 & 消除类游戏 & 二维扫描 & 标记后统一清除 & B 题矩阵模拟 \\
201509-2 & 日期计算 & 日期 & 月份天数、闰年 & A/B 日期模板 \\
201503-2 & 数字排序 & 排序 / 频率 & 频次降序、值升序 & B 题排序比较器 \\
201412-2 & Z字形扫描 & 矩阵遍历 & 对角线方向切换 & B 题矩阵扫描 \\
201409-2 & 画图 & 二维标记 & 矩形覆盖面积 & B 题二维表格 \\
201403-2 & 窗口 & 模拟 / 栈序 & 点击命中、置顶 & B 题模拟、列表 \\
201312-2 & ISBN号码 & 字符串 / 校验 & 加权求和、校验码 & A 题字符串、公式 \\
\bottomrule
\end{longtable}
\endgroup

\section{C 题索引：大模拟与中档算法入口}

C 题按输入中的对象、规则和字符串格式查找。表中的做法可以作参考，具体实现还要看数据范围。

\begingroup
\small
\setlength{\tabcolsep}{3pt}
\begin{longtable}{p{1.75cm}p{2.55cm}p{3.0cm}p{3.6cm}p{3.0cm}}
\toprule
编号 & 题名 & 第一识别信号 & 模板入口 & 易错提醒 \\
\midrule
\endfirsthead
\toprule
编号 & 题名 & 第一识别信号 & 模板入口 & 易错提醒 \\
\midrule
\endhead
202406-3 & 文本分词 & 词表、token、合并规则 & \Tref{601}、\Tref{604}、字符串解析 & 输入文本和词表顺序不要混。 \\
202403-3 & 化学方程式配平 & 方程、元素、系数未知 & 表达式解析、线性方程 & 元素名大小写和括号倍数。 \\
202312-3 & 树上搜索 & 树、子树、搜索顺序 & DFS、树模拟 & 子树权值和、删除/保留状态。 \\
202309-3 & 梯度求解 & 表达式、变量、后缀计算 & 后缀栈求导 & 操作数顺序、取模和导数规则。 \\
202305-3 & 解压缩 & 编码串、块、长度 & 字符串解析、位运算 & 十六进制、长度字段、边界。 \\
202303-3 & LDAP & 递归条件表达式 & 递归下降、集合/bitset & 与/或/非优先级和括号。 \\
202212-3 & JPEG 解码 & 8x8 矩阵、Z 字形、变换 & 矩阵模拟、浮点/取整 & 下标顺序、四舍五入、截断范围。 \\
202209-3 & 防疫大数据 & 时间窗口、风险地区、人员轨迹 & \Tref{601}、set/map & 生效时间和过期时间。 \\
202206-3 & 角色授权 & 用户、角色、权限、资源 & map/set、规则匹配 & 通配符、用户/用户组关联、权限匹配。 \\
202203-3 & 计算资源调度器 & 节点选择、多级约束 & \Tref{601}、排序、set & 优先级和约束过滤顺序。 \\
202112-3 & 登机牌号码 & 编码、校验、字符串转换 & 进制/编码、模拟 & 字符集映射和补位。 \\
202109-3 & 脉冲神经网络 & 多轮状态更新 & 大模拟、滚动数组 & 同一轮更新要先累计后生效。 \\
202104-3 & DHCP服务器 & 请求/响应、租约状态 & \Tref{603} 状态机、事件时间 & 到期时间和状态转移顺序。 \\
202012-3 & 带配额的文件系统 & 目录树、文件大小、配额 & \Tref{606}、树、map & 路径创建失败要回滚。 \\
202009-3 & 点亮数字人生 & 逻辑门、依赖顺序 & \Tref{205} 拓扑排序 & 环检测、多个输入输出。 \\
202006-3 & Markdown渲染器 & 文本块、列表、段落 & \Tref{601}、字符串解析 & 空行分块和宽度换行。 \\
201912-3 & 化学方程式 & 化学式、元素守恒 & 字符串解析、map & 多位系数、括号嵌套。 \\
201909-3 & 字符画 & 像素块、RGB、转义输出 & 矩阵模拟、字符串输出 & 转义序列和连续颜色压缩。 \\
201903-3 & 损坏的RAID5 & 磁盘、条带、异或恢复 & 模拟、位运算 & 十六进制转换和缺盘数量。 \\
201812-3 & CIDR合并 & IP/前缀、排序、合并 & \Tref{513}、排序扫描 & IP 转整数、包含关系。 \\
201809-3 & 元素选择器 & CSS 选择器、层级匹配 & 字符串解析、树/栈 & 大小写、祖先关系。 \\
201803-3 & URL映射 & 路由规则、参数提取 & 字符串匹配、类型参数检查 & 参数类型、匹配失败。 \\
201712-3 & Crontab & 时间范围、周期触发 & \Tref{514}、事件枚举 & 星期、月份、闭区间。 \\
201709-3 & JSON查询 & 键值、嵌套对象、查询 & 字符串解析、栈/map & 转义字符和层级路径。 \\
201703-3 & Markdown & 标题、列表、强调 & 字符串解析 & 块级规则优先于行内规则。 \\
201612-3 & 权限查询 & 用户/角色/权限等级 & map/set、规则判断 & 通配等级和最大权限。 \\
201609-3 & 炉石传说 & 角色、随从、攻击 & \Tref{603} 状态机 & 死亡结算顺序。 \\
201604-3 & 路径解析 & 当前目录、相对路径 & \Tref{606} 路径规范化 & 多个斜杠、\Code{..} 到根目录。 \\
201509-3 & 模板生成系统 & 变量替换、文本模板 & \Tref{501} split、map & 引号内容、未定义变量为空。 \\
201403-3 & 命令行选项 & 选项、参数、命令行 & \Tref{602} 命令分发 & 有参/无参选项区别。 \\
\bottomrule
\end{longtable}
\endgroup

\section{D/E 题索引与小数据做法}

D/E 题先看有无适合暴力的小数据，或链、树、无修改等特殊条件。表中列出一些可尝试的方向。

\begingroup
\small
\setlength{\tabcolsep}{3pt}
\begin{longtable}{p{2.0cm}p{3.0cm}p{4.0cm}p{4.9cm}}
\toprule
编号 & 题名 & 参考方向 & 小数据做法 \\
\midrule
\endfirsthead
\toprule
编号 & 题名 & 参考方向 & 小数据做法 \\
\midrule
\endhead
202403-4 & 十滴水 & set/map、优先队列、连锁反应 & 小数据暴力模拟每次爆裂。 \\
202403-5 & 文件夹合并 & DFS 序、链表、线段树 & 只做小树 DFS 合并，或无修改子任务。 \\
202312-4 & 宝藏 & 分块、前缀统计 & 暴力枚举查询，小范围预处理。 \\
202312-5 & 彩色路径 & 状压 DP、折半搜索 & 颜色数小就状压，否则 DFS 小图。 \\
202309-4 & 阴阳龙 & 动态 set / 平衡树 & 小数据每次排序或线性扫描。 \\
202309-5 & 阻击 & 树形 DP、线段树、树剖 & 链/小树暴力 DFS。 \\
202305-4 & 电力网络 & 图论、最短路/割点信号 & 小图建图后暴力枚举。 \\
202305-5 & 闪耀巡航 & 最短路、状压 DP & 点数小用 Floyd + 状压。 \\
202303-4 & 星际网络II & 线段树、离散化 & 小范围直接数组模拟。 \\
202303-5 & 施肥 & 分治、线段树/树状数组 & 差分可拿部分分。 \\
202212-4 & 聚集方差 & 启发式合并 & 小数据暴力维护集合。 \\
202212-5 & 星际网络 & 线段树建图、高精度 & 小图建边后跑最短路。 \\
202112-4 & 磁盘文件操作 & 线段树、区间覆盖 & 小范围数组模拟。 \\
202104-4 & 校门外的树 & DP & 先写 $O(n^2)$ 或按区间暴力。 \\
202009-4 & 星际旅行 & 图/几何/矩阵类优化 & 小数据按定义模拟。 \\
201812-4 & 数据中心 & 最小生成树 & \Tref{207} Kruskal。 \\
201412-4 & 最优灌溉 & 最小生成树 & \Tref{207} Kruskal。 \\
\bottomrule
\end{longtable}
\endgroup

\section{题目反查模板索引}

如果你记得“题看起来像什么”，但想不起算法名，按这一表反查。

\begin{longtable}{p{4.2cm}p{5.0cm}p{5.0cm}}
\toprule
题面样子 & 优先翻看 & 典型真题 \\
\midrule
\endfirsthead
\toprule
题面样子 & 优先翻看 & 典型真题 \\
\midrule
\endhead
一堆对象，每个对象有多个属性，要求按规则选择 & C 题状态机、大模拟骨架、排序比较器 & 计算资源调度器、DHCP服务器、文本分词 \\
路径、目录、文件、配额、URL、JSON、Markdown & 字符串解析、路径规范化、递归结构 & 路径解析、带配额的文件系统、JSON查询、URL映射 \\
矩阵、图像、8x8、Z 字形、旋转、重塑 & 二维数组、矩阵扫描、下标变换 & JPEG 解码、矩阵重塑、Z字形扫描、图像旋转 \\
多次询问，每次问区间/矩形/时间覆盖 & 前缀和、差分、线段树 & 邻域均值、出行计划、磁盘文件操作 \\
要求最小时间、最大收益、阈值最优，且答案越大越容易/越难 & 二分答案 & 垦田计划、资源压缩题 \\
依赖关系、先后顺序、逻辑门、任务计划 & 拓扑排序 & 训练计划、点亮数字人生 \\
集合合并、连通块、最小连接代价 & 并查集、Kruskal & 数据中心、最优灌溉 \\
字符串反复变换、替换函数、循环节 & 映射图、环、倍增/取模 & 字符串变换类题 \\
\bottomrule
\end{longtable}

\section{按题面信号反查真题}

记不清题号时，可以按题面中的操作、对象或数据结构查找。

\begingroup
\small
\setlength{\tabcolsep}{3pt}
\begin{longtable}{p{3.0cm}p{3.6cm}p{4.2cm}p{4.0cm}}
\toprule
题面信号 & 第一判断 & 类似真题 & 模板入口 \\
\midrule
\endfirsthead
\toprule
题面信号 & 第一判断 & 类似真题 & 模板入口 \\
\midrule
\endhead
只要求统计数量、最大最小、频次 & A 题扫描/计数 & 出现次数最多的数、灰度直方图、词频统计 & A 题统计、\Tref{106} \\
相邻元素、连续段、分段数量 & 一遍扫描或排序后扫描 & 数列分段、折点计数、非零段划分 & A 题扫描、B 题排序扫描 \\
矩形、坐标点、覆盖面积 & 坐标公式或哈希点集 & 田地丈量、画图、回收站选址、寻宝大冒险 & 坐标、\Tref{106}、\Tref{109} \\
矩阵重塑、旋转、Z 字形、局部平均 & 二维下标或二维前缀 & 图像旋转、Z字形扫描、邻域均值、矩阵重塑 & \Tref{102}、矩阵扫描 \\
时间段、红绿灯、租约、到期 & 时间模拟或事件排序 & 小明上学、小明放学、公共钥匙盒、DHCP服务器 & \Tref{110}、\Tref{514}、\Tref{603} \\
稀疏数据，两列编号和值 & map 或双指针 & 稀疏向量、坐标变换、回收站选址 & \Tref{106}、\Tref{108} \\
阈值、最优分界点、排序后评价 & 排序扫描或前缀统计 & 最佳阈值、非零段划分 & \Tref{105}、\Tref{101} \\
最小时间、最大可行值、资源压缩 & 二分答案 & 垦田计划、资源压缩类题 & \Tref{107} \\
选择若干物品，金额达到目标 & 可达金额 DP & 何以包邮 & \Tref{303} \\
依赖任务、前置关系、逻辑门 & 拓扑排序 & 训练计划、点亮数字人生 & \Tref{205} \\
路径、目录、文件、配额 & 字符串解析 + 树形结构 & 路径解析、带配额的文件系统、文件夹合并 & \Tref{606}、\Tref{601}、\Tref{402} \\
JSON、URL、Markdown、模板替换 & 格式解析 & JSON查询、URL映射、Markdown、模板生成系统 & \Tref{501}、\Tref{607}、\Tref{601} \\
化学式、表达式、括号、方程 & 按语法和任务区分 & 化学方程式、配平、梯度求解、二十四点 & 算术见 \Tref{610}；化学/求导需专用实现 \\
缓存、最近最少使用、淘汰 & 链表/map 或队列模拟 & 缓存类模拟题 & \Tref{605} \\
对象有状态，命令改变状态 & 状态机 & DHCP服务器、炉石传说、角色授权、计算资源调度器 & \Tref{603}、\Tref{602} \\
每次取最优对象，可能删除失效对象 & 堆或 set，必要时懒删除 & 文本分词、任务调度类题、十滴水 & \Tref{604} \\
树上子树、祖先、路径、合并 & DFS 序、LCA、树形 DP & 树上搜索、文件夹合并、阻击 & \Tref{208}、\Tref{402} \\
动态区间修改查询 & 线段树/树状数组 & 磁盘文件操作、星际网络II、施肥 & \Tref{401}、\Tref{402} \\
最小连接代价、网络建设 & 最小生成树 & 最优灌溉、数据中心 & \Tref{207} \\
点少、状态多、颜色/集合限制 & 状压 DP & 彩色路径、闪耀巡航 & 第 8 章状压入口 \\
\bottomrule
\end{longtable}
\endgroup

\section{练习后怎样复盘}

做完一道题，可以对照最初的思路和实际做法，记下差别。

\begin{longtable}{p{3.0cm}p{5.2cm}p{5.4cm}}
\toprule
检查项 & 如果答不上来 & 应补到哪里 \\
\midrule
\endfirsthead
\toprule
检查项 & 如果答不上来 & 应补到哪里 \\
\midrule
\endhead
这题第一眼信号是什么？ & 只记得题名，不知道题面特征 & 本章“按题面信号反查真题”。 \\
数据范围卡掉了哪些暴力？ & 没估复杂度就写代码 & 第 3 章数据范围判断树。 \\
满分模板是什么？ & 只会题解，不会迁移 & 第 9 章模板库和第 3 章关键词索引。 \\
有没有更简单部分分？ & D/E 空着没写 & 第 8 章 D/E 部分分策略。 \\
我误判成了什么？ & 例如把二分题当贪心，把表达式当普通字符串 & 第 3 章容易误判分支。 \\
这题有没有固定易错点？ & 样例过了但提交错 & 第 10 章易错点总表。 \\
\bottomrule
\end{longtable}

\section{A/B 题型归纳}

\subsection{A 题高频题型}

\begin{longtable}{p{4cm}p{5cm}p{5cm}}
\toprule
题型 & 常见信号 & 对应章节 \\
\midrule
\endfirsthead
\toprule
题型 & 常见信号 & 对应章节 \\
\midrule
\endhead
计数统计 & 出现次数、直方图、满足条件数量 & A 题简单统计、频率统计 \\
公式计算 & 安全指数、现值、加权求和 & A 题公式题 \\
坐标处理 & 检测点、坐标变换、矩形 & A 题坐标和距离 \\
字符串/哈希 & 重复局面、字符串计数 & A 题字符串、频率统计 \\
二维数组 & 图像旋转、矩阵重塑、卖菜 & A 题数组边界、二维表格 \\
简单模拟 & 按题意判断、逐项检查 & A 题简单模拟 \\
\bottomrule
\end{longtable}

\subsection{B 题高频题型}

\begin{longtable}{p{4cm}p{5cm}p{5cm}}
\toprule
题型 & 常见信号 & 对应章节 \\
\midrule
\endfirsthead
\toprule
题型 & 常见信号 & 对应章节 \\
\midrule
\endhead
前缀和 / 二维前缀和 & 多次区间或矩形查询 & B 题前缀和 \\
差分 & 时间区间覆盖、区间修改 & B 题差分 \\
排序扫描 & 阈值、排名、分段统计 & B 题排序加扫描 \\
哈希 / 集合 & 坐标点集、稀疏数据、相似度 & B 题哈希统计 \\
二分答案 & 最小化最大时间、资源压缩 & B 题二分答案 \\
背包 DP & 选择若干使总和达到要求 & B 题简单 DP \\
拓扑 / 依赖 & 任务计划、前置关系 & B 题拓扑排序 \\
矩阵模拟 & 重塑、扫描、局部统计 & B 题矩阵处理 \\
数学分解 & 因子、约数、质因数 & B 题基础数学 \\
\bottomrule
\end{longtable}


\backmatter

\chapter*{版本记录}
\addcontentsline{toc}{chapter}{版本记录}

\begin{longtable}{p{2.5cm}p{3cm}p{8cm}}
\toprule
版本 & 日期 & 内容 \\
\midrule
v0.1 & 2026-05-27 & 建立 LaTeX 项目骨架，完成手册入口、C++ 语法、A/B/C/D/E 策略、真题索引、模板库、易错点、详细判断树和考场极速索引初版。 \\
v0.2 & 2026-06-04 & 重排全书结构为“导航与判断、基础与题型、模板与检查、真题校准”，统一章节文件编号；补强 C++ 基础、整数取整和标准取模。 \\
\bottomrule
\end{longtable}

\end{document}
